\documentclass[11pt,a4paper]{article}
\usepackage[utf8]{inputenc}
\usepackage[T1]{fontenc}
\usepackage[spanish,es-noshorthands,es-nodecimaldot]{babel}
\usepackage{lmodern,microtype}
\usepackage{amsmath,amssymb,mathtools}
\usepackage[margin=2.3cm]{geometry}
\usepackage{enumitem,xcolor,titlesec}
\usepackage[colorlinks,linkcolor=acento,urlcolor=acento,citecolor=acento,bookmarksnumbered]{hyperref}
\usepackage{needspace}
\usepackage{xcolor}\definecolor{acento}{HTML}{1F5F8B}\definecolor{ok}{HTML}{1F7A4A}\definecolor{bad}{HTML}{B23A2E}
\newcommand{\R}{\mathbb{R}}\newcommand{\N}{\mathbb{N}}\newcommand{\Q}{\mathbb{Q}}
\newcommand{\K}{\mathcal{K}}\newcommand{\M}{\mathbb{M}}
\DeclareMathOperator{\diam}{diam}\DeclareMathOperator{\dis}{dis}\DeclareMathOperator{\dist}{dist}
\newcommand{\cl}[1]{\overline{#1}}\newcommand{\norm}[1]{\lVert #1\rVert}\newcommand{\abs}[1]{\lvert #1\rvert}
\providecommand{\lt}{<}\providecommand{\gt}{>}
\newlist{pasos}{itemize}{2}\setlist[pasos]{label={},leftmargin=2.4em,labelsep=0.6em,itemsep=2pt,topsep=3pt}
\newlist{subpasos}{itemize}{2}\setlist[subpasos]{label={},leftmargin=1.6em,itemsep=1pt,topsep=1pt}
\titleformat{\section}{\Large\bfseries\color{acento}}{}{0pt}{}
\titleformat{\subsection}{\large\bfseries}{}{0pt}{}
\newcommand{\parte}[1]{\section*{#1}\addcontentsline{toc}{section}{#1}}
\emergencystretch=3em
\newcommand{\bloque}[1]{\par\smallskip\noindent\textbf{\color{acento}#1}\par}
\setlength{\parindent}{0pt}\setlength{\parskip}{4pt}
\title{Análisis Real (MAT2514)\\[4pt]\large Quiz de repaso: verdadero o falso, con respuestas, ejemplos y ejercicios}
\author{Prof. Mircea Petrache --- 2026, segundo semestre}
\date{}
\begin{document}
\maketitle
\thispagestyle{empty}
\noindent Esta es la versión en PDF del quiz del sitio de repaso
(\url{https://decostruttivismo.github.io/MAT2514-repaso/}).
Cada afirmación es verdadera o falsa, y cada una corresponde a un error que descuenta puntos en los problemas de la guía del examen.

\begin{itemize}[leftmargin=1.4em,itemsep=1pt]
\item \textbf{Parte 1.} Las 33 afirmaciones, para responder sin ver la solución.
\item \textbf{Parte 2.} La respuesta de cada una con su justificación paso a paso, cuándo la afirmación sí es verdadera, y dónde aparece en la guía y en los apuntes.
\item \textbf{Parte 3.} Para cada pregunta, un ejemplo tomado de los apuntes y un ejercicio.
\item \textbf{Parte 4.} Las soluciones de los ejercicios.
\item \textbf{Apéndice.} Los enunciados de los problemas de la guía del examen, a los que se refieren las respuestas.
\end{itemize}
\textbf{Cómo navegar.} El documento es clickeable. En la Parte 1, el número de cada afirmación lleva a su respuesta; desde la respuesta se va al ejemplo y ejercicio, y desde el ejercicio a su solución. Las menciones de un problema de la guía llevan a su enunciado en el Apéndice. Las referencias del tipo ``Teorema 20.16'' y los números de clase abren los \href{https://decostruttivismo.github.io/MAT2514-repaso/pdf/MAT2514_apuntes_completos.pdf}{apuntes completos} en el lugar exacto (se necesita conexión a internet).

\tableofcontents\newpage
\parte{Parte 1. Las afirmaciones}
Marque V o F. La respuesta está en la Parte 2, con el mismo número: haga clic en el número para ir a ella.
\begin{enumerate}[leftmargin=2.2em,itemsep=8pt,label=\textbf{\arabic*.}]
\item[\hyperlink{r-1}{\textbf{1.}}] \hypertarget{q-1}{}{\small \textsc{Compacidad y distancia} · Clase \href{https://decostruttivismo.github.io/MAT2514-repaso/pdf/MAT2514_apuntes_completos.pdf\#nameddest=section.20}{20}}\hfill \textbf{V}\,\fbox{\rule{0pt}{6pt}\hspace{6pt}}\quad\textbf{F}\,\fbox{\rule{0pt}{6pt}\hspace{6pt}}\par Si $A\subseteq X$ es cerrado y no vacío, entonces para todo $x\in X$ existe $a_0\in A$ con $d(x,a_0)=d(x,A)$.\hfill{\footnotesize\hyperlink{r-1}{ver respuesta}}
\item[\hyperlink{r-2}{\textbf{2.}}] \hypertarget{q-2}{}{\small \textsc{Compacidad y distancia} · Clase \href{https://decostruttivismo.github.io/MAT2514-repaso/pdf/MAT2514_apuntes_completos.pdf\#nameddest=section.20}{20}}\hfill \textbf{V}\,\fbox{\rule{0pt}{6pt}\hspace{6pt}}\quad\textbf{F}\,\fbox{\rule{0pt}{6pt}\hspace{6pt}}\par Si $F_1,F_2\subseteq X$ son cerrados, no vacíos y disjuntos, entonces $d(F_1,F_2)\gt 0$.\hfill{\footnotesize\hyperlink{r-2}{ver respuesta}}
\item[\hyperlink{r-3}{\textbf{3.}}] \hypertarget{q-3}{}{\small \textsc{Compacidad y conexidad} · Clases \href{https://decostruttivismo.github.io/MAT2514-repaso/pdf/MAT2514_apuntes_completos.pdf\#nameddest=section.24}{24}, \href{https://decostruttivismo.github.io/MAT2514-repaso/pdf/MAT2514_apuntes_completos.pdf\#nameddest=section.20}{20}}\hfill \textbf{V}\,\fbox{\rule{0pt}{6pt}\hspace{6pt}}\quad\textbf{F}\,\fbox{\rule{0pt}{6pt}\hspace{6pt}}\par Si $K$ es compacto y disconexo, entonces $K=K_1\cup K_2$ con $K_1,K_2$ compactos, disjuntos y no vacíos; en particular $d(K_1,K_2)\gt 0$.\hfill{\footnotesize\hyperlink{r-3}{ver respuesta}}
\item[\hyperlink{r-4}{\textbf{4.}}] \hypertarget{q-4}{}{\small \textsc{Compacidad y distancia} · Clases \href{https://decostruttivismo.github.io/MAT2514-repaso/pdf/MAT2514_apuntes_completos.pdf\#nameddest=section.18}{18}, \href{https://decostruttivismo.github.io/MAT2514-repaso/pdf/MAT2514_apuntes_completos.pdf\#nameddest=section.19}{19}}\hfill \textbf{V}\,\fbox{\rule{0pt}{6pt}\hspace{6pt}}\quad\textbf{F}\,\fbox{\rule{0pt}{6pt}\hspace{6pt}}\par Todo espacio métrico totalmente acotado es compacto.\hfill{\footnotesize\hyperlink{r-4}{ver respuesta}}
\item[\hyperlink{r-5}{\textbf{5.}}] \hypertarget{q-5}{}{\small \textsc{Compacidad y Hausdorff} · Clases \href{https://decostruttivismo.github.io/MAT2514-repaso/pdf/MAT2514_apuntes_completos.pdf\#nameddest=section.18}{18}, \href{https://decostruttivismo.github.io/MAT2514-repaso/pdf/MAT2514_apuntes_completos.pdf\#nameddest=section.36}{36}}\hfill \textbf{V}\,\fbox{\rule{0pt}{6pt}\hspace{6pt}}\quad\textbf{F}\,\fbox{\rule{0pt}{6pt}\hspace{6pt}}\par Si $K$ es compacto, para cada $\varepsilon\gt 0$ tiene una $\varepsilon$-red finita $S\subseteq K$, y entonces $d_H(S,K)\le\varepsilon$.\hfill{\footnotesize\hyperlink{r-5}{ver respuesta}}
\item[\hyperlink{r-6}{\textbf{6.}}] \hypertarget{q-6}{}{\small \textsc{Completitud} · Clases \href{https://decostruttivismo.github.io/MAT2514-repaso/pdf/MAT2514_apuntes_completos.pdf\#nameddest=section.15}{15}, \href{https://decostruttivismo.github.io/MAT2514-repaso/pdf/MAT2514_apuntes_completos.pdf\#nameddest=section.37}{37}}\hfill \textbf{V}\,\fbox{\rule{0pt}{6pt}\hspace{6pt}}\quad\textbf{F}\,\fbox{\rule{0pt}{6pt}\hspace{6pt}}\par Un subconjunto cerrado de un espacio métrico completo es completo.\hfill{\footnotesize\hyperlink{r-6}{ver respuesta}}
\item[\hyperlink{r-7}{\textbf{7.}}] \hypertarget{q-7}{}{\small \textsc{Distancia de Hausdorff} · Clase \href{https://decostruttivismo.github.io/MAT2514-repaso/pdf/MAT2514_apuntes_completos.pdf\#nameddest=section.36}{36}}\hfill \textbf{V}\,\fbox{\rule{0pt}{6pt}\hspace{6pt}}\quad\textbf{F}\,\fbox{\rule{0pt}{6pt}\hspace{6pt}}\par Para calcular $d_H(A,B)$ basta calcular $\sup_{a\in A}d(a,B)$.\hfill{\footnotesize\hyperlink{r-7}{ver respuesta}}
\item[\hyperlink{r-8}{\textbf{8.}}] \hypertarget{q-8}{}{\small \textsc{Distancia de Hausdorff} · Clase \href{https://decostruttivismo.github.io/MAT2514-repaso/pdf/MAT2514_apuntes_completos.pdf\#nameddest=section.36}{36}}\hfill \textbf{V}\,\fbox{\rule{0pt}{6pt}\hspace{6pt}}\quad\textbf{F}\,\fbox{\rule{0pt}{6pt}\hspace{6pt}}\par Si $C\subseteq D$, entonces $d_H(D,C)=0$.\hfill{\footnotesize\hyperlink{r-8}{ver respuesta}}
\item[\hyperlink{r-9}{\textbf{9.}}] \hypertarget{q-9}{}{\small \textsc{Distancia de Hausdorff} · Clase \href{https://decostruttivismo.github.io/MAT2514-repaso/pdf/MAT2514_apuntes_completos.pdf\#nameddest=section.36}{36}}\hfill \textbf{V}\,\fbox{\rule{0pt}{6pt}\hspace{6pt}}\quad\textbf{F}\,\fbox{\rule{0pt}{6pt}\hspace{6pt}}\par $d_H\bigl((0,1),[0,1]\bigr)=0$, aunque los conjuntos son distintos.\hfill{\footnotesize\hyperlink{r-9}{ver respuesta}}
\item[\hyperlink{r-10}{\textbf{10.}}] \hypertarget{q-10}{}{\small \textsc{Distancia de Hausdorff} · Clases \href{https://decostruttivismo.github.io/MAT2514-repaso/pdf/MAT2514_apuntes_completos.pdf\#nameddest=section.36}{36}, \href{https://decostruttivismo.github.io/MAT2514-repaso/pdf/MAT2514_apuntes_completos.pdf\#nameddest=section.20}{20}}\hfill \textbf{V}\,\fbox{\rule{0pt}{6pt}\hspace{6pt}}\quad\textbf{F}\,\fbox{\rule{0pt}{6pt}\hspace{6pt}}\par Para $A,B\in\mathcal K(Z)$, con $\varphi_A(z)=d(z,A)$, se tiene $\sup_{z\in Z}\abs{\varphi_A(z)-\varphi_B(z)}=d_H(A,B)$.\hfill{\footnotesize\hyperlink{r-10}{ver respuesta}}
\item[\hyperlink{r-11}{\textbf{11.}}] \hypertarget{q-11}{}{\small \textsc{Distancia de Hausdorff} · Clase \href{https://decostruttivismo.github.io/MAT2514-repaso/pdf/MAT2514_apuntes_completos.pdf\#nameddest=section.37}{37}}\hfill \textbf{V}\,\fbox{\rule{0pt}{6pt}\hspace{6pt}}\quad\textbf{F}\,\fbox{\rule{0pt}{6pt}\hspace{6pt}}\par La desigualdad del Problema \hyperlink{prob-10}{10(a)} es $d_H(A_1\cup A_2,\,B_1\cup B_2)\le d_H(A_1,B_1)+d_H(A_2,B_2)$.\hfill{\footnotesize\hyperlink{r-11}{ver respuesta}}
\item[\hyperlink{r-12}{\textbf{12.}}] \hypertarget{q-12}{}{\small \textsc{Distancia de Hausdorff} · Clases \href{https://decostruttivismo.github.io/MAT2514-repaso/pdf/MAT2514_apuntes_completos.pdf\#nameddest=section.37}{37}, \href{https://decostruttivismo.github.io/MAT2514-repaso/pdf/MAT2514_apuntes_completos.pdf\#nameddest=section.19}{19}}\hfill \textbf{V}\,\fbox{\rule{0pt}{6pt}\hspace{6pt}}\quad\textbf{F}\,\fbox{\rule{0pt}{6pt}\hspace{6pt}}\par Si $A_1\supseteq A_2\supseteq\cdots$ son compactos no vacíos y $A=\bigcap_nA_n$, entonces para cada $\varepsilon\gt 0$ hay un $N$ con $A_n\subseteq A^\varepsilon$ para todo $n\ge N$.\hfill{\footnotesize\hyperlink{r-12}{ver respuesta}}
\item[\hyperlink{r-13}{\textbf{13.}}] \hypertarget{q-13}{}{\small \textsc{Distancia de Hausdorff} · Clase \href{https://decostruttivismo.github.io/MAT2514-repaso/pdf/MAT2514_apuntes_completos.pdf\#nameddest=section.36}{36}}\hfill \textbf{V}\,\fbox{\rule{0pt}{6pt}\hspace{6pt}}\quad\textbf{F}\,\fbox{\rule{0pt}{6pt}\hspace{6pt}}\par La sucesión creciente $[0,n]$ converge en $\mathcal K(\R)$ a $[0,\infty)$.\hfill{\footnotesize\hyperlink{r-13}{ver respuesta}}
\item[\hyperlink{r-14}{\textbf{14.}}] \hypertarget{q-14}{}{\small \textsc{Operador de Hutchinson} · Clase \href{https://decostruttivismo.github.io/MAT2514-repaso/pdf/MAT2514_apuntes_completos.pdf\#nameddest=section.37}{37}}\hfill \textbf{V}\,\fbox{\rule{0pt}{6pt}\hspace{6pt}}\quad\textbf{F}\,\fbox{\rule{0pt}{6pt}\hspace{6pt}}\par Si $W\bigl(\bigcap_nA_n\bigr)\subseteq\bigcap_nA_n$, entonces $\bigcap_nA_n$ es el atractor.\hfill{\footnotesize\hyperlink{r-14}{ver respuesta}}
\item[\hyperlink{r-15}{\textbf{15.}}] \hypertarget{q-15}{}{\small \textsc{Operador de Hutchinson} · Clase \href{https://decostruttivismo.github.io/MAT2514-repaso/pdf/MAT2514_apuntes_completos.pdf\#nameddest=section.37}{37}}\hfill \textbf{V}\,\fbox{\rule{0pt}{6pt}\hspace{6pt}}\quad\textbf{F}\,\fbox{\rule{0pt}{6pt}\hspace{6pt}}\par $d_H\bigl([0,1],\ [0,\frac14]\cup[\frac34,1]\bigr)=\frac12$.\hfill{\footnotesize\hyperlink{r-15}{ver respuesta}}
\item[\hyperlink{r-16}{\textbf{16.}}] \hypertarget{q-16}{}{\small \textsc{Punto fijo de Banach} · Clase \href{https://decostruttivismo.github.io/MAT2514-repaso/pdf/MAT2514_apuntes_completos.pdf\#nameddest=section.31}{31}}\hfill \textbf{V}\,\fbox{\rule{0pt}{6pt}\hspace{6pt}}\quad\textbf{F}\,\fbox{\rule{0pt}{6pt}\hspace{6pt}}\par La cota a posteriori $d(x_n,x^*)\le\frac k{1-k}d(x_n,x_{n-1})$ permite decidir de antemano cuántas iteraciones hacen falta.\hfill{\footnotesize\hyperlink{r-16}{ver respuesta}}
\item[\hyperlink{r-17}{\textbf{17.}}] \hypertarget{q-17}{}{\small \textsc{Compacidad} · Clases \href{https://decostruttivismo.github.io/MAT2514-repaso/pdf/MAT2514_apuntes_completos.pdf\#nameddest=section.19}{19}, \href{https://decostruttivismo.github.io/MAT2514-repaso/pdf/MAT2514_apuntes_completos.pdf\#nameddest=section.20}{20}}\hfill \textbf{V}\,\fbox{\rule{0pt}{6pt}\hspace{6pt}}\quad\textbf{F}\,\fbox{\rule{0pt}{6pt}\hspace{6pt}}\par En un espacio métrico completo, toda bola cerrada es compacta.\hfill{\footnotesize\hyperlink{r-17}{ver respuesta}}
\item[\hyperlink{r-18}{\textbf{18.}}] \hypertarget{q-18}{}{\small \textsc{Punto fijo de Banach} · Clase \href{https://decostruttivismo.github.io/MAT2514-repaso/pdf/MAT2514_apuntes_completos.pdf\#nameddest=section.31}{31}}\hfill \textbf{V}\,\fbox{\rule{0pt}{6pt}\hspace{6pt}}\quad\textbf{F}\,\fbox{\rule{0pt}{6pt}\hspace{6pt}}\par Si $X$ es completo y $T\colon X\to X$ cumple $d(Tx,Ty)\lt d(x,y)$ para todo $x\ne y$, entonces $T$ tiene un punto fijo.\hfill{\footnotesize\hyperlink{r-18}{ver respuesta}}
\item[\hyperlink{r-19}{\textbf{19.}}] \hypertarget{q-19}{}{\small \textsc{Funcionales lineales} · Clases \href{https://decostruttivismo.github.io/MAT2514-repaso/pdf/MAT2514_apuntes_completos.pdf\#nameddest=section.32}{32}, \href{https://decostruttivismo.github.io/MAT2514-repaso/pdf/MAT2514_apuntes_completos.pdf\#nameddest=section.34}{34}}\hfill \textbf{V}\,\fbox{\rule{0pt}{6pt}\hspace{6pt}}\quad\textbf{F}\,\fbox{\rule{0pt}{6pt}\hspace{6pt}}\par Si $t_1,\dots,t_m\in[0,1]$ son distintos, el funcional $\Lambda f=\sum_ic_if(t_i)$ tiene norma $\sum_i\abs{c_i}$ en $C([0,1],\R)$.\hfill{\footnotesize\hyperlink{r-19}{ver respuesta}}
\item[\hyperlink{r-20}{\textbf{20.}}] \hypertarget{q-20}{}{\small \textsc{Hahn–Banach} · Clase \href{https://decostruttivismo.github.io/MAT2514-repaso/pdf/MAT2514_apuntes_completos.pdf\#nameddest=section.35}{35}}\hfill \textbf{V}\,\fbox{\rule{0pt}{6pt}\hspace{6pt}}\quad\textbf{F}\,\fbox{\rule{0pt}{6pt}\hspace{6pt}}\par La distancia de $f\in C([0,1],\R)$ al subespacio $M$ de las funciones constantes es $\norm f_\infty$.\hfill{\footnotesize\hyperlink{r-20}{ver respuesta}}
\item[\hyperlink{r-21}{\textbf{21.}}] \hypertarget{q-21}{}{\small \textsc{Hahn–Banach} · Clase \href{https://decostruttivismo.github.io/MAT2514-repaso/pdf/MAT2514_apuntes_completos.pdf\#nameddest=section.35}{35}}\hfill \textbf{V}\,\fbox{\rule{0pt}{6pt}\hspace{6pt}}\quad\textbf{F}\,\fbox{\rule{0pt}{6pt}\hspace{6pt}}\par El corolario de distancia de Hahn–Banach se aplica a cualquier subespacio $M$ y cualquier $x_0\notin M$.\hfill{\footnotesize\hyperlink{r-21}{ver respuesta}}
\item[\hyperlink{r-22}{\textbf{22.}}] \hypertarget{q-22}{}{\small \textsc{Funcionales lineales} · Clase \href{https://decostruttivismo.github.io/MAT2514-repaso/pdf/MAT2514_apuntes_completos.pdf\#nameddest=section.34}{34}}\hfill \textbf{V}\,\fbox{\rule{0pt}{6pt}\hspace{6pt}}\quad\textbf{F}\,\fbox{\rule{0pt}{6pt}\hspace{6pt}}\par En $c_{00}$ (sucesiones finalmente nulas, con $\norm\cdot_\infty$), la fórmula $f(x)=\sum_nn\,x_n$ no define un funcional lineal.\hfill{\footnotesize\hyperlink{r-22}{ver respuesta}}
\item[\hyperlink{r-23}{\textbf{23.}}] \hypertarget{q-23}{}{\small \textsc{Funcionales lineales} · Clase \href{https://decostruttivismo.github.io/MAT2514-repaso/pdf/MAT2514_apuntes_completos.pdf\#nameddest=section.34}{34}}\hfill \textbf{V}\,\fbox{\rule{0pt}{6pt}\hspace{6pt}}\quad\textbf{F}\,\fbox{\rule{0pt}{6pt}\hspace{6pt}}\par Un funcional lineal que no es continuo en $0$ no es continuo en ningún punto.\hfill{\footnotesize\hyperlink{r-23}{ver respuesta}}
\item[\hyperlink{r-24}{\textbf{24.}}] \hypertarget{q-24}{}{\small \textsc{Funcionales lineales} · Clase \href{https://decostruttivismo.github.io/MAT2514-repaso/pdf/MAT2514_apuntes_completos.pdf\#nameddest=section.34}{34}}\hfill \textbf{V}\,\fbox{\rule{0pt}{6pt}\hspace{6pt}}\quad\textbf{F}\,\fbox{\rule{0pt}{6pt}\hspace{6pt}}\par El dual $X^*$ es un espacio de Banach porque $X$ lo es.\hfill{\footnotesize\hyperlink{r-24}{ver respuesta}}
\item[\hyperlink{r-25}{\textbf{25.}}] \hypertarget{q-25}{}{\small \textsc{Gromov–Hausdorff} · Clase \href{https://decostruttivismo.github.io/MAT2514-repaso/pdf/MAT2514_apuntes_completos.pdf\#nameddest=section.40}{40}}\hfill \textbf{V}\,\fbox{\rule{0pt}{6pt}\hspace{6pt}}\quad\textbf{F}\,\fbox{\rule{0pt}{6pt}\hspace{6pt}}\par Para calcular $\dis R$ basta mirar pares $(x,y),(x',y')\in R$ con $x\ne x'$.\hfill{\footnotesize\hyperlink{r-25}{ver respuesta}}
\item[\hyperlink{r-26}{\textbf{26.}}] \hypertarget{q-26}{}{\small \textsc{Gromov–Hausdorff} · Clases \href{https://decostruttivismo.github.io/MAT2514-repaso/pdf/MAT2514_apuntes_completos.pdf\#nameddest=section.40}{40}, \href{https://decostruttivismo.github.io/MAT2514-repaso/pdf/MAT2514_apuntes_completos.pdf\#nameddest=section.39}{39}}\hfill \textbf{V}\,\fbox{\rule{0pt}{6pt}\hspace{6pt}}\quad\textbf{F}\,\fbox{\rule{0pt}{6pt}\hspace{6pt}}\par Para probar $d_{GH}(X,Y)\ge c$ basta exhibir una correspondencia $R$ con $\dis R\ge2c$.\hfill{\footnotesize\hyperlink{r-26}{ver respuesta}}
\item[\hyperlink{r-27}{\textbf{27.}}] \hypertarget{q-27}{}{\small \textsc{Gromov–Hausdorff} · Clases \href{https://decostruttivismo.github.io/MAT2514-repaso/pdf/MAT2514_apuntes_completos.pdf\#nameddest=section.39}{39}, \href{https://decostruttivismo.github.io/MAT2514-repaso/pdf/MAT2514_apuntes_completos.pdf\#nameddest=section.40}{40}}\hfill \textbf{V}\,\fbox{\rule{0pt}{6pt}\hspace{6pt}}\quad\textbf{F}\,\fbox{\rule{0pt}{6pt}\hspace{6pt}}\par Si $T$ son tres puntos a distancia mutua $1$ y $P$ son dos puntos a distancia $1$, la cota por diámetros da $d_{GH}(T,P)\ge\frac12$.\hfill{\footnotesize\hyperlink{r-27}{ver respuesta}}
\item[\hyperlink{r-28}{\textbf{28.}}] \hypertarget{q-28}{}{\small \textsc{Gromov–Hausdorff} · Clase \href{https://decostruttivismo.github.io/MAT2514-repaso/pdf/MAT2514_apuntes_completos.pdf\#nameddest=section.39}{39}}\hfill \textbf{V}\,\fbox{\rule{0pt}{6pt}\hspace{6pt}}\quad\textbf{F}\,\fbox{\rule{0pt}{6pt}\hspace{6pt}}\par Si $A,B$ son compactos de un mismo espacio $Z$, con la métrica restringida, entonces $d_{GH}(A,B)\le d_H(A,B)$.\hfill{\footnotesize\hyperlink{r-28}{ver respuesta}}
\item[\hyperlink{r-29}{\textbf{29.}}] \hypertarget{q-29}{}{\small \textsc{Gromov–Hausdorff} · Clase \href{https://decostruttivismo.github.io/MAT2514-repaso/pdf/MAT2514_apuntes_completos.pdf\#nameddest=section.41}{41}}\hfill \textbf{V}\,\fbox{\rule{0pt}{6pt}\hspace{6pt}}\quad\textbf{F}\,\fbox{\rule{0pt}{6pt}\hspace{6pt}}\par Si $X_n\xrightarrow{GH}X$ y todos los $X_n$ son disconexos (por ejemplo, finitos con al menos dos puntos), entonces $X$ es disconexo.\hfill{\footnotesize\hyperlink{r-29}{ver respuesta}}
\item[\hyperlink{r-30}{\textbf{30.}}] \hypertarget{q-30}{}{\small \textsc{Gromov–Hausdorff} · Clase \href{https://decostruttivismo.github.io/MAT2514-repaso/pdf/MAT2514_apuntes_completos.pdf\#nameddest=section.41}{41}}\hfill \textbf{V}\,\fbox{\rule{0pt}{6pt}\hspace{6pt}}\quad\textbf{F}\,\fbox{\rule{0pt}{6pt}\hspace{6pt}}\par Si $f\colon X\to Y$ cumple $\dis f\le\varepsilon$, entonces $d_{GH}(X,Y)\le\frac32\varepsilon$.\hfill{\footnotesize\hyperlink{r-30}{ver respuesta}}
\item[\hyperlink{r-31}{\textbf{31.}}] \hypertarget{q-31}{}{\small \textsc{Gromov–Hausdorff} · Clase \href{https://decostruttivismo.github.io/MAT2514-repaso/pdf/MAT2514_apuntes_completos.pdf\#nameddest=section.41}{41}}\hfill \textbf{V}\,\fbox{\rule{0pt}{6pt}\hspace{6pt}}\quad\textbf{F}\,\fbox{\rule{0pt}{6pt}\hspace{6pt}}\par Toda $\varepsilon$-isometría es continua.\hfill{\footnotesize\hyperlink{r-31}{ver respuesta}}
\item[\hyperlink{r-32}{\textbf{32.}}] \hypertarget{q-32}{}{\small \textsc{Gromov–Hausdorff} · Clases \href{https://decostruttivismo.github.io/MAT2514-repaso/pdf/MAT2514_apuntes_completos.pdf\#nameddest=section.41}{41}, \href{https://decostruttivismo.github.io/MAT2514-repaso/pdf/MAT2514_apuntes_completos.pdf\#nameddest=section.19}{19}}\hfill \textbf{V}\,\fbox{\rule{0pt}{6pt}\hspace{6pt}}\quad\textbf{F}\,\fbox{\rule{0pt}{6pt}\hspace{6pt}}\par Como $G_n\xrightarrow{GH}[0,1]$, el conjunto $\{[G_n]:n\in\N\}$ es compacto en $(\mathbb M,d_{GH})$.\hfill{\footnotesize\hyperlink{r-32}{ver respuesta}}
\item[\hyperlink{r-33}{\textbf{33.}}] \hypertarget{q-33}{}{\small \textsc{Gromov–Hausdorff} · Clases \href{https://decostruttivismo.github.io/MAT2514-repaso/pdf/MAT2514_apuntes_completos.pdf\#nameddest=section.41}{41}, \href{https://decostruttivismo.github.io/MAT2514-repaso/pdf/MAT2514_apuntes_completos.pdf\#nameddest=section.40}{40}}\hfill \textbf{V}\,\fbox{\rule{0pt}{6pt}\hspace{6pt}}\quad\textbf{F}\,\fbox{\rule{0pt}{6pt}\hspace{6pt}}\par Para compactos $X,Y$: $\abs{\diam X-\diam Y}\le2\,d_{GH}(X,Y)$.\hfill{\footnotesize\hyperlink{r-33}{ver respuesta}}
\end{enumerate}
\newpage\parte{Parte 2. Respuestas y justificaciones}
\needspace{6\baselineskip}
\subsection*{\hypertarget{r-1}{}1.\ \textcolor{bad}{Falso}}
{\small \textsc{Compacidad y distancia} · Clase \href{https://decostruttivismo.github.io/MAT2514-repaso/pdf/MAT2514_apuntes_completos.pdf\#nameddest=section.20}{20}}\par
\emph{Si $A\subseteq X$ es cerrado y no vacío, entonces para todo $x\in X$ existe $a_0\in A$ con $d(x,a_0)=d(x,A)$.}
\bloque{Justificación}
\begin{pasos}
\item[] La afirmación dice: cuando $A$ es cerrado, el ínfimo $d(x,A)=\inf_{a\in A}d(x,a)$ siempre se alcanza en algún punto $a_0\in A$.
\item[] Para ver que es falsa basta exhibir un espacio $X$, un cerrado no vacío $A\subseteq X$ y un punto $x\in X$ tales que ningún $a\in A$ cumple $d(x,a)=d(x,A)$.
\item[] \textbf{Contraejemplo.} Sea $X=\Q$ con la distancia $d(p,q)=\abs{p-q}$, sea $A=\{q\in\Q:\ q\ge\sqrt2\}$ y sea $x=0$.
\item[(1)] $A$ es no vacío: por ejemplo $2\in A$.
\item[(2)] $A$ es cerrado en $\Q$.
\begin{subpasos}
\item Sea $q\in\Q\setminus A$, es decir, $q\in\Q$ y $q\lt \sqrt2$. Sea $r=\sqrt2-q\gt 0$.
\item Si $p\in\Q$ cumple $\abs{p-q}\lt r$, entonces $p\lt q+r=\sqrt2$, luego $p\notin A$.
\item Así la bola de centro $q$ y radio $r$ en $\Q$ está contenida en $\Q\setminus A$: el complemento de $A$ es abierto en $\Q$.
\end{subpasos}
\item[(3)] $d(0,A)\ge\sqrt2$: cada $q\in A$ cumple $d(0,q)=q\ge\sqrt2$.
\item[(4)] $d(0,A)\le\sqrt2$.
\begin{subpasos}
\item Sea $\varepsilon\gt 0$. Por la densidad de $\Q$ en $\R$ (Clase 5) existe $q\in\Q$ con $\sqrt2\lt q\lt \sqrt2+\varepsilon$.
\item Ese $q$ está en $A$ y cumple $d(0,q)\lt \sqrt2+\varepsilon$. Como $\varepsilon$ es arbitrario, $d(0,A)\le\sqrt2$.
\end{subpasos}
\item[(5)] Por (3) y (4), $d(0,A)=\sqrt2$.
\item[(6)] Ningún $a\in A$ cumple $d(0,a)=\sqrt2$: eso diría $a=\sqrt2$, y $\sqrt2\notin\Q$ (Clase 1).
\item[] Por (1), (2), (5) y (6): $A$ es un cerrado no vacío de $\Q$ y la distancia $d(0,A)=\sqrt2$ no se alcanza en ningún punto de $A$. La afirmación es falsa.
\end{pasos}
\bloque{¿Y cuándo es verdadero?}
\begin{pasos}
\item[] Si $A$ es compacto y no vacío, la distancia $d(x,A)$ sí se alcanza (\href{https://decostruttivismo.github.io/MAT2514-repaso/pdf/MAT2514_apuntes_completos.pdf\#nameddest=teorema.20.16}{Teorema 20.16}).
\item[] La razón: la función $a\mapsto d(x,a)$ es continua sobre el compacto $A$.
\item[] Por el teorema de Weierstrass, esa función alcanza su mínimo en algún $a_0\in A$, y ese mínimo es $d(x,A)$.
\end{pasos}
\bloque{En la guía del examen}
\begin{pasos}
\item[] En el Problema \hyperlink{prob-1}{1}, antes de escribir \guillemotleft{}sea $a_0\in A$ el punto más cercano a $x$\guillemotright{} hay que decir por qué ese punto existe: porque $A$ es compacto.
\end{pasos}
\par\smallskip{\small \textbf{En los apuntes:} \href{https://decostruttivismo.github.io/MAT2514-repaso/pdf/MAT2514_apuntes_completos.pdf\#nameddest=teorema.20.16}{Teorema 20.16} · \href{https://decostruttivismo.github.io/MAT2514-repaso/pdf/MAT2514_apuntes_completos.pdf\#nameddest=teorema.20.10}{Observación 20.10} \quad \textbf{Guía:} Problema \hyperlink{prob-1}{1}}
\par{\small \hyperlink{e-1}{Ejemplo y ejercicio de esta pregunta} \quad·\quad \hyperlink{q-1}{volver a la afirmación}}
\needspace{6\baselineskip}
\subsection*{\hypertarget{r-2}{}2.\ \textcolor{bad}{Falso}}
{\small \textsc{Compacidad y distancia} · Clase \href{https://decostruttivismo.github.io/MAT2514-repaso/pdf/MAT2514_apuntes_completos.pdf\#nameddest=section.20}{20}}\par
\emph{Si $F_1,F_2\subseteq X$ son cerrados, no vacíos y disjuntos, entonces $d(F_1,F_2)\gt 0$.}
\bloque{Justificación}
\begin{pasos}
\item[] Recordamos la definición: $d(F_1,F_2)=\inf\{d(a,b):\ a\in F_1,\ b\in F_2\}$.
\item[] Que $F_1$ y $F_2$ sean disjuntos dice que cada número $d(a,b)$ es positivo.
\item[] Pero el ínfimo de un conjunto de números positivos puede ser $0$. Eso es lo que pasa en el ejemplo siguiente.
\item[] \textbf{Contraejemplo.} (\href{https://decostruttivismo.github.io/MAT2514-repaso/pdf/MAT2514_apuntes_completos.pdf\#nameddest=teorema.20.21}{Ejemplo 20.21}.) En $\R^2$ con la distancia euclidiana, sea $F_1=\{(t,0):t\in\R\}$ el eje horizontal y sea $F_2=\{(t,\frac1t):t\gt 0\}$ una rama de la hipérbola $y=\frac1x$.
\item[(1)] $F_1$ es cerrado: es el conjunto de ceros de la función continua $(x,y)\mapsto y$.
\item[(2)] $F_2$ es cerrado: es el conjunto de ceros de la función continua $(x,y)\mapsto xy-1$ intersecado con el cerrado $\{x\ge0\}$.
\item[(3)] $F_1$ y $F_2$ son disjuntos: los puntos de $F_1$ tienen segunda coordenada $0$, y los de $F_2$ tienen segunda coordenada $\frac1t\neq0$.
\item[(4)] Para cada $t\gt 0$, los puntos $a_t=(t,0)\in F_1$ y $b_t=(t,\frac1t)\in F_2$ están a distancia $\norm{a_t-b_t}=\frac1t$.
\item[(5)] Por (4), $d(F_1,F_2)\le\frac1t$ para todo $t\gt 0$. Haciendo $t\to\infty$, $d(F_1,F_2)=0$.
\item[] Por (1), (2), (3) y (5): $F_1$ y $F_2$ son cerrados, no vacíos y disjuntos, y sin embargo $d(F_1,F_2)=0$. La afirmación es falsa.
\end{pasos}
\bloque{¿Y cuándo es verdadero?}
\begin{pasos}
\item[] Si además uno de los dos conjuntos es compacto, la afirmación es verdadera (\href{https://decostruttivismo.github.io/MAT2514-repaso/pdf/MAT2514_apuntes_completos.pdf\#nameddest=teorema.20.20}{Teorema 20.20}):
\item[] si $K$ es compacto, $F$ es cerrado, ambos son no vacíos y $K\cap F=\varnothing$, entonces $d(K,F)\gt 0$.
\item[] Intuición: en el contraejemplo, $F_1$ y $F_2$ se acercan \guillemotleft{}en el infinito\guillemotright{}. Un compacto es acotado y no puede escaparse al infinito.
\end{pasos}
\bloque{En la guía del examen}
\begin{pasos}
\item[] En el Problema \hyperlink{prob-2}{2(a)} se usa este hecho con dos compactos para obtener $\eta\gt 0$.
\end{pasos}
\par\smallskip{\small \textbf{En los apuntes:} \href{https://decostruttivismo.github.io/MAT2514-repaso/pdf/MAT2514_apuntes_completos.pdf\#nameddest=teorema.20.20}{Teorema 20.20} · \href{https://decostruttivismo.github.io/MAT2514-repaso/pdf/MAT2514_apuntes_completos.pdf\#nameddest=teorema.20.21}{Ejemplo 20.21} \quad \textbf{Guía:} Problema \hyperlink{prob-2}{2}}
\par{\small \hyperlink{e-2}{Ejemplo y ejercicio de esta pregunta} \quad·\quad \hyperlink{q-2}{volver a la afirmación}}
\needspace{6\baselineskip}
\subsection*{\hypertarget{r-3}{}3.\ \textcolor{ok}{Verdadero}}
{\small \textsc{Compacidad y conexidad} · Clases \href{https://decostruttivismo.github.io/MAT2514-repaso/pdf/MAT2514_apuntes_completos.pdf\#nameddest=section.24}{24}, \href{https://decostruttivismo.github.io/MAT2514-repaso/pdf/MAT2514_apuntes_completos.pdf\#nameddest=section.20}{20}}\par
\emph{Si $K$ es compacto y disconexo, entonces $K=K_1\cup K_2$ con $K_1,K_2$ compactos, disjuntos y no vacíos; en particular $d(K_1,K_2)\gt 0$.}
\bloque{Justificación}
\begin{pasos}
\item[(1)] Que $K$ sea disconexo significa (\href{https://decostruttivismo.github.io/MAT2514-repaso/pdf/MAT2514_apuntes_completos.pdf\#nameddest=teorema.24.1}{Definición 24.1}): existen $U,V\subseteq K$ abiertos en $K$, disjuntos y no vacíos, con $U\cup V=K$.
\item[] Ponemos $K_1=U$ y $K_2=V$. Hay que ver que $K_1$ y $K_2$ son compactos y que $d(K_1,K_2)\gt 0$.
\item[(2)] Por (1), $K_1=K\setminus K_2$. Como $K_2$ es abierto en $K$, su complemento $K_1$ es cerrado en $K$.
\item[(3)] Del mismo modo, $K_2=K\setminus K_1$ es cerrado en $K$.
\item[(4)] Un subconjunto cerrado de un compacto es compacto (\href{https://decostruttivismo.github.io/MAT2514-repaso/pdf/MAT2514_apuntes_completos.pdf\#nameddest=teorema.20.3}{Corolario 20.3}). Por (2) y (3), $K_1$ y $K_2$ son compactos.
\item[(5)] Por (1), $K_1$ y $K_2$ son disjuntos y no vacíos.
\item[(6)] $K_2$ es compacto, luego cerrado en $X$.
\item[] Por (4), (5), (6) y el \href{https://decostruttivismo.github.io/MAT2514-repaso/pdf/MAT2514_apuntes_completos.pdf\#nameddest=teorema.20.20}{Teorema 20.20} (aplicado al compacto $K_1$ y al cerrado $K_2$): $d(K_1,K_2)\gt 0$. La afirmación es verdadera.
\end{pasos}
\bloque{En la guía del examen}
\begin{pasos}
\item[] Es el primer paso del Problema \hyperlink{prob-2}{2(a)}.
\item[] Se descuenta usar $\eta=d(K_1,K_2)\gt 0$ sin justificarlo, o justificarlo solo diciendo que $K_1$ y $K_2$ son cerrados disjuntos, que no basta (pregunta 2).
\end{pasos}
\par\smallskip{\small \textbf{En los apuntes:} \href{https://decostruttivismo.github.io/MAT2514-repaso/pdf/MAT2514_apuntes_completos.pdf\#nameddest=teorema.24.1}{Definición 24.1} · \href{https://decostruttivismo.github.io/MAT2514-repaso/pdf/MAT2514_apuntes_completos.pdf\#nameddest=teorema.20.3}{Corolario 20.3} · \href{https://decostruttivismo.github.io/MAT2514-repaso/pdf/MAT2514_apuntes_completos.pdf\#nameddest=teorema.20.20}{Teorema 20.20} \quad \textbf{Guía:} Problema \hyperlink{prob-2}{2}}
\par{\small \hyperlink{e-3}{Ejemplo y ejercicio de esta pregunta} \quad·\quad \hyperlink{q-3}{volver a la afirmación}}
\needspace{6\baselineskip}
\subsection*{\hypertarget{r-4}{}4.\ \textcolor{bad}{Falso}}
{\small \textsc{Compacidad y distancia} · Clases \href{https://decostruttivismo.github.io/MAT2514-repaso/pdf/MAT2514_apuntes_completos.pdf\#nameddest=section.18}{18}, \href{https://decostruttivismo.github.io/MAT2514-repaso/pdf/MAT2514_apuntes_completos.pdf\#nameddest=section.19}{19}}\par
\emph{Todo espacio métrico totalmente acotado es compacto.}
\bloque{Justificación}
\begin{pasos}
\item[] La caracterización correcta es (\href{https://decostruttivismo.github.io/MAT2514-repaso/pdf/MAT2514_apuntes_completos.pdf\#nameddest=teorema.19.6}{Corolario 19.6}): un espacio métrico es compacto si y solo si es completo y totalmente acotado.
\item[] La afirmación omite la completitud. Para ver que es falsa basta un espacio totalmente acotado que no sea compacto.
\item[] \textbf{Contraejemplo.} (\href{https://decostruttivismo.github.io/MAT2514-repaso/pdf/MAT2514_apuntes_completos.pdf\#nameddest=teorema.18.14}{Observación 18.14}.) El intervalo $(0,1)$ con la distancia $d(x,y)=\abs{x-y}$.
\item[(1)] $(0,1)$ es totalmente acotado.
\begin{subpasos}
\item Dado $\varepsilon\gt 0$, elegimos $n$ con $\frac1n\lt \varepsilon$.
\item Todo $x\in(0,1)$ está a distancia $\lt \frac1n$ de alguno de los puntos $\frac1{2n},\frac2{2n},\dots,\frac{2n-1}{2n}$, que están en $(0,1)$.
\item Esos $2n-1$ puntos forman una $\varepsilon$-red finita de $(0,1)$.
\end{subpasos}
\item[(2)] $(0,1)$ no es completo.
\begin{subpasos}
\item La sucesión $x_n=\frac1n$ ($n\ge2$) está en $(0,1)$ y es de Cauchy, porque converge a $0$ en $\R$.
\item Si convergiera a algún $x\in(0,1)$, por unicidad del límite en $\R$ sería $x=0$. Pero $0\notin(0,1)$.
\end{subpasos}
\item[(3)] Por (2) y el \href{https://decostruttivismo.github.io/MAT2514-repaso/pdf/MAT2514_apuntes_completos.pdf\#nameddest=teorema.19.6}{Corolario 19.6}, $(0,1)$ no es compacto.
\item[] Por (1) y (3), $(0,1)$ es totalmente acotado y no es compacto. La afirmación es falsa.
\end{pasos}
\bloque{¿Y cuándo es verdadero?}
\begin{pasos}
\item[] Si el espacio, además de totalmente acotado, es completo, entonces es compacto (\href{https://decostruttivismo.github.io/MAT2514-repaso/pdf/MAT2514_apuntes_completos.pdf\#nameddest=teorema.19.6}{Corolario 19.6}).
\end{pasos}
\bloque{En la guía del examen}
\begin{pasos}
\item[] En el Problema \hyperlink{prob-4}{4(c)} se descuenta concluir \guillemotleft{}compacto\guillemotright{} solo a partir de la acotación total, sin probar la completitud.
\end{pasos}
\par\smallskip{\small \textbf{En los apuntes:} \href{https://decostruttivismo.github.io/MAT2514-repaso/pdf/MAT2514_apuntes_completos.pdf\#nameddest=teorema.19.6}{Corolario 19.6} · \href{https://decostruttivismo.github.io/MAT2514-repaso/pdf/MAT2514_apuntes_completos.pdf\#nameddest=teorema.18.14}{Observación 18.14} \quad \textbf{Guía:} Problemas \hyperlink{prob-1}{1}, \hyperlink{prob-4}{4}}
\par{\small \hyperlink{e-4}{Ejemplo y ejercicio de esta pregunta} \quad·\quad \hyperlink{q-4}{volver a la afirmación}}
\needspace{6\baselineskip}
\subsection*{\hypertarget{r-5}{}5.\ \textcolor{ok}{Verdadero}}
{\small \textsc{Compacidad y Hausdorff} · Clases \href{https://decostruttivismo.github.io/MAT2514-repaso/pdf/MAT2514_apuntes_completos.pdf\#nameddest=section.18}{18}, \href{https://decostruttivismo.github.io/MAT2514-repaso/pdf/MAT2514_apuntes_completos.pdf\#nameddest=section.36}{36}}\par
\emph{Si $K$ es compacto, para cada $\varepsilon\gt 0$ tiene una $\varepsilon$-red finita $S\subseteq K$, y entonces $d_H(S,K)\le\varepsilon$.}
\bloque{Justificación}
\begin{pasos}
\item[(1)] $K$ es compacto, luego totalmente acotado (\href{https://decostruttivismo.github.io/MAT2514-repaso/pdf/MAT2514_apuntes_completos.pdf\#nameddest=teorema.19.6}{Corolario 19.6}).
\item[(2)] Por (1) y el \href{https://decostruttivismo.github.io/MAT2514-repaso/pdf/MAT2514_apuntes_completos.pdf\#nameddest=teorema.18.6}{Lema 18.6}, la red finita puede tomarse dentro de $K$: hay un conjunto finito $S\subseteq K$ tal que todo $x\in K$ está a distancia $\lt \varepsilon$ de algún punto de $S$.
\item[] Recordamos: $d_H(S,K)=\max\bigl\{\sup_{s\in S}d(s,K),\ \sup_{x\in K}d(x,S)\bigr\}$. Hay que acotar las dos desviaciones.
\item[(3)] Primera desviación: cada $s\in S$ está en $K$, luego $d(s,K)=0$. Así $\sup_{s\in S}d(s,K)=0$.
\item[(4)] Segunda desviación: por (2), cada $x\in K$ cumple $d(x,S)\lt \varepsilon$. Así $\sup_{x\in K}d(x,S)\le\varepsilon$.
\item[] Por (3) y (4), $d_H(S,K)\le\varepsilon$. La afirmación es verdadera.
\end{pasos}
\bloque{En la guía del examen}
\begin{pasos}
\item[] Es el Problema \hyperlink{prob-1}{1(b)}: tomando $\varepsilon=\frac1n$ se obtienen conjuntos finitos $S_n$ con $S_n\to K$ en $(\mathcal K(X),d_H)$.
\item[] Se descuenta usar redes que no son finitas.
\end{pasos}
\par\smallskip{\small \textbf{En los apuntes:} \href{https://decostruttivismo.github.io/MAT2514-repaso/pdf/MAT2514_apuntes_completos.pdf\#nameddest=teorema.18.5}{Definición 18.5} · \href{https://decostruttivismo.github.io/MAT2514-repaso/pdf/MAT2514_apuntes_completos.pdf\#nameddest=teorema.18.6}{Lema 18.6} · \href{https://decostruttivismo.github.io/MAT2514-repaso/pdf/MAT2514_apuntes_completos.pdf\#nameddest=teorema.36.1}{Definición 36.1} \quad \textbf{Guía:} Problema \hyperlink{prob-1}{1}}
\par{\small \hyperlink{e-5}{Ejemplo y ejercicio de esta pregunta} \quad·\quad \hyperlink{q-5}{volver a la afirmación}}
\needspace{6\baselineskip}
\subsection*{\hypertarget{r-6}{}6.\ \textcolor{ok}{Verdadero}}
{\small \textsc{Completitud} · Clases \href{https://decostruttivismo.github.io/MAT2514-repaso/pdf/MAT2514_apuntes_completos.pdf\#nameddest=section.15}{15}, \href{https://decostruttivismo.github.io/MAT2514-repaso/pdf/MAT2514_apuntes_completos.pdf\#nameddest=section.37}{37}}\par
\emph{Un subconjunto cerrado de un espacio métrico completo es completo.}
\bloque{Justificación}
\begin{pasos}
\item[] Sea $X$ completo y $F\subseteq X$ cerrado. Hay que ver que toda sucesión de Cauchy en $F$ converge a un punto de $F$.
\item[] Sea $(x_n)$ una sucesión de Cauchy con todos sus términos en $F$.
\item[(1)] La distancia de $F$ es la misma que la de $X$. Luego $(x_n)$ es de Cauchy en $X$.
\item[(2)] $X$ es completo. Por (1), existe $x\in X$ con $x_n\to x$.
\item[(3)] $F$ es cerrado y todos los $x_n$ están en $F$. Por la caracterización secuencial de los cerrados (\href{https://decostruttivismo.github.io/MAT2514-repaso/pdf/MAT2514_apuntes_completos.pdf\#nameddest=teorema.12.10}{Corolario 12.10}), el límite $x$ está en $F$.
\item[] Por (2) y (3), $(x_n)$ converge a un punto de $F$. Luego $F$ es completo (\href{https://decostruttivismo.github.io/MAT2514-repaso/pdf/MAT2514_apuntes_completos.pdf\#nameddest=teorema.15.12}{Proposición 15.12}(b)). La afirmación es verdadera.
\end{pasos}
\bloque{En la guía del examen}
\begin{pasos}
\item[] Se usa en el Problema \hyperlink{prob-2}{2(c)}. Primero se prueba que $\mathcal{CN}(X)$ es cerrado en $\mathcal K(X)$.
\item[] Después se cita que $\mathcal K(X)$ es completo (\href{https://decostruttivismo.github.io/MAT2514-repaso/pdf/MAT2514_apuntes_completos.pdf\#nameddest=teorema.37.1}{Teorema 37.1}).
\item[] Si falta una de las dos cosas, el paso queda incompleto.
\end{pasos}
\par\smallskip{\small \textbf{En los apuntes:} \href{https://decostruttivismo.github.io/MAT2514-repaso/pdf/MAT2514_apuntes_completos.pdf\#nameddest=teorema.15.12}{Proposición 15.12} · \href{https://decostruttivismo.github.io/MAT2514-repaso/pdf/MAT2514_apuntes_completos.pdf\#nameddest=teorema.37.1}{Teorema 37.1} \quad \textbf{Guía:} Problemas \hyperlink{prob-2}{2}, \hyperlink{prob-11}{11}}
\par{\small \hyperlink{e-6}{Ejemplo y ejercicio de esta pregunta} \quad·\quad \hyperlink{q-6}{volver a la afirmación}}
\needspace{6\baselineskip}
\subsection*{\hypertarget{r-7}{}7.\ \textcolor{bad}{Falso}}
{\small \textsc{Distancia de Hausdorff} · Clase \href{https://decostruttivismo.github.io/MAT2514-repaso/pdf/MAT2514_apuntes_completos.pdf\#nameddest=section.36}{36}}\par
\emph{Para calcular $d_H(A,B)$ basta calcular $\sup_{a\in A}d(a,B)$.}
\bloque{Justificación}
\begin{pasos}
\item[(1)] Por definición (\href{https://decostruttivismo.github.io/MAT2514-repaso/pdf/MAT2514_apuntes_completos.pdf\#nameddest=teorema.36.1}{Definición 36.1}), $d_H(A,B)=\max\bigl\{\sup_{a\in A}d(a,B),\ \sup_{b\in B}d(b,A)\bigr\}$.
\item[] Por (1), $d_H(A,B)$ tiene dos desviaciones: la de $A$ respecto de $B$ y la de $B$ respecto de $A$. Calcular solo una puede dar un valor demasiado chico.
\item[] \textbf{Contraejemplo.} Sean $A=\{0,1\}$ y $B=[0,1]$ en $\R$.
\item[(2)] Desviación de $A$: los puntos $0$ y $1$ están en $B$, luego $d(0,B)=d(1,B)=0$. Así $\sup_{a\in A}d(a,B)=0$.
\item[(3)] Desviación de $B$: para $x\in[0,1]$, el punto de $A$ más cercano es $0$ o $1$, y $d(x,A)=\min\{x,\,1-x\}$.
\item[(4)] La función $\min\{x,1-x\}$ alcanza su máximo en $x=\frac12$, donde vale $\frac12$. Así $\sup_{b\in B}d(b,A)=\frac12$.
\item[] Por (1), (2) y (4): $d_H(A,B)=\max\{0,\frac12\}=\frac12$. En cambio, la fórmula de la afirmación daría $0$. La afirmación es falsa.
\end{pasos}
\bloque{¿Y cuándo es verdadero?}
\begin{pasos}
\item[] Calcular una sola desviación basta solo cuando se sabe que la otra vale $0$.
\item[] Por ejemplo, si $B\subseteq A$, entonces $\sup_{b\in B}d(b,A)=0$ y $d_H(A,B)=\sup_{a\in A}d(a,B)$.
\item[] Aun así, en la redacción hay que decir explícitamente por qué la otra desviación es $0$.
\end{pasos}
\bloque{En la guía del examen}
\begin{pasos}
\item[] Es el Problema \hyperlink{prob-9}{9(a)}.
\item[] Olvidar una de las dos desviaciones es el error más frecuente en los Problemas \hyperlink{prob-1}{1}, 3 y \hyperlink{prob-9}{9}.
\end{pasos}
\par\smallskip{\small \textbf{En los apuntes:} \href{https://decostruttivismo.github.io/MAT2514-repaso/pdf/MAT2514_apuntes_completos.pdf\#nameddest=teorema.36.1}{Definición 36.1} · \href{https://decostruttivismo.github.io/MAT2514-repaso/pdf/MAT2514_apuntes_completos.pdf\#nameddest=teorema.36.5}{Proposición 36.5} \quad \textbf{Guía:} Problemas \hyperlink{prob-1}{1}, 3, \hyperlink{prob-9}{9}}
\par{\small \hyperlink{e-7}{Ejemplo y ejercicio de esta pregunta} \quad·\quad \hyperlink{q-7}{volver a la afirmación}}
\needspace{6\baselineskip}
\subsection*{\hypertarget{r-8}{}8.\ \textcolor{bad}{Falso}}
{\small \textsc{Distancia de Hausdorff} · Clase \href{https://decostruttivismo.github.io/MAT2514-repaso/pdf/MAT2514_apuntes_completos.pdf\#nameddest=section.36}{36}}\par
\emph{Si $C\subseteq D$, entonces $d_H(D,C)=0$.}
\bloque{Justificación}
\begin{pasos}
\item[(1)] La inclusión $C\subseteq D$ dice que cada $c\in C$ cumple $d(c,D)=0$. Luego la desviación $\sup_{c\in C}d(c,D)$ vale $0$.
\item[(2)] La inclusión no dice nada sobre la otra desviación, $\sup_{x\in D}d(x,C)$, que puede ser positiva.
\item[] \textbf{Contraejemplo.} (Problema \hyperlink{prob-9}{9(b)}.) En $\R^2$, sean $D=\{x:\norm x\le1\}$ el disco unitario cerrado y $C=\{x:\norm x=1\}$ la circunferencia unitaria. Se tiene $C\subseteq D$.
\item[] Calculamos $d(x,C)$ para $x\in D$.
\item[(3)] Si $x\ne0$, el punto $\frac x{\norm x}$ está en $C$ y su distancia a $x$ es $\bigl\lVert x-\frac x{\norm x}\bigr\rVert=1-\norm x$. Luego $d(x,C)\le1-\norm x$.
\item[(4)] Si $y\in C$, por la triangular inversa $\norm{x-y}\ge\norm y-\norm x=1-\norm x$. Luego $d(x,C)\ge1-\norm x$.
\item[(5)] Por (3) y (4), $d(x,C)=1-\norm x$ para $x\ne0$.
\item[(6)] Si $x=0$, todo $y\in C$ cumple $\norm{0-y}=1$. Luego $d(0,C)=1=1-\norm0$: la fórmula de (5) también vale en $x=0$.
\item[(7)] Por (5) y (6), $\sup_{x\in D}d(x,C)=\sup_{x\in D}(1-\norm x)=1$, alcanzado en el centro $x=0$.
\item[] Por (1) y (7): $d_H(D,C)=\max\{0,1\}=1\ne0$. La afirmación es falsa.
\end{pasos}
\bloque{¿Y cuándo es verdadero?}
\begin{pasos}
\item[] Si $C\subseteq D$, entonces $d_H(D,C)=0$ exactamente cuando además $D\subseteq\cl C$, es decir, cuando $C$ es denso en $D$.
\end{pasos}
\bloque{En la guía del examen}
\begin{pasos}
\item[] En el Problema \hyperlink{prob-9}{9(b)} hay que justificar la fórmula $d(x,C)=1-\norm x$ con las dos desigualdades (3) y (4).
\item[] Darla sin la triangular inversa es un descuento frecuente.
\end{pasos}
\par\smallskip{\small \textbf{En los apuntes:} \href{https://decostruttivismo.github.io/MAT2514-repaso/pdf/MAT2514_apuntes_completos.pdf\#nameddest=teorema.36.1}{Definición 36.1} \quad \textbf{Guía:} Problema \hyperlink{prob-9}{9}}
\par{\small \hyperlink{e-8}{Ejemplo y ejercicio de esta pregunta} \quad·\quad \hyperlink{q-8}{volver a la afirmación}}
\needspace{6\baselineskip}
\subsection*{\hypertarget{r-9}{}9.\ \textcolor{ok}{Verdadero}}
{\small \textsc{Distancia de Hausdorff} · Clase \href{https://decostruttivismo.github.io/MAT2514-repaso/pdf/MAT2514_apuntes_completos.pdf\#nameddest=section.36}{36}}\par
\emph{$d_H\bigl((0,1),[0,1]\bigr)=0$, aunque los conjuntos son distintos.}
\bloque{Justificación}
\begin{pasos}
\item[] Calculamos las dos desviaciones entre $A=(0,1)$ y $B=[0,1]$.
\item[(1)] Desviación de $A$: como $(0,1)\subseteq[0,1]$, cada $a\in A$ cumple $d(a,B)=0$. Vale $0$.
\item[(2)] Si $x\in(0,1)$, entonces $x\in A$ y $d(x,A)=0$.
\item[(3)] Para $x=0$: $d(0,A)=\inf_{a\in(0,1)}a=0$, porque los puntos $\frac1n\in A$ están a distancia $\frac1n$ de $0$.
\item[(4)] Para $x=1$: del mismo modo, los puntos $1-\frac1n\in A$ dan $d(1,A)=0$.
\item[(5)] Por (2), (3) y (4), la desviación de $B$ es $\sup_{x\in[0,1]}d(x,A)=0$.
\item[] Por (1) y (5): $d_H(A,B)=0$, aunque $A\ne B$. La afirmación es verdadera.
\end{pasos}
\bloque{En la guía del examen}
\begin{pasos}
\item[] En el Problema \hyperlink{prob-9}{9(c)} los ejemplos deben ser conjuntos acotados, para que $d_H$ sea finita.
\end{pasos}
\par\smallskip{\small \textbf{En los apuntes:} \href{https://decostruttivismo.github.io/MAT2514-repaso/pdf/MAT2514_apuntes_completos.pdf\#nameddest=teorema.36.3}{Teorema 36.3} · \href{https://decostruttivismo.github.io/MAT2514-repaso/pdf/MAT2514_apuntes_completos.pdf\#nameddest=teorema.36.4}{Observación 36.4} \quad \textbf{Guía:} Problema \hyperlink{prob-9}{9}}
\par{\small \hyperlink{e-9}{Ejemplo y ejercicio de esta pregunta} \quad·\quad \hyperlink{q-9}{volver a la afirmación}}
\needspace{6\baselineskip}
\subsection*{\hypertarget{r-10}{}10.\ \textcolor{ok}{Verdadero}}
{\small \textsc{Distancia de Hausdorff} · Clases \href{https://decostruttivismo.github.io/MAT2514-repaso/pdf/MAT2514_apuntes_completos.pdf\#nameddest=section.36}{36}, \href{https://decostruttivismo.github.io/MAT2514-repaso/pdf/MAT2514_apuntes_completos.pdf\#nameddest=section.20}{20}}\par
\emph{Para $A,B\in\mathcal K(Z)$, con $\varphi_A(z)=d(z,A)$, se tiene $\sup_{z\in Z}\abs{\varphi_A(z)-\varphi_B(z)}=d_H(A,B)$.}
\bloque{Justificación}
\begin{pasos}
\item[] Hay que probar dos desigualdades: \guillemotleft{}$\le$\guillemotright{} y \guillemotleft{}$\ge$\guillemotright{}.
\item[] <b>Desigualdad \guillemotleft{}$\le$\guillemotright{}.</b> Fijamos $z\in Z$.
\item[(1)] Sea $b\in B$. Por la triangular para la distancia a un conjunto, $d(z,A)\le d(z,b)+d(b,A)$.
\item[(2)] Como $b\in B$, se tiene $d(b,A)\le\sup_{b'\in B}d(b',A)\le d_H(A,B)$.
\item[(3)] Por (1) y (2), $d(z,A)\le d(z,b)+d_H(A,B)$ para todo $b\in B$.
\item[(4)] Tomando ínfimo en $b\in B$ en (3): $d(z,A)\le d(z,B)+d_H(A,B)$, es decir, $\varphi_A(z)-\varphi_B(z)\le d_H(A,B)$.
\item[(5)] Intercambiando $A$ y $B$ en (1)–(4): $\varphi_B(z)-\varphi_A(z)\le d_H(A,B)$.
\item[(6)] Por (4) y (5), $\abs{\varphi_A(z)-\varphi_B(z)}\le d_H(A,B)$ para todo $z\in Z$. Luego el supremo es $\le d_H(A,B)$.
\item[] <b>Desigualdad \guillemotleft{}$\ge$\guillemotright{}.</b> Llamamos $S=\sup_{z\in Z}\abs{\varphi_A(z)-\varphi_B(z)}$.
\item[(7)] Si $a\in A$, entonces $\varphi_A(a)=0$ y $\abs{\varphi_A(a)-\varphi_B(a)}=d(a,B)$. Luego $S\ge d(a,B)$.
\item[(8)] Tomando supremo en $a\in A$ en (7): $S\ge\sup_{a\in A}d(a,B)$.
\item[(9)] Del mismo modo, evaluando en puntos $b\in B$: $S\ge\sup_{b\in B}d(b,A)$.
\item[(10)] Por (8) y (9), $S\ge\max\{\sup_{a\in A}d(a,B),\sup_{b\in B}d(b,A)\}=d_H(A,B)$.
\item[] Por (6) y (10), $S=d_H(A,B)$. La afirmación es verdadera.
\end{pasos}
\bloque{En la guía del examen}
\begin{pasos}
\item[] Es el Problema \hyperlink{prob-4}{4(b)}.
\item[] El error típico es probar solo \guillemotleft{}$\le$\guillemotright{} y dar la igualdad por obvia.
\end{pasos}
\par\smallskip{\small \textbf{En los apuntes:} \href{https://decostruttivismo.github.io/MAT2514-repaso/pdf/MAT2514_apuntes_completos.pdf\#nameddest=teorema.36.1}{Definición 36.1} · \href{https://decostruttivismo.github.io/MAT2514-repaso/pdf/MAT2514_apuntes_completos.pdf\#nameddest=teorema.20.11}{Proposición 20.11} \quad \textbf{Guía:} Problema \hyperlink{prob-4}{4}}
\par{\small \hyperlink{e-10}{Ejemplo y ejercicio de esta pregunta} \quad·\quad \hyperlink{q-10}{volver a la afirmación}}
\needspace{6\baselineskip}
\subsection*{\hypertarget{r-11}{}11.\ \textcolor{bad}{Falso}}
{\small \textsc{Distancia de Hausdorff} · Clase \href{https://decostruttivismo.github.io/MAT2514-repaso/pdf/MAT2514_apuntes_completos.pdf\#nameddest=section.37}{37}}\par
\emph{La desigualdad del Problema \hyperlink{prob-10}{10(a)} es $d_H(A_1\cup A_2,\,B_1\cup B_2)\le d_H(A_1,B_1)+d_H(A_2,B_2)$.}
\bloque{Justificación}
\begin{pasos}
\item[] La desigualdad del Problema \hyperlink{prob-10}{10(a)} (\href{https://decostruttivismo.github.io/MAT2514-repaso/pdf/MAT2514_apuntes_completos.pdf\#nameddest=teorema.37.5}{Lema 37.5}(a)) lleva un máximo, no una suma:
\item[] \[d_H(A_1\cup A_2,\,B_1\cup B_2)\le\max\{d_H(A_1,B_1),\,d_H(A_2,B_2)\}.\]
\item[] La versión con suma también es verdadera (porque el máximo de dos números no negativos es menor o igual que su suma), pero es más débil y no sirve para la contracción. Probamos la versión con máximo.
\item[] Sea $\delta=\max\{d_H(A_1,B_1),d_H(A_2,B_2)\}$.
\item[(1)] Sea $a\in A_1$. Como $B_1\subseteq B_1\cup B_2$, se tiene $d(a,B_1\cup B_2)\le d(a,B_1)$.
\item[(2)] Como $a\in A_1$, se tiene $d(a,B_1)\le d_H(A_1,B_1)\le\delta$.
\item[(3)] Por (1) y (2), $d(a,B_1\cup B_2)\le\delta$. Lo mismo vale si $a\in A_2$, usando $B_2$.
\item[(4)] Por (3), la desviación de $A_1\cup A_2$ respecto de $B_1\cup B_2$ es $\le\delta$.
\item[(5)] Intercambiando los papeles de los $A_i$ y los $B_i$, la otra desviación también es $\le\delta$.
\item[] Por (4) y (5), $d_H(A_1\cup A_2,B_1\cup B_2)\le\delta$.
\end{pasos}
\bloque{¿Y cuándo es verdadero?}
\begin{pasos}
\item[] La versión con suma es verdadera, pero es la equivocada para este problema.
\item[] Con la suma, el operador de Hutchinson de $m$ contracciones de razón $k$ daría la cota $m\,k\,d_H(A,B)$ en vez de $k\,d_H(A,B)$.
\item[] Si $mk\ge1$, eso no prueba que sea una contracción.
\end{pasos}
\bloque{En la guía del examen}
\begin{pasos}
\item[] Es el Problema \hyperlink{prob-10}{10(a)}.
\end{pasos}
\par\smallskip{\small \textbf{En los apuntes:} \href{https://decostruttivismo.github.io/MAT2514-repaso/pdf/MAT2514_apuntes_completos.pdf\#nameddest=teorema.37.5}{Lema 37.5} · \href{https://decostruttivismo.github.io/MAT2514-repaso/pdf/MAT2514_apuntes_completos.pdf\#nameddest=teorema.37.6}{Teorema 37.6} \quad \textbf{Guía:} Problema \hyperlink{prob-10}{10}}
\par{\small \hyperlink{e-11}{Ejemplo y ejercicio de esta pregunta} \quad·\quad \hyperlink{q-11}{volver a la afirmación}}
\needspace{6\baselineskip}
\subsection*{\hypertarget{r-12}{}12.\ \textcolor{ok}{Verdadero}}
{\small \textsc{Distancia de Hausdorff} · Clases \href{https://decostruttivismo.github.io/MAT2514-repaso/pdf/MAT2514_apuntes_completos.pdf\#nameddest=section.37}{37}, \href{https://decostruttivismo.github.io/MAT2514-repaso/pdf/MAT2514_apuntes_completos.pdf\#nameddest=section.19}{19}}\par
\emph{Si $A_1\supseteq A_2\supseteq\cdots$ son compactos no vacíos y $A=\bigcap_nA_n$, entonces para cada $\varepsilon\gt 0$ hay un $N$ con $A_n\subseteq A^\varepsilon$ para todo $n\ge N$.}
\bloque{Justificación}
\begin{pasos}
\item[] Aquí $A^\varepsilon=\{x:\ d(x,A)\lt \varepsilon\}$ es el $\varepsilon$-engrosamiento de $A$, que es un conjunto abierto.
\item[(1)] $A$ es no vacío y compacto, por el \href{https://decostruttivismo.github.io/MAT2514-repaso/pdf/MAT2514_apuntes_completos.pdf\#nameddest=teorema.19.13}{Corolario 19.13} (compactos encajados no vacíos).
\item[] Fijamos $\varepsilon\gt 0$ y definimos los restos $R_n=A_n\setminus A^\varepsilon$.
\item[(2)] Cada $R_n$ es compacto: es la intersección de $A_n$ con el cerrado $X\setminus A^\varepsilon$, es decir, un cerrado dentro del compacto $A_n$.
\item[(3)] $R_1\supseteq R_2\supseteq\cdots$, porque $A_1\supseteq A_2\supseteq\cdots$.
\item[(4)] $\bigcap_nR_n=\bigl(\bigcap_nA_n\bigr)\setminus A^\varepsilon=A\setminus A^\varepsilon$.
\item[(5)] $A\subseteq A^\varepsilon$, porque cada $a\in A$ cumple $d(a,A)=0\lt \varepsilon$. Por (4), $\bigcap_nR_n=\varnothing$.
\item[(6)] Si todos los $R_n$ fueran no vacíos, por (2), (3) y el \href{https://decostruttivismo.github.io/MAT2514-repaso/pdf/MAT2514_apuntes_completos.pdf\#nameddest=teorema.19.13}{Corolario 19.13} su intersección sería no vacía. Eso contradice (5).
\item[(7)] Por (6), existe $N$ con $R_N=\varnothing$. Por (3), $R_n=\varnothing$ para todo $n\ge N$.
\item[] Por (7), $A_n\subseteq A^\varepsilon$ para todo $n\ge N$. La afirmación es verdadera.
\end{pasos}
\bloque{En la guía del examen}
\begin{pasos}
\item[] Es la parte difícil del Problema \hyperlink{prob-10}{10(b)}, y se usa en los Problemas \hyperlink{prob-11}{11} y \hyperlink{prob-12}{12}. La otra desviación es inmediata, porque $A\subseteq A_n$.
\item[] Se descuenta tomar un $N_x$ distinto para cada punto $x$ sin obtener un $N$ común, o dar el paso por obvio.
\end{pasos}
\par\smallskip{\small \textbf{En los apuntes:} \href{https://decostruttivismo.github.io/MAT2514-repaso/pdf/MAT2514_apuntes_completos.pdf\#nameddest=teorema.37.3}{Proposición 37.3} · \href{https://decostruttivismo.github.io/MAT2514-repaso/pdf/MAT2514_apuntes_completos.pdf\#nameddest=teorema.19.13}{Corolario 19.13} \quad \textbf{Guía:} Problemas \hyperlink{prob-10}{10}, \hyperlink{prob-11}{11}, \hyperlink{prob-12}{12}}
\par{\small \hyperlink{e-12}{Ejemplo y ejercicio de esta pregunta} \quad·\quad \hyperlink{q-12}{volver a la afirmación}}
\needspace{6\baselineskip}
\subsection*{\hypertarget{r-13}{}13.\ \textcolor{bad}{Falso}}
{\small \textsc{Distancia de Hausdorff} · Clase \href{https://decostruttivismo.github.io/MAT2514-repaso/pdf/MAT2514_apuntes_completos.pdf\#nameddest=section.36}{36}}\par
\emph{La sucesión creciente $[0,n]$ converge en $\mathcal K(\R)$ a $[0,\infty)$.}
\bloque{Justificación}
\begin{pasos}
\item[(1)] $[0,\infty)$ no es compacto (no es acotado). Luego no es un elemento de $\mathcal K(\R)$, y la sucesión no puede converger a él en $\mathcal K(\R)$.
\item[] Más aún, la sucesión $([0,n])_n$ no converge a ningún elemento de $\mathcal K(\R)$. Lo probamos.
\item[] Sea $B\in\mathcal K(\R)$. Como $B$ es acotado, existe $R\gt 0$ con $B\subseteq[-R,R]$.
\item[(2)] Para $n\gt R$, el punto $n\in[0,n]$ cumple $d(n,B)\ge n-R$, porque todo $b\in B$ cumple $b\le R$.
\item[(3)] Por (2), la desviación de $[0,n]$ respecto de $B$ es $\ge n-R$. Luego $d_H([0,n],B)\ge n-R$.
\item[(4)] Por (3), $d_H([0,n],B)\to\infty$. La sucesión no converge a $B$.
\item[] Por (1) y (4), la afirmación es falsa.
\end{pasos}
\bloque{¿Y cuándo es verdadero?}
\begin{pasos}
\item[] Una sucesión creciente de compactos contenida en un compacto fijo sí converge: su límite es la clausura de la unión.
\end{pasos}
\bloque{En la guía del examen}
\begin{pasos}
\item[] Es el Problema \hyperlink{prob-10}{10(c)}.
\end{pasos}
\par\smallskip{\small \textbf{En los apuntes:} \href{https://decostruttivismo.github.io/MAT2514-repaso/pdf/MAT2514_apuntes_completos.pdf\#nameddest=teorema.36.3}{Teorema 36.3} \quad \textbf{Guía:} Problema \hyperlink{prob-10}{10}}
\par{\small \hyperlink{e-13}{Ejemplo y ejercicio de esta pregunta} \quad·\quad \hyperlink{q-13}{volver a la afirmación}}
\needspace{6\baselineskip}
\subsection*{\hypertarget{r-14}{}14.\ \textcolor{bad}{Falso}}
{\small \textsc{Operador de Hutchinson} · Clase \href{https://decostruttivismo.github.io/MAT2514-repaso/pdf/MAT2514_apuntes_completos.pdf\#nameddest=section.37}{37}}\par
\emph{Si $W\bigl(\bigcap_nA_n\bigr)\subseteq\bigcap_nA_n$, entonces $\bigcap_nA_n$ es el atractor.}
\bloque{Justificación}
\begin{pasos}
\item[] Aquí $W$ es el operador de Hutchinson de un sistema de contracciones, $A_0$ es un compacto y $A_n=W^n(A_0)$.
\item[(1)] El atractor $K$ es el único compacto no vacío con $W(K)=K$ (\href{https://decostruttivismo.github.io/MAT2514-repaso/pdf/MAT2514_apuntes_completos.pdf\#nameddest=teorema.37.6}{Teorema 37.6}).
\item[(2)] La condición $W(A)\subseteq A$ (ser invariante) la cumplen muchos conjuntos, no solo el atractor.
\item[] \textbf{Contraejemplo.} Para el sistema de Cantor, $f_1(x)=\frac x3$, $f_2(x)=\frac x3+\frac23$ en $[0,1]$, tomamos $A=[0,1]$.
\item[(3)] $W([0,1])=[0,\frac13]\cup[\frac23,1]\subseteq[0,1]$: el conjunto $[0,1]$ es invariante.
\item[(4)] $W([0,1])\ne[0,1]$. Por (1), $[0,1]$ no es el atractor.
\item[] Por (3) y (4), ser invariante no basta para ser el atractor. La afirmación es falsa.
\end{pasos}
\bloque{¿Y cuándo es verdadero?}
\begin{pasos}
\item[] La identificación correcta usa dos límites y la unicidad del límite (\href{https://decostruttivismo.github.io/MAT2514-repaso/pdf/MAT2514_apuntes_completos.pdf\#nameddest=teorema.37.7}{Ejemplo 37.7}):
\item[(1)] Si $(A_n)$ es decreciente, entonces $A_n\to\bigcap_nA_n$ en $d_H$ (\href{https://decostruttivismo.github.io/MAT2514-repaso/pdf/MAT2514_apuntes_completos.pdf\#nameddest=teorema.37.3}{Proposición 37.3}).
\item[(2)] Por el \href{https://decostruttivismo.github.io/MAT2514-repaso/pdf/MAT2514_apuntes_completos.pdf\#nameddest=teorema.37.6}{Teorema 37.6}, $A_n=W^n(A_0)\to K$ en $d_H$.
\item[] Por (1), (2) y la unicidad del límite en el espacio métrico $(\mathcal K(X),d_H)$: $\bigcap_nA_n=K$.
\end{pasos}
\bloque{En la guía del examen}
\begin{pasos}
\item[] Es la vía que piden los Problemas \hyperlink{prob-11}{11(b)} y \hyperlink{prob-12}{12}.
\end{pasos}
\par\smallskip{\small \textbf{En los apuntes:} \href{https://decostruttivismo.github.io/MAT2514-repaso/pdf/MAT2514_apuntes_completos.pdf\#nameddest=teorema.37.6}{Teorema 37.6} · \href{https://decostruttivismo.github.io/MAT2514-repaso/pdf/MAT2514_apuntes_completos.pdf\#nameddest=teorema.37.3}{Proposición 37.3} · \href{https://decostruttivismo.github.io/MAT2514-repaso/pdf/MAT2514_apuntes_completos.pdf\#nameddest=teorema.37.7}{Ejemplo 37.7} \quad \textbf{Guía:} Problemas \hyperlink{prob-11}{11}, \hyperlink{prob-12}{12}}
\par{\small \hyperlink{e-14}{Ejemplo y ejercicio de esta pregunta} \quad·\quad \hyperlink{q-14}{volver a la afirmación}}
\needspace{6\baselineskip}
\subsection*{\hypertarget{r-15}{}15.\ \textcolor{bad}{Falso}}
{\small \textsc{Operador de Hutchinson} · Clase \href{https://decostruttivismo.github.io/MAT2514-repaso/pdf/MAT2514_apuntes_completos.pdf\#nameddest=section.37}{37}}\par
\emph{$d_H\bigl([0,1],\ [0,\frac14]\cup[\frac34,1]\bigr)=\frac12$.}
\bloque{Justificación}
\begin{pasos}
\item[] Sea $A_1=[0,\frac14]\cup[\frac34,1]$. El número $\frac12$ es el largo del hueco $(\frac14,\frac34)$; veamos que $d_H$ no es eso.
\item[(1)] Desviación de $A_1$: como $A_1\subseteq[0,1]$, vale $0$.
\item[(2)] Si $x\in A_1$, entonces $d(x,A_1)=0$.
\item[(3)] Si $x\in[\frac14,\frac34]$, los puntos de $A_1$ más cercanos son $\frac14$ y $\frac34$, y $d(x,A_1)=\min\{x-\frac14,\ \frac34-x\}$.
\item[(4)] La expresión de (3) es máxima en el centro del hueco, $x=\frac12$, donde vale $\frac14$.
\item[(5)] Por (2), (3) y (4), la desviación de $[0,1]$ es $\frac14$.
\item[] Por (1) y (5): $d_H=\frac14$, la mitad del hueco, y no $\frac12$. La afirmación es falsa.
\end{pasos}
\bloque{En la guía del examen}
\begin{pasos}
\item[] En el Problema \hyperlink{prob-11}{11}, con $k=\frac14$ y $d_H(A_0,A_1)=\frac14$, la cota a priori del \href{https://decostruttivismo.github.io/MAT2514-repaso/pdf/MAT2514_apuntes_completos.pdf\#nameddest=teorema.37.6}{Teorema 37.6} es
\item[] \[d_H(A_n,K)\le\frac{k^n}{1-k}\,d_H(A_0,A_1)\]
\item[] \[\qquad=\frac{(1/4)^n}{3/4}\cdot\frac14=\frac{4^{-n}}3.\]
\end{pasos}
\par\smallskip{\small \textbf{En los apuntes:} \href{https://decostruttivismo.github.io/MAT2514-repaso/pdf/MAT2514_apuntes_completos.pdf\#nameddest=teorema.37.7}{Ejemplo 37.7} · \href{https://decostruttivismo.github.io/MAT2514-repaso/pdf/MAT2514_apuntes_completos.pdf\#nameddest=teorema.37.6}{Teorema 37.6} \quad \textbf{Guía:} Problema \hyperlink{prob-11}{11}}
\par{\small \hyperlink{e-15}{Ejemplo y ejercicio de esta pregunta} \quad·\quad \hyperlink{q-15}{volver a la afirmación}}
\needspace{6\baselineskip}
\subsection*{\hypertarget{r-16}{}16.\ \textcolor{bad}{Falso}}
{\small \textsc{Punto fijo de Banach} · Clase \href{https://decostruttivismo.github.io/MAT2514-repaso/pdf/MAT2514_apuntes_completos.pdf\#nameddest=section.31}{31}}\par
\emph{La cota a posteriori $d(x_n,x^*)\le\frac k{1-k}d(x_n,x_{n-1})$ permite decidir de antemano cuántas iteraciones hacen falta.}
\bloque{Justificación}
\begin{pasos}
\item[(1)] La cota a posteriori (\href{https://decostruttivismo.github.io/MAT2514-repaso/pdf/MAT2514_apuntes_completos.pdf\#nameddest=teorema.31.4}{Corolario 31.4}) usa $x_n$ y $x_{n-1}$.
\item[(2)] Los valores $x_n$ y $x_{n-1}$ solo se conocen después de calcular $n$ iteraciones.
\item[(3)] Por (1) y (2), con la cota a posteriori no se puede saber de antemano qué $n$ basta: solo se puede verificar, una vez calculado $x_n$, si ya se alcanzó la tolerancia.
\item[] Por (3), la afirmación es falsa.
\end{pasos}
\bloque{¿Y cuándo es verdadero?}
\begin{pasos}
\item[] Para fijar $n$ antes de iterar se usa la cota a priori: $d(x_n,x^*)\le\frac{k^n}{1-k}\,d(x_1,x_0)$.
\item[] Esa cota solo usa $x_0$ y $x_1$, que se conocen desde el principio.
\item[] Receta: se calculan $k$ y $d(x_1,x_0)$, se exige $\frac{k^n}{1-k}\,d(x_1,x_0)\lt $ tolerancia, y se despeja el menor $n$ que sirve.
\end{pasos}
\bloque{En la guía del examen}
\begin{pasos}
\item[] Los apuntes lo resumen así (\href{https://decostruttivismo.github.io/MAT2514-repaso/pdf/MAT2514_apuntes_completos.pdf\#nameddest=teorema.31.8}{Observación 31.8}): la cota a priori sirve para garantizar antes de iterar; la cota a posteriori, para parar a tiempo.
\item[] En el Problema \hyperlink{prob-11}{11(c)} se pide la cota a priori.
\end{pasos}
\par\smallskip{\small \textbf{En los apuntes:} \href{https://decostruttivismo.github.io/MAT2514-repaso/pdf/MAT2514_apuntes_completos.pdf\#nameddest=teorema.31.4}{Corolario 31.4} · \href{https://decostruttivismo.github.io/MAT2514-repaso/pdf/MAT2514_apuntes_completos.pdf\#nameddest=teorema.31.7}{Ejemplo 31.7} \quad \textbf{Guía:} Problema \hyperlink{prob-11}{11}}
\par{\small \hyperlink{e-16}{Ejemplo y ejercicio de esta pregunta} \quad·\quad \hyperlink{q-16}{volver a la afirmación}}
\needspace{6\baselineskip}
\subsection*{\hypertarget{r-17}{}17.\ \textcolor{bad}{Falso}}
{\small \textsc{Compacidad} · Clases \href{https://decostruttivismo.github.io/MAT2514-repaso/pdf/MAT2514_apuntes_completos.pdf\#nameddest=section.19}{19}, \href{https://decostruttivismo.github.io/MAT2514-repaso/pdf/MAT2514_apuntes_completos.pdf\#nameddest=section.20}{20}}\par
\emph{En un espacio métrico completo, toda bola cerrada es compacta.}
\bloque{Justificación}
\begin{pasos}
\item[] Que toda bola cerrada sea compacta es el teorema de Heine–Borel (\href{https://decostruttivismo.github.io/MAT2514-repaso/pdf/MAT2514_apuntes_completos.pdf\#nameddest=teorema.19.9}{Teorema 19.9}). Vale en $\R^n$, no en un espacio completo cualquiera.
\item[] \textbf{Contraejemplo.} (\href{https://decostruttivismo.github.io/MAT2514-repaso/pdf/MAT2514_apuntes_completos.pdf\#nameddest=teorema.20.27}{Corolario 20.27}.) En $\ell^\infty$, que es completo, la bola cerrada $\bar B(0,1)=\{x:\norm x_\infty\le1\}$ no es compacta.
\item[] Sea $e^{(n)}\in\ell^\infty$ la sucesión con un $1$ en la posición $n$ y ceros en el resto.
\item[(1)] Cada $e^{(n)}$ cumple $\norm{e^{(n)}}_\infty=1$. Luego $e^{(n)}\in\bar B(0,1)$.
\item[(2)] Si $n\ne m$, entonces $e^{(n)}-e^{(m)}$ tiene un $1$ y un $-1$, y $d_\infty(e^{(n)},e^{(m)})=1$.
\item[(3)] Por (2), ninguna subsucesión de $(e^{(n)})_n$ es de Cauchy: dos términos distintos siempre están a distancia $1$.
\item[(4)] Por (3), ninguna subsucesión de $(e^{(n)})_n$ converge.
\item[] Por (1) y (4), $\bar B(0,1)$ no es secuencialmente compacta, luego no es compacta. La afirmación es falsa.
\end{pasos}
\bloque{¿Y cuándo es verdadero?}
\begin{pasos}
\item[] En $\R^n$ con la distancia euclidiana la afirmación es verdadera: por Heine–Borel, toda bola cerrada es cerrada y acotada, luego compacta.
\end{pasos}
\bloque{En la guía del examen}
\begin{pasos}
\item[] En el Problema \hyperlink{prob-12}{12(b)}, la compacidad de $\cl B(x_0,R)$ se justifica diciendo que se está en $\R^d$ y citando Heine–Borel.
\end{pasos}
\par\smallskip{\small \textbf{En los apuntes:} \href{https://decostruttivismo.github.io/MAT2514-repaso/pdf/MAT2514_apuntes_completos.pdf\#nameddest=teorema.19.9}{Teorema 19.9} · \href{https://decostruttivismo.github.io/MAT2514-repaso/pdf/MAT2514_apuntes_completos.pdf\#nameddest=teorema.20.26}{Ejemplo 20.26} · \href{https://decostruttivismo.github.io/MAT2514-repaso/pdf/MAT2514_apuntes_completos.pdf\#nameddest=teorema.20.27}{Corolario 20.27} \quad \textbf{Guía:} Problema \hyperlink{prob-12}{12}}
\par{\small \hyperlink{e-17}{Ejemplo y ejercicio de esta pregunta} \quad·\quad \hyperlink{q-17}{volver a la afirmación}}
\needspace{6\baselineskip}
\subsection*{\hypertarget{r-18}{}18.\ \textcolor{bad}{Falso}}
{\small \textsc{Punto fijo de Banach} · Clase \href{https://decostruttivismo.github.io/MAT2514-repaso/pdf/MAT2514_apuntes_completos.pdf\#nameddest=section.31}{31}}\par
\emph{Si $X$ es completo y $T\colon X\to X$ cumple $d(Tx,Ty)\lt d(x,y)$ para todo $x\ne y$, entonces $T$ tiene un punto fijo.}
\bloque{Justificación}
\begin{pasos}
\item[] El \href{https://decostruttivismo.github.io/MAT2514-repaso/pdf/MAT2514_apuntes_completos.pdf\#nameddest=teorema.31.3}{Teorema 31.3} pide una razón $k\lt 1$ común a todos los pares: $d(Tx,Ty)\le k\,d(x,y)$. La hipótesis de la afirmación es más débil.
\item[] \textbf{Contraejemplo.} (\href{https://decostruttivismo.github.io/MAT2514-repaso/pdf/MAT2514_apuntes_completos.pdf\#nameddest=teorema.31.6}{Observación 31.6}(c).) Sean $X=[0,\infty)$ y $T(x)=\sqrt{x^2+1}$.
\item[(1)] $X$ es completo, por ser cerrado en $\R$.
\item[(2)] $T$ manda $X$ en $X$: $T(x)\ge1\ge0$.
\item[(3)] Para $x\ne y$ en $X$, multiplicando y dividiendo por el conjugado: $\abs{T(x)-T(y)}=\abs{x-y}\cdot\dfrac{x+y}{\sqrt{x^2+1}+\sqrt{y^2+1}}$.
\item[(4)] Como $x\lt \sqrt{x^2+1}$ e $y\lt \sqrt{y^2+1}$, el cociente de (3) es $\lt 1$. Luego $\abs{T(x)-T(y)}\lt \abs{x-y}$.
\item[(5)] $T$ no tiene punto fijo: $T(x)=x$ daría $x^2+1=x^2$, que es imposible.
\item[] Por (1), (2), (4) y (5): se cumplen las hipótesis y no hay punto fijo. La afirmación es falsa.
\end{pasos}
\bloque{¿Y cuándo es verdadero?}
\begin{pasos}
\item[] Es verdadera si hay una razón $k\lt 1$ común: $d(Tx,Ty)\le k\,d(x,y)$ para todos $x,y$ (\href{https://decostruttivismo.github.io/MAT2514-repaso/pdf/MAT2514_apuntes_completos.pdf\#nameddest=teorema.31.3}{Teorema 31.3}).
\item[] También es verdadera si además $X$ es compacto (ver el ejercicio).
\end{pasos}
\par\smallskip{\small \textbf{En los apuntes:} \href{https://decostruttivismo.github.io/MAT2514-repaso/pdf/MAT2514_apuntes_completos.pdf\#nameddest=teorema.31.3}{Teorema 31.3} · \href{https://decostruttivismo.github.io/MAT2514-repaso/pdf/MAT2514_apuntes_completos.pdf\#nameddest=teorema.31.6}{Observación 31.6} \quad \textbf{Guía:} Problema \hyperlink{prob-11}{11}}
\par{\small \hyperlink{e-18}{Ejemplo y ejercicio de esta pregunta} \quad·\quad \hyperlink{q-18}{volver a la afirmación}}
\needspace{6\baselineskip}
\subsection*{\hypertarget{r-19}{}19.\ \textcolor{ok}{Verdadero}}
{\small \textsc{Funcionales lineales} · Clases \href{https://decostruttivismo.github.io/MAT2514-repaso/pdf/MAT2514_apuntes_completos.pdf\#nameddest=section.32}{32}, \href{https://decostruttivismo.github.io/MAT2514-repaso/pdf/MAT2514_apuntes_completos.pdf\#nameddest=section.34}{34}}\par
\emph{Si $t_1,\dots,t_m\in[0,1]$ son distintos, el funcional $\Lambda f=\sum_ic_if(t_i)$ tiene norma $\sum_i\abs{c_i}$ en $C([0,1],\R)$.}
\bloque{Justificación}
\begin{pasos}
\item[] Para calcular una norma hacen falta dos cosas: una cota $\abs{\Lambda f}\le C\norm f_\infty$ para toda $f$, y una función $g$ con $\norm g_\infty=1$ que alcance $\abs{\Lambda g}=C$.
\item[] \textbf{Cota.}
\item[(1)] Por la triangular, $\abs{\Lambda f}\le\sum_i\abs{c_i}\abs{f(t_i)}$.
\item[(2)] Cada $\abs{f(t_i)}\le\norm f_\infty$. Por (1), $\abs{\Lambda f}\le\bigl(\sum_i\abs{c_i}\bigr)\norm f_\infty$.
\item[(3)] Por (2), $\norm\Lambda\le\sum_i\abs{c_i}$.
\item[] \textbf{Alcance.} Ordenamos los puntos: $t_1\lt t_2\lt \dots\lt t_m$.
\item[(4)] Elegimos signos $\sigma_i\in\{1,-1\}$ con $\sigma_ic_i=\abs{c_i}$.
\item[(5)] Sea $g$ la poligonal que vale $\sigma_i$ en $t_i$, es lineal en cada $[t_i,t_{i+1}]$, y es constante en $[0,t_1]$ y en $[t_m,1]$.
\item[(6)] $g$ es continua: es lineal en cada tramo y los tramos coinciden en los nodos.
\item[(7)] $\norm g_\infty=1$: en cada tramo, $g$ está entre los valores en los extremos, que son $\pm1$.
\item[(8)] Por (4) y (5), $\Lambda g=\sum_ic_i\sigma_i=\sum_i\abs{c_i}$.
\item[] Por (3), (6), (7) y (8): $\norm\Lambda=\sum_i\abs{c_i}$. La afirmación es verdadera.
\end{pasos}
\bloque{En la guía del examen}
\begin{pasos}
\item[] En el Problema \hyperlink{prob-5}{5} se descuenta definir $g$ \guillemotleft{}por casos\guillemotright{} sin justificar que es continua, como en (6).
\end{pasos}
\par\smallskip{\small \textbf{En los apuntes:} \href{https://decostruttivismo.github.io/MAT2514-repaso/pdf/MAT2514_apuntes_completos.pdf\#nameddest=teorema.32.10}{Observación 32.10} · \href{https://decostruttivismo.github.io/MAT2514-repaso/pdf/MAT2514_apuntes_completos.pdf\#nameddest=teorema.34.5}{Ejemplo 34.5} \quad \textbf{Guía:} Problema \hyperlink{prob-5}{5}}
\par{\small \hyperlink{e-19}{Ejemplo y ejercicio de esta pregunta} \quad·\quad \hyperlink{q-19}{volver a la afirmación}}
\needspace{6\baselineskip}
\subsection*{\hypertarget{r-20}{}20.\ \textcolor{bad}{Falso}}
{\small \textsc{Hahn–Banach} · Clase \href{https://decostruttivismo.github.io/MAT2514-repaso/pdf/MAT2514_apuntes_completos.pdf\#nameddest=section.35}{35}}\par
\emph{La distancia de $f\in C([0,1],\R)$ al subespacio $M$ de las funciones constantes es $\norm f_\infty$.}
\bloque{Justificación}
\begin{pasos}
\item[] La distancia es $\dist(f,M)=\inf_{c\in\R}\norm{f-c}_\infty$. Con la constante $c=0$ se obtiene $\norm f_\infty$, pero otra constante puede dar menos.
\item[] Sean $m=\min f$ y $M_f=\max f$, que existen por el teorema de Weierstrass.
\item[(1)] Con $c=\frac{m+M_f}2$: cada valor $f(x)$ está en $[m,M_f]$, luego $\abs{f(x)-c}\le\frac{M_f-m}2$. Así $\dist(f,M)\le\frac{M_f-m}2$.
\item[(2)] Sea $c\in\R$ cualquiera. Entonces $\norm{f-c}_\infty\ge\max\{M_f-c,\ c-m\}$.
\item[(3)] Los dos números de (2) suman $M_f-m$, así que el mayor es al menos $\frac{M_f-m}2$. Luego $\dist(f,M)\ge\frac{M_f-m}2$.
\item[(4)] Por (1) y (3), $\dist(f,M)=\frac12(\max f-\min f)$.
\item[] Por ejemplo, para $f(x)=x$, (4) da $\dist(f,M)=\frac12$, mientras que $\norm f_\infty=1$. La afirmación es falsa.
\end{pasos}
\bloque{¿Y cuándo es verdadero?}
\begin{pasos}
\item[] Cuando $\max f=-\min f$: entonces la constante óptima de (1) es $c=0$, y $\dist(f,M)=\norm f_\infty$.
\end{pasos}
\bloque{En la guía del examen}
\begin{pasos}
\item[] En el Problema \hyperlink{prob-5}{5} se pide además el funcional del \href{https://decostruttivismo.github.io/MAT2514-repaso/pdf/MAT2514_apuntes_completos.pdf\#nameddest=teorema.35.7}{Corolario 35.7}.
\item[] Es $F=\frac12(\delta_{x^+}-\delta_{x^-})$, donde $x^+$ y $x^-$ son puntos en que $f$ alcanza su máximo y su mínimo.
\item[] $F$ anula las constantes, cumple $F(f)=\frac12(\max f-\min f)$ y tiene norma $1$ si $x^+\ne x^-$, es decir, si $f$ no es constante.
\end{pasos}
\par\smallskip{\small \textbf{En los apuntes:} \href{https://decostruttivismo.github.io/MAT2514-repaso/pdf/MAT2514_apuntes_completos.pdf\#nameddest=teorema.35.7}{Corolario 35.7} · \href{https://decostruttivismo.github.io/MAT2514-repaso/pdf/MAT2514_apuntes_completos.pdf\#nameddest=teorema.35.8}{Ejemplo 35.8} \quad \textbf{Guía:} Problema \hyperlink{prob-5}{5}}
\par{\small \hyperlink{e-20}{Ejemplo y ejercicio de esta pregunta} \quad·\quad \hyperlink{q-20}{volver a la afirmación}}
\needspace{6\baselineskip}
\subsection*{\hypertarget{r-21}{}21.\ \textcolor{bad}{Falso}}
{\small \textsc{Hahn–Banach} · Clase \href{https://decostruttivismo.github.io/MAT2514-repaso/pdf/MAT2514_apuntes_completos.pdf\#nameddest=section.35}{35}}\par
\emph{El corolario de distancia de Hahn–Banach se aplica a cualquier subespacio $M$ y cualquier $x_0\notin M$.}
\bloque{Justificación}
\begin{pasos}
\item[] El corolario (\href{https://decostruttivismo.github.io/MAT2514-repaso/pdf/MAT2514_apuntes_completos.pdf\#nameddest=teorema.35.7}{Corolario 35.7}) dice: si $M$ es un subespacio cerrado y $x_0\notin M$, hay $F\in X^*$ con $F|_M=0$, $\norm F=1$ y $F(x_0)=\dist(x_0,M)\gt 0$.
\item[] Sin la hipótesis \guillemotleft{}cerrado\guillemotright{} la conclusión puede fallar.
\item[] \textbf{Contraejemplo.} Sea $M$ un subespacio no cerrado y sea $x_0\in\cl M\setminus M$.
\item[(1)] Como $x_0\in\cl M$, se tiene $\dist(x_0,M)=0$.
\item[(2)] Hay una sucesión $(m_n)$ en $M$ con $m_n\to x_0$.
\item[(3)] Sea $F\in X^*$ con $F|_M=0$. Por continuidad de $F$ y (2), $F(x_0)=\lim_nF(m_n)=0$.
\item[] Por (1) y (3): ningún $F$ continuo que anula $M$ cumple $F(x_0)\gt 0$, ni siquiera $F(x_0)\ne0$. La afirmación es falsa.
\end{pasos}
\bloque{¿Y cuándo es verdadero?}
\begin{pasos}
\item[] Si $M$ es cerrado, la afirmación es verdadera (\href{https://decostruttivismo.github.io/MAT2514-repaso/pdf/MAT2514_apuntes_completos.pdf\#nameddest=teorema.35.7}{Corolario 35.7}).
\item[] Si $M$ no es cerrado, se aplica el corolario a $\cl M$, que es un subespacio cerrado (ver el ejercicio), con $x_0\notin\cl M$.
\end{pasos}
\par\smallskip{\small \textbf{En los apuntes:} \href{https://decostruttivismo.github.io/MAT2514-repaso/pdf/MAT2514_apuntes_completos.pdf\#nameddest=teorema.35.7}{Corolario 35.7} \quad \textbf{Guía:} Problema \hyperlink{prob-5}{5}}
\par{\small \hyperlink{e-21}{Ejemplo y ejercicio de esta pregunta} \quad·\quad \hyperlink{q-21}{volver a la afirmación}}
\needspace{6\baselineskip}
\subsection*{\hypertarget{r-22}{}22.\ \textcolor{bad}{Falso}}
{\small \textsc{Funcionales lineales} · Clase \href{https://decostruttivismo.github.io/MAT2514-repaso/pdf/MAT2514_apuntes_completos.pdf\#nameddest=section.34}{34}}\par
\emph{En $c_{00}$ (sucesiones finalmente nulas, con $\norm\cdot_\infty$), la fórmula $f(x)=\sum_nn\,x_n$ no define un funcional lineal.}
\bloque{Justificación}
\begin{pasos}
\item[(1)] La fórmula está bien definida: si $x\in c_{00}$, solo finitos $x_n$ son distintos de $0$, y la suma $\sum_nn\,x_n$ es una suma finita.
\item[(2)] $f$ es lineal: $f(x+\lambda y)=\sum_nn(x_n+\lambda y_n)=f(x)+\lambda f(y)$.
\item[] Por (1) y (2), $f$ sí es un funcional lineal. La afirmación es falsa.
\item[] Lo que falla es otra cosa: $f$ no es acotado.
\item[(3)] Sea $e^{(m)}$ la sucesión con un $1$ en el lugar $m$ y ceros en el resto. Está en $c_{00}$ y $\norm{e^{(m)}}_\infty=1$.
\item[(4)] $f(e^{(m)})=m$.
\item[(5)] Por (3) y (4), $\norm f\ge m$ para todo $m$. Luego $f$ no es acotado.
\item[(6)] Por (5) y el \href{https://decostruttivismo.github.io/MAT2514-repaso/pdf/MAT2514_apuntes_completos.pdf\#nameddest=teorema.34.3}{Teorema 34.3} (para funcionales lineales, continuo equivale a acotado), $f$ no es continuo en ningún punto.
\end{pasos}
\bloque{En la guía del examen}
\begin{pasos}
\item[] En el Problema \hyperlink{prob-6}{6} hay que separar las dos cosas: $f$ está bien definido y es lineal, pero no es acotado.
\end{pasos}
\par\smallskip{\small \textbf{En los apuntes:} \href{https://decostruttivismo.github.io/MAT2514-repaso/pdf/MAT2514_apuntes_completos.pdf\#nameddest=teorema.34.6}{Ejemplo 34.6} · \href{https://decostruttivismo.github.io/MAT2514-repaso/pdf/MAT2514_apuntes_completos.pdf\#nameddest=teorema.34.3}{Teorema 34.3} \quad \textbf{Guía:} Problema \hyperlink{prob-6}{6}}
\par{\small \hyperlink{e-22}{Ejemplo y ejercicio de esta pregunta} \quad·\quad \hyperlink{q-22}{volver a la afirmación}}
\needspace{6\baselineskip}
\subsection*{\hypertarget{r-23}{}23.\ \textcolor{ok}{Verdadero}}
{\small \textsc{Funcionales lineales} · Clase \href{https://decostruttivismo.github.io/MAT2514-repaso/pdf/MAT2514_apuntes_completos.pdf\#nameddest=section.34}{34}}\par
\emph{Un funcional lineal que no es continuo en $0$ no es continuo en ningún punto.}
\bloque{Justificación}
\begin{pasos}
\item[] Probamos el contrarrecíproco: si un funcional lineal $f$ es continuo en un punto $x_0$, entonces es continuo en $0$.
\item[] Sea $f$ continuo en $x_0$ y sea $h_n\to0$.
\item[(1)] $x_0+h_n\to x_0$.
\item[(2)] Por (1) y la continuidad de $f$ en $x_0$: $f(x_0+h_n)\to f(x_0)$.
\item[(3)] Por linealidad, $f(h_n)=f(x_0+h_n)-f(x_0)$.
\item[(4)] Por (2) y (3), $f(h_n)\to0=f(0)$. Luego $f$ es continuo en $0$.
\item[] Por (4), si $f$ no es continuo en $0$, no puede ser continuo en ningún $x_0$. La afirmación es verdadera.
\item[] Esto es parte del \href{https://decostruttivismo.github.io/MAT2514-repaso/pdf/MAT2514_apuntes_completos.pdf\#nameddest=teorema.34.3}{Teorema 34.3}: para $f$ lineal son equivalentes ser continuo en algún punto, continuo en $0$, continuo en todo $X$, acotado y Lipschitz.
\end{pasos}
\bloque{En la guía del examen}
\begin{pasos}
\item[] En el Problema \hyperlink{prob-6}{6(b)}, concluir solo \guillemotleft{}no es continuo en $0$\guillemotright{} pierde puntos: hay que citar esta equivalencia para concluir que no es continuo en ningún punto.
\end{pasos}
\par\smallskip{\small \textbf{En los apuntes:} \href{https://decostruttivismo.github.io/MAT2514-repaso/pdf/MAT2514_apuntes_completos.pdf\#nameddest=teorema.34.3}{Teorema 34.3} \quad \textbf{Guía:} Problema \hyperlink{prob-6}{6}}
\par{\small \hyperlink{e-23}{Ejemplo y ejercicio de esta pregunta} \quad·\quad \hyperlink{q-23}{volver a la afirmación}}
\needspace{6\baselineskip}
\subsection*{\hypertarget{r-24}{}24.\ \textcolor{bad}{Falso}}
{\small \textsc{Funcionales lineales} · Clase \href{https://decostruttivismo.github.io/MAT2514-repaso/pdf/MAT2514_apuntes_completos.pdf\#nameddest=section.34}{34}}\par
\emph{El dual $X^*$ es un espacio de Banach porque $X$ lo es.}
\bloque{Justificación}
\begin{pasos}
\item[] El \href{https://decostruttivismo.github.io/MAT2514-repaso/pdf/MAT2514_apuntes_completos.pdf\#nameddest=teorema.34.7}{Teorema 34.7} dice que $X^*$ es de Banach siempre, aunque $X$ no sea completo. Revisamos dónde se usa la completitud en la prueba.
\item[] Sea $(f_n)$ una sucesión de Cauchy en $X^*$ y sea $x\in X$.
\item[(1)] $\abs{f_n(x)-f_m(x)}\le\norm{f_n-f_m}\,\norm x$.
\item[(2)] Por (1), la sucesión de números reales $(f_n(x))_n$ es de Cauchy.
\item[(3)] $\R$ es completo. Por (2), existe $\lim_nf_n(x)$; lo llamamos $f(x)$.
\item[(4)] $f$ es lineal, porque cada $f_n$ lo es y el límite respeta sumas y productos por escalar.
\item[(5)] Se prueba después que $f$ es acotado y que $\norm{f_n-f}\to0$, usando solo (1) y (3).
\item[] Por (3) y (5): la única completitud que se usa es la de $\R$, nunca la de $X$. La afirmación es falsa.
\end{pasos}
\bloque{¿Y cuándo es verdadero?}
\begin{pasos}
\item[] Es verdadero que $X^*$ es de Banach, pero la razón es la completitud de $\R$, no la de $X$.
\end{pasos}
\bloque{En la guía del examen}
\begin{pasos}
\item[] En el Problema \hyperlink{prob-6}{6} se aplica a $X=c_{00}$, que no es completo: su dual sí es de Banach.
\end{pasos}
\par\smallskip{\small \textbf{En los apuntes:} \href{https://decostruttivismo.github.io/MAT2514-repaso/pdf/MAT2514_apuntes_completos.pdf\#nameddest=teorema.34.7}{Teorema 34.7} \quad \textbf{Guía:} Problema \hyperlink{prob-6}{6}}
\par{\small \hyperlink{e-24}{Ejemplo y ejercicio de esta pregunta} \quad·\quad \hyperlink{q-24}{volver a la afirmación}}
\needspace{6\baselineskip}
\subsection*{\hypertarget{r-25}{}25.\ \textcolor{bad}{Falso}}
{\small \textsc{Gromov–Hausdorff} · Clase \href{https://decostruttivismo.github.io/MAT2514-repaso/pdf/MAT2514_apuntes_completos.pdf\#nameddest=section.40}{40}}\par
\emph{Para calcular $\dis R$ basta mirar pares $(x,y),(x',y')\in R$ con $x\ne x'$.}
\bloque{Justificación}
\begin{pasos}
\item[(1)] Por definición (\href{https://decostruttivismo.github.io/MAT2514-repaso/pdf/MAT2514_apuntes_completos.pdf\#nameddest=teorema.40.1}{Definición 40.1}), $\dis R=\sup\bigl\{\abs{d_X(x,x')-d_Y(y,y')}:\ (x,y),(x',y')\in R\bigr\}$, sobre todos los pares de elementos de $R$.
\item[(2)] Por (1), también cuentan los pares con $x=x'$ e $y\ne y'$. Un tal par aporta $\abs{d_X(x,x)-d_Y(y,y')}=d_Y(y,y')$.
\item[] Ese término puede ser el más grande. Lo vemos en un ejemplo.
\item[] \textbf{Contraejemplo.} (\href{https://decostruttivismo.github.io/MAT2514-repaso/pdf/MAT2514_apuntes_completos.pdf\#nameddest=teorema.40.2}{Ejemplo 40.2}.) Sean $X=\{0,1,2\}$ e $Y=\{0,2\}$, con la distancia de $\R$, y $R=\{(0,0),(1,0),(1,2),(2,2)\}$.
\item[(3)] Los elementos $(1,0)$ y $(1,2)$ comparten la primera coordenada. Aportan $\abs{d(1,1)-d(0,2)}=\abs{0-2}=2$.
\item[(4)] Los otros cinco pares de elementos distintos de $R$ aportan $1,1,0,1,1$.
\item[] Por (3) y (4): $\dis R=2$, pero mirando solo pares con $x\ne x'$ se obtendría $1$. La afirmación es falsa.
\end{pasos}
\bloque{En la guía del examen}
\begin{pasos}
\item[] Olvidar los pares que comparten una coordenada es el error frecuente del Problema \hyperlink{prob-13}{13}.
\end{pasos}
\par\smallskip{\small \textbf{En los apuntes:} \href{https://decostruttivismo.github.io/MAT2514-repaso/pdf/MAT2514_apuntes_completos.pdf\#nameddest=teorema.40.1}{Definición 40.1} · \href{https://decostruttivismo.github.io/MAT2514-repaso/pdf/MAT2514_apuntes_completos.pdf\#nameddest=teorema.40.2}{Ejemplo 40.2} \quad \textbf{Guía:} Problema \hyperlink{prob-13}{13}}
\par{\small \hyperlink{e-25}{Ejemplo y ejercicio de esta pregunta} \quad·\quad \hyperlink{q-25}{volver a la afirmación}}
\needspace{6\baselineskip}
\subsection*{\hypertarget{r-26}{}26.\ \textcolor{bad}{Falso}}
{\small \textsc{Gromov–Hausdorff} · Clases \href{https://decostruttivismo.github.io/MAT2514-repaso/pdf/MAT2514_apuntes_completos.pdf\#nameddest=section.40}{40}, \href{https://decostruttivismo.github.io/MAT2514-repaso/pdf/MAT2514_apuntes_completos.pdf\#nameddest=section.39}{39}}\par
\emph{Para probar $d_{GH}(X,Y)\ge c$ basta exhibir una correspondencia $R$ con $\dis R\ge2c$.}
\bloque{Justificación}
\begin{pasos}
\item[(1)] Por el \href{https://decostruttivismo.github.io/MAT2514-repaso/pdf/MAT2514_apuntes_completos.pdf\#nameddest=teorema.40.4}{Teorema 40.4}, $d_{GH}(X,Y)=\frac12\inf_R\dis R$, donde el ínfimo recorre todas las correspondencias entre $X$ e $Y$.
\item[(2)] Por (1), una correspondencia concreta $R$ da una cota superior: $d_{GH}(X,Y)\le\frac12\dis R$.
\item[(3)] Por (1), para una cota inferior $d_{GH}(X,Y)\ge c$ hay que probar $\dis R\ge2c$ para todas las correspondencias $R$.
\item[] Por (2) y (3), exhibir una sola correspondencia con distorsión grande no da ninguna cota inferior (ver el ejercicio). La afirmación es falsa.
\end{pasos}
\bloque{¿Y cuándo es verdadero?}
\begin{pasos}
\item[] Es verdadera si la desigualdad $\dis R\ge2c$ se prueba para toda correspondencia $R$.
\item[] Las herramientas de los apuntes para eso son: la cota por diámetros (\href{https://decostruttivismo.github.io/MAT2514-repaso/pdf/MAT2514_apuntes_completos.pdf\#nameddest=teorema.39.6}{Proposición 39.6}(c)), el conteo tipo palomar (\href{https://decostruttivismo.github.io/MAT2514-repaso/pdf/MAT2514_apuntes_completos.pdf\#nameddest=teorema.40.7}{Ejemplos 40.7} y 40.8) y, para ver que $d_{GH}\gt 0$, la rigidez.
\end{pasos}
\bloque{En la guía del examen}
\begin{pasos}
\item[] Es el error de fondo del Problema \hyperlink{prob-13}{13}.
\end{pasos}
\par\smallskip{\small \textbf{En los apuntes:} \href{https://decostruttivismo.github.io/MAT2514-repaso/pdf/MAT2514_apuntes_completos.pdf\#nameddest=teorema.40.4}{Teorema 40.4} · \href{https://decostruttivismo.github.io/MAT2514-repaso/pdf/MAT2514_apuntes_completos.pdf\#nameddest=teorema.40.7}{Ejemplo 40.7} \quad \textbf{Guía:} Problema \hyperlink{prob-13}{13}}
\par{\small \hyperlink{e-26}{Ejemplo y ejercicio de esta pregunta} \quad·\quad \hyperlink{q-26}{volver a la afirmación}}
\needspace{6\baselineskip}
\subsection*{\hypertarget{r-27}{}27.\ \textcolor{bad}{Falso}}
{\small \textsc{Gromov–Hausdorff} · Clases \href{https://decostruttivismo.github.io/MAT2514-repaso/pdf/MAT2514_apuntes_completos.pdf\#nameddest=section.39}{39}, \href{https://decostruttivismo.github.io/MAT2514-repaso/pdf/MAT2514_apuntes_completos.pdf\#nameddest=section.40}{40}}\par
\emph{Si $T$ son tres puntos a distancia mutua $1$ y $P$ son dos puntos a distancia $1$, la cota por diámetros da $d_{GH}(T,P)\ge\frac12$.}
\bloque{Justificación}
\begin{pasos}
\item[(1)] La cota por diámetros (\href{https://decostruttivismo.github.io/MAT2514-repaso/pdf/MAT2514_apuntes_completos.pdf\#nameddest=teorema.39.6}{Proposición 39.6}(c)) es $d_{GH}(T,P)\ge\frac12\abs{\diam T-\diam P}$.
\item[(2)] Todas las distancias entre puntos distintos valen $1$. Luego $\diam T=1$ y $\diam P=1$.
\item[(3)] Por (1) y (2), la cota da $d_{GH}(T,P)\ge\frac12\abs{1-1}=0$, que no dice nada.
\item[] Por (3), la cota por diámetros no da $\frac12$. La afirmación es falsa.
\end{pasos}
\bloque{¿Y cuándo es verdadero?}
\begin{pasos}
\item[] El valor $d_{GH}(T,P)\ge\frac12$ es verdadero, pero sale por conteo (palomar), como en el \href{https://decostruttivismo.github.io/MAT2514-repaso/pdf/MAT2514_apuntes_completos.pdf\#nameddest=teorema.40.7}{Ejemplo 40.7}.
\item[] La cota por diámetros solo sirve cuando los diámetros son distintos. La \href{https://decostruttivismo.github.io/MAT2514-repaso/pdf/MAT2514_apuntes_completos.pdf\#nameddest=teorema.39.11}{Observación 39.11} advierte que cuando coinciden no dicen nada.
\end{pasos}
\bloque{En la guía del examen}
\begin{pasos}
\item[] Es el Problema \hyperlink{prob-13}{13(b)}: el enunciado mismo pide observar que la cota por diámetros no sirve.
\end{pasos}
\par\smallskip{\small \textbf{En los apuntes:} \href{https://decostruttivismo.github.io/MAT2514-repaso/pdf/MAT2514_apuntes_completos.pdf\#nameddest=teorema.39.6}{Proposición 39.6} · \href{https://decostruttivismo.github.io/MAT2514-repaso/pdf/MAT2514_apuntes_completos.pdf\#nameddest=teorema.39.11}{Observación 39.11} \quad \textbf{Guía:} Problema \hyperlink{prob-13}{13}}
\par{\small \hyperlink{e-27}{Ejemplo y ejercicio de esta pregunta} \quad·\quad \hyperlink{q-27}{volver a la afirmación}}
\needspace{6\baselineskip}
\subsection*{\hypertarget{r-28}{}28.\ \textcolor{ok}{Verdadero}}
{\small \textsc{Gromov–Hausdorff} · Clase \href{https://decostruttivismo.github.io/MAT2514-repaso/pdf/MAT2514_apuntes_completos.pdf\#nameddest=section.39}{39}}\par
\emph{Si $A,B$ son compactos de un mismo espacio $Z$, con la métrica restringida, entonces $d_{GH}(A,B)\le d_H(A,B)$.}
\bloque{Justificación}
\begin{pasos}
\item[(1)] Por definición (\href{https://decostruttivismo.github.io/MAT2514-repaso/pdf/MAT2514_apuntes_completos.pdf\#nameddest=teorema.39.1}{Definición 39.1}), $d_{GH}(A,B)$ es el ínfimo de $d_H^W(\varphi(A),\psi(B))$ sobre todos los espacios métricos $W$ y todos los encajes isométricos $\varphi\colon A\to W$, $\psi\colon B\to W$.
\item[(2)] Tomamos $W=Z$, con $\varphi$ y $\psi$ las inclusiones de $A$ y de $B$ en $Z$.
\item[(3)] Las inclusiones de (2) son encajes isométricos, porque la métrica de $A$ y la de $B$ son la restricción de la métrica de $Z$.
\item[(4)] Para la elección de (2), el valor $d_H^W(\varphi(A),\psi(B))$ es $d_H(A,B)$.
\item[] Por (1), (3) y (4): el ínfimo es menor o igual que este valor, es decir, $d_{GH}(A,B)\le d_H(A,B)$ (\href{https://decostruttivismo.github.io/MAT2514-repaso/pdf/MAT2514_apuntes_completos.pdf\#nameddest=teorema.39.4}{Proposición 39.4}). La afirmación es verdadera.
\end{pasos}
\bloque{En la guía del examen}
\begin{pasos}
\item[] En el Problema \hyperlink{prob-14}{14(a)} hay que decir explícitamente cuál es el ambiente común $Z$ y por qué las inclusiones son encajes isométricos.
\end{pasos}
\par\smallskip{\small \textbf{En los apuntes:} \href{https://decostruttivismo.github.io/MAT2514-repaso/pdf/MAT2514_apuntes_completos.pdf\#nameddest=teorema.39.4}{Proposición 39.4} · \href{https://decostruttivismo.github.io/MAT2514-repaso/pdf/MAT2514_apuntes_completos.pdf\#nameddest=teorema.39.1}{Definición 39.1} \quad \textbf{Guía:} Problema \hyperlink{prob-14}{14}}
\par{\small \hyperlink{e-28}{Ejemplo y ejercicio de esta pregunta} \quad·\quad \hyperlink{q-28}{volver a la afirmación}}
\needspace{6\baselineskip}
\subsection*{\hypertarget{r-29}{}29.\ \textcolor{bad}{Falso}}
{\small \textsc{Gromov–Hausdorff} · Clase \href{https://decostruttivismo.github.io/MAT2514-repaso/pdf/MAT2514_apuntes_completos.pdf\#nameddest=section.41}{41}}\par
\emph{Si $X_n\xrightarrow{GH}X$ y todos los $X_n$ son disconexos (por ejemplo, finitos con al menos dos puntos), entonces $X$ es disconexo.}
\bloque{Justificación}
\begin{pasos}
\item[] La convergencia de Gromov–Hausdorff controla distancias, no la topología (\href{https://decostruttivismo.github.io/MAT2514-repaso/pdf/MAT2514_apuntes_completos.pdf\#nameddest=teorema.41.11}{Observación 41.11}).
\item[] \textbf{Contraejemplo.} (\href{https://decostruttivismo.github.io/MAT2514-repaso/pdf/MAT2514_apuntes_completos.pdf\#nameddest=teorema.41.7}{Ejemplo 41.7}(a).) Sea $X_n=\{0,\frac1n,\frac2n,\dots,1\}$, la grilla de paso $\frac1n$ en $[0,1]$, con la distancia de $\R$.
\item[(1)] Cada $X_n$ es finito y tiene al menos dos puntos. Luego es disconexo (ver el ejercicio).
\item[(2)] $X_n$ es una $\frac1{2n}$-red de $[0,1]$: todo $x\in[0,1]$ está a distancia $\le\frac1{2n}$ del punto de la grilla más cercano.
\item[(3)] Por (2) y la \href{https://decostruttivismo.github.io/MAT2514-repaso/pdf/MAT2514_apuntes_completos.pdf\#nameddest=teorema.39.4}{Proposición 39.4}, $d_{GH}(X_n,[0,1])\le\frac1{2n}\to0$. Luego $X_n\xrightarrow{GH}[0,1]$.
\item[(4)] $[0,1]$ es conexo (Clase 24).
\item[] Por (1), (3) y (4): una sucesión de espacios disconexos converge a un espacio conexo. La afirmación es falsa.
\end{pasos}
\bloque{¿Y cuándo es verdadero?}
\begin{pasos}
\item[] Lo que sí pasa al límite es el diámetro: $\abs{\diam X_n-\diam X}\le2\,d_{GH}(X_n,X)\to0$ (\href{https://decostruttivismo.github.io/MAT2514-repaso/pdf/MAT2514_apuntes_completos.pdf\#nameddest=teorema.41.9}{Proposición 41.9}).
\end{pasos}
\bloque{En la guía del examen}
\begin{pasos}
\item[] En los Problemas \hyperlink{prob-14}{14} y \hyperlink{prob-17}{17} no se puede argumentar con propiedades topológicas de los $X_n$ (conexidad, dimensión, cantidad de puntos).
\end{pasos}
\par\smallskip{\small \textbf{En los apuntes:} \href{https://decostruttivismo.github.io/MAT2514-repaso/pdf/MAT2514_apuntes_completos.pdf\#nameddest=teorema.41.11}{Observación 41.11} · \href{https://decostruttivismo.github.io/MAT2514-repaso/pdf/MAT2514_apuntes_completos.pdf\#nameddest=teorema.41.7}{Ejemplo 41.7} \quad \textbf{Guía:} Problemas \hyperlink{prob-14}{14}, \hyperlink{prob-17}{17}}
\par{\small \hyperlink{e-29}{Ejemplo y ejercicio de esta pregunta} \quad·\quad \hyperlink{q-29}{volver a la afirmación}}
\needspace{6\baselineskip}
\subsection*{\hypertarget{r-30}{}30.\ \textcolor{bad}{Falso}}
{\small \textsc{Gromov–Hausdorff} · Clase \href{https://decostruttivismo.github.io/MAT2514-repaso/pdf/MAT2514_apuntes_completos.pdf\#nameddest=section.41}{41}}\par
\emph{Si $f\colon X\to Y$ cumple $\dis f\le\varepsilon$, entonces $d_{GH}(X,Y)\le\frac32\varepsilon$.}
\bloque{Justificación}
\begin{pasos}
\item[(1)] La \href{https://decostruttivismo.github.io/MAT2514-repaso/pdf/MAT2514_apuntes_completos.pdf\#nameddest=teorema.41.3}{Proposición 41.3}(a) pide que $f$ sea una $\varepsilon$-isometría (\href{https://decostruttivismo.github.io/MAT2514-repaso/pdf/MAT2514_apuntes_completos.pdf\#nameddest=teorema.41.1}{Definición 41.1}): $\dis f\le\varepsilon$ y además $f(X)$ es una $\varepsilon$-red de $Y$.
\item[] Sin la condición de red, $f$ puede ver solo un trozo de $Y$.
\item[] \textbf{Contraejemplo.} Sean $X=\{0\}$, $Y=[0,1]$ y $f\colon X\to Y$ la inclusión, $f(0)=0$.
\item[(2)] $\dis f=0$: el único par de puntos de $X$ es $(0,0)$, que aporta $0$.
\item[(3)] $d_{GH}(\{0\},[0,1])=\frac12\diam[0,1]=\frac12$ (\href{https://decostruttivismo.github.io/MAT2514-repaso/pdf/MAT2514_apuntes_completos.pdf\#nameddest=teorema.39.6}{Proposición 39.6}(b)).
\item[(4)] Por (2), la afirmación con $\varepsilon=0$ daría $d_{GH}(\{0\},[0,1])\le0$, lo que contradice (3).
\item[] Por (4), la afirmación es falsa.
\end{pasos}
\bloque{¿Y cuándo es verdadero?}
\begin{pasos}
\item[] Si además $f(X)$ es una $\varepsilon$-red de $Y$, es decir, si $f$ es una $\varepsilon$-isometría (\href{https://decostruttivismo.github.io/MAT2514-repaso/pdf/MAT2514_apuntes_completos.pdf\#nameddest=teorema.41.3}{Proposición 41.3}(a)).
\end{pasos}
\bloque{En la guía del examen}
\begin{pasos}
\item[] En el Problema \hyperlink{prob-17}{17(b)} hay que verificar las dos condiciones de $\varepsilon$-isometría: la distorsión y la red.
\end{pasos}
\par\smallskip{\small \textbf{En los apuntes:} \href{https://decostruttivismo.github.io/MAT2514-repaso/pdf/MAT2514_apuntes_completos.pdf\#nameddest=teorema.41.1}{Definición 41.1} · \href{https://decostruttivismo.github.io/MAT2514-repaso/pdf/MAT2514_apuntes_completos.pdf\#nameddest=teorema.41.3}{Proposición 41.3} \quad \textbf{Guía:} Problema \hyperlink{prob-17}{17}}
\par{\small \hyperlink{e-30}{Ejemplo y ejercicio de esta pregunta} \quad·\quad \hyperlink{q-30}{volver a la afirmación}}
\needspace{6\baselineskip}
\subsection*{\hypertarget{r-31}{}31.\ \textcolor{bad}{Falso}}
{\small \textsc{Gromov–Hausdorff} · Clase \href{https://decostruttivismo.github.io/MAT2514-repaso/pdf/MAT2514_apuntes_completos.pdf\#nameddest=section.41}{41}}\par
\emph{Toda $\varepsilon$-isometría es continua.}
\bloque{Justificación}
\begin{pasos}
\item[(1)] La \href{https://decostruttivismo.github.io/MAT2514-repaso/pdf/MAT2514_apuntes_completos.pdf\#nameddest=teorema.41.1}{Definición 41.1} pide solo dos cosas: distorsión $\le\varepsilon$ y que la imagen sea una $\varepsilon$-red. No pide continuidad.
\item[] \textbf{Contraejemplo.} (\href{https://decostruttivismo.github.io/MAT2514-repaso/pdf/MAT2514_apuntes_completos.pdf\#nameddest=teorema.41.2}{Ejemplo 41.2}.) Sea $f\colon[0,1]\to X_n=\{0,\frac1n,\dots,1\}$, $f(x)=\lfloor nx\rfloor/n$, el redondeo hacia abajo a la grilla.
\item[(2)] $f$ es una $\frac1n$-isometría (ver el ejemplo).
\item[(3)] $f$ no es continua en $x=\frac1n$: a la izquierda vale $0$ y en $\frac1n$ vale $\frac1n$.
\item[] Por (2) y (3), hay $\varepsilon$-isometrías discontinuas. La afirmación es falsa.
\end{pasos}
\bloque{En la guía del examen}
\begin{pasos}
\item[] En el Problema \hyperlink{prob-17}{17} no hay que intentar probar la continuidad de $f_n$ ni usarla.
\end{pasos}
\par\smallskip{\small \textbf{En los apuntes:} \href{https://decostruttivismo.github.io/MAT2514-repaso/pdf/MAT2514_apuntes_completos.pdf\#nameddest=teorema.41.1}{Definición 41.1} · \href{https://decostruttivismo.github.io/MAT2514-repaso/pdf/MAT2514_apuntes_completos.pdf\#nameddest=teorema.41.2}{Ejemplo 41.2} \quad \textbf{Guía:} Problema \hyperlink{prob-17}{17}}
\par{\small \hyperlink{e-31}{Ejemplo y ejercicio de esta pregunta} \quad·\quad \hyperlink{q-31}{volver a la afirmación}}
\needspace{6\baselineskip}
\subsection*{\hypertarget{r-32}{}32.\ \textcolor{bad}{Falso}}
{\small \textsc{Gromov–Hausdorff} · Clases \href{https://decostruttivismo.github.io/MAT2514-repaso/pdf/MAT2514_apuntes_completos.pdf\#nameddest=section.41}{41}, \href{https://decostruttivismo.github.io/MAT2514-repaso/pdf/MAT2514_apuntes_completos.pdf\#nameddest=section.19}{19}}\par
\emph{Como $G_n\xrightarrow{GH}[0,1]$, el conjunto $\{[G_n]:n\in\N\}$ es compacto en $(\mathbb M,d_{GH})$.}
\bloque{Justificación}
\begin{pasos}
\item[] Aquí $G_n$ es la escalera del Problema \hyperlink{prob-17}{17}, $\mathbb M$ es el espacio de las clases de isometría de compactos, y $[G_n]$ es la clase de $G_n$.
\item[(1)] Por hipótesis, $[G_n]\to[[0,1]]$ en $(\mathbb M,d_{GH})$.
\item[(2)] Por (1), toda subsucesión de $([G_n])_n$ converge también a $[[0,1]]$.
\item[(3)] $[0,1]$ no es isométrico a ningún $G_n$: $[0,1]$ es infinito y cada $G_n$ es finito. Luego $[[0,1]]\notin\{[G_n]:n\in\N\}$.
\item[(4)] Por (2) y (3), por la unicidad del límite, ninguna subsucesión de $([G_n])_n$ converge a un elemento de $\{[G_n]:n\in\N\}$.
\item[] Por (4), el conjunto $\{[G_n]:n\in\N\}$ no es secuencialmente compacto, luego no es compacto. La afirmación es falsa.
\end{pasos}
\bloque{¿Y cuándo es verdadero?}
\begin{pasos}
\item[] Agregando el límite sí se obtiene un compacto: $\mathcal C=\{[G_n]:n\in\N\}\cup\{[[0,1]]\}$, una sucesión convergente junto con su límite (ver el ejercicio).
\end{pasos}
\bloque{En la guía del examen}
\begin{pasos}
\item[] Es el Problema \hyperlink{prob-17}{17(c)}.
\end{pasos}
\par\smallskip{\small \textbf{En los apuntes:} \href{https://decostruttivismo.github.io/MAT2514-repaso/pdf/MAT2514_apuntes_completos.pdf\#nameddest=teorema.41.5}{Definición 41.5} · \href{https://decostruttivismo.github.io/MAT2514-repaso/pdf/MAT2514_apuntes_completos.pdf\#nameddest=teorema.39.13}{Corolario 39.13} \quad \textbf{Guía:} Problema \hyperlink{prob-17}{17}}
\par{\small \hyperlink{e-32}{Ejemplo y ejercicio de esta pregunta} \quad·\quad \hyperlink{q-32}{volver a la afirmación}}
\needspace{6\baselineskip}
\subsection*{\hypertarget{r-33}{}33.\ \textcolor{ok}{Verdadero}}
{\small \textsc{Gromov–Hausdorff} · Clases \href{https://decostruttivismo.github.io/MAT2514-repaso/pdf/MAT2514_apuntes_completos.pdf\#nameddest=section.41}{41}, \href{https://decostruttivismo.github.io/MAT2514-repaso/pdf/MAT2514_apuntes_completos.pdf\#nameddest=section.40}{40}}\par
\emph{Para compactos $X,Y$: $\abs{\diam X-\diam Y}\le2\,d_{GH}(X,Y)$.}
\bloque{Justificación}
\begin{pasos}
\item[] Es la \href{https://decostruttivismo.github.io/MAT2514-repaso/pdf/MAT2514_apuntes_completos.pdf\#nameddest=teorema.41.9}{Proposición 41.9}. La probamos con correspondencias.
\item[] Sea $R$ una correspondencia entre $X$ e $Y$ y sea $\delta=\dis R$.
\item[(1)] Sean $x,x'\in X$. Como $R$ es correspondencia, existen $y,y'\in Y$ con $(x,y),(x',y')\in R$.
\item[(2)] Por definición de distorsión, $\abs{d_X(x,x')-d_Y(y,y')}\le\delta$. Luego $d_X(x,x')\le d_Y(y,y')+\delta$.
\item[(3)] $d_Y(y,y')\le\diam Y$. Por (2), $d_X(x,x')\le\diam Y+\delta$.
\item[(4)] Tomando supremo en $x,x'\in X$ en (3): $\diam X\le\diam Y+\delta$.
\item[(5)] Intercambiando los papeles de $X$ e $Y$ en (1)–(4): $\diam Y\le\diam X+\delta$.
\item[(6)] Por (4) y (5), $\abs{\diam X-\diam Y}\le\dis R$, para toda correspondencia $R$.
\item[(7)] Tomando ínfimo sobre $R$ en (6) y usando el \href{https://decostruttivismo.github.io/MAT2514-repaso/pdf/MAT2514_apuntes_completos.pdf\#nameddest=teorema.40.4}{Teorema 40.4}: $\abs{\diam X-\diam Y}\le\inf_R\dis R=2\,d_{GH}(X,Y)$.
\item[] La afirmación es verdadera.
\end{pasos}
\bloque{En la guía del examen}
\begin{pasos}
\item[] El error frecuente del Problema \hyperlink{prob-14}{14(c)} es probar (6) para una correspondencia y olvidar el paso (7): tomar el ínfimo.
\end{pasos}
\par\smallskip{\small \textbf{En los apuntes:} \href{https://decostruttivismo.github.io/MAT2514-repaso/pdf/MAT2514_apuntes_completos.pdf\#nameddest=teorema.41.9}{Proposición 41.9} · \href{https://decostruttivismo.github.io/MAT2514-repaso/pdf/MAT2514_apuntes_completos.pdf\#nameddest=teorema.40.4}{Teorema 40.4} \quad \textbf{Guía:} Problemas \hyperlink{prob-14}{14}, \hyperlink{prob-17}{17}}
\par{\small \hyperlink{e-33}{Ejemplo y ejercicio de esta pregunta} \quad·\quad \hyperlink{q-33}{volver a la afirmación}}
\newpage\parte{Parte 3. Ejemplos y ejercicios}
Para cada pregunta del quiz, un ejemplo tomado de los apuntes y un ejercicio. Las soluciones de los ejercicios están en la Parte 4.
\needspace{6\baselineskip}
\subsection*{\hypertarget{e-1}{}Pregunta 1 \normalfont\small(\textsc{Compacidad y distancia} · Clase \href{https://decostruttivismo.github.io/MAT2514-repaso/pdf/MAT2514_apuntes_completos.pdf\#nameddest=section.20}{20})}
\bloque{Ejemplo}
\textit{Un ínfimo que no es mínimo (Observación de la Clase 20)} {\small(en los apuntes: \href{https://decostruttivismo.github.io/MAT2514-repaso/pdf/MAT2514_apuntes_completos.pdf\#nameddest=teorema.20.10}{Observación 20.10})}
\begin{pasos}
\item[] En $X=\R$ sean $A=(0,1)$ y $x=0$.
\item[(1)] Para cada $a\in A$, $d(0,a)=a\gt 0$.
\item[(2)] Los puntos $a_n=\frac1n\in A$ (con $n\ge2$) cumplen $d(0,a_n)=\frac1n\to0$; luego $d(0,A)=0$.
\item[] Por (1) y (2), ningún $a\in A$ cumple $d(0,a)=d(0,A)$.
\item[] Aquí $A$ no es cerrado. El contraejemplo de la justificación muestra que ser cerrado tampoco basta.
\end{pasos}
\bloque{\hypertarget{x-1}{}Ejercicio 1}
(\href{https://decostruttivismo.github.io/MAT2514-repaso/pdf/MAT2514_apuntes_completos.pdf\#nameddest=ejercicio.A12B.6}{Ejercicio A12B.6}.) Sea $K\subseteq X$ compacto no vacío y $x\in X$. Demuestre que existe $y_0\in K$ con $d(x,y_0)=d(x,K)$.
\par{\small \hyperlink{s-1}{Ver solución} \quad·\quad \hyperlink{r-1}{volver a la respuesta de la pregunta 1}}
\needspace{6\baselineskip}
\subsection*{\hypertarget{e-2}{}Pregunta 2 \normalfont\small(\textsc{Compacidad y distancia} · Clase \href{https://decostruttivismo.github.io/MAT2514-repaso/pdf/MAT2514_apuntes_completos.pdf\#nameddest=section.20}{20})}
\bloque{Ejemplo}
\textit{\href{https://decostruttivismo.github.io/MAT2514-repaso/pdf/MAT2514_apuntes_completos.pdf\#nameddest=teorema.20.21}{Ejemplo 20.21}, mirado en el dibujo} {\small(en los apuntes: \href{https://decostruttivismo.github.io/MAT2514-repaso/pdf/MAT2514_apuntes_completos.pdf\#nameddest=teorema.20.21}{Ejemplo 20.21})}
\begin{pasos}
\item[] Los puntos $(t,0)$ y $(t,\frac1t)$ están uno exactamente encima del otro, separados por la altura $\frac1t$.
\item[] Al avanzar hacia la derecha ($t\to\infty$), esa altura tiende a $0$.
\item[] Por lo tanto la hipérbola se acerca al eje tanto como se quiera, sin tocarlo nunca.
\end{pasos}
\bloque{\hypertarget{x-2}{}Ejercicio 2}
Demuestre: si $K$ es compacto, $F$ es cerrado, ambos son no vacíos y $K\cap F=\varnothing$, entonces $d(K,F)\gt 0$.
\par{\small \hyperlink{s-2}{Ver solución} \quad·\quad \hyperlink{r-2}{volver a la respuesta de la pregunta 2}}
\needspace{6\baselineskip}
\subsection*{\hypertarget{e-3}{}Pregunta 3 \normalfont\small(\textsc{Compacidad y conexidad} · Clases \href{https://decostruttivismo.github.io/MAT2514-repaso/pdf/MAT2514_apuntes_completos.pdf\#nameddest=section.24}{24}, \href{https://decostruttivismo.github.io/MAT2514-repaso/pdf/MAT2514_apuntes_completos.pdf\#nameddest=section.20}{20})}
\bloque{Ejemplo}
\textit{Un caso concreto} {\small(en los apuntes: \href{https://decostruttivismo.github.io/MAT2514-repaso/pdf/MAT2514_apuntes_completos.pdf\#nameddest=teorema.24.4}{Observación 24.4})}
\begin{pasos}
\item[] Sea $K=[0,1]\cup[2,3]\subseteq\R$. Es compacto, por ser cerrado y acotado en $\R$.
\item[(1)] Los abiertos $U=(-1,\frac32)$ y $V=(\frac32,4)$ de $\R$ dan una disconexión de $K$ (\href{https://decostruttivismo.github.io/MAT2514-repaso/pdf/MAT2514_apuntes_completos.pdf\#nameddest=teorema.24.4}{Observación 24.4}).
\item[(2)] Por (1), $K_1=U\cap K=[0,1]$ y $K_2=V\cap K=[2,3]$.
\item[(3)] El par más cercano entre $K_1$ y $K_2$ es $1\in K_1$, $2\in K_2$. Luego $\eta=d(K_1,K_2)=1$.
\end{pasos}
\bloque{\hypertarget{x-3}{}Ejercicio 3}
Sea $K=[0,1]\cup[2,3]$ y sea $L\subseteq\R$ compacto con $d_H(L,K)\lt \frac13$. Pruebe que $L$ es disconexo.
\par{\small \hyperlink{s-3}{Ver solución} \quad·\quad \hyperlink{r-3}{volver a la respuesta de la pregunta 3}}
\needspace{6\baselineskip}
\subsection*{\hypertarget{e-4}{}Pregunta 4 \normalfont\small(\textsc{Compacidad y distancia} · Clases \href{https://decostruttivismo.github.io/MAT2514-repaso/pdf/MAT2514_apuntes_completos.pdf\#nameddest=section.18}{18}, \href{https://decostruttivismo.github.io/MAT2514-repaso/pdf/MAT2514_apuntes_completos.pdf\#nameddest=section.19}{19})}
\bloque{Ejemplo}
\textit{Las dos hipótesis son independientes (\href{https://decostruttivismo.github.io/MAT2514-repaso/pdf/MAT2514_apuntes_completos.pdf\#nameddest=teorema.18.14}{Observación 18.14})} {\small(en los apuntes: \href{https://decostruttivismo.github.io/MAT2514-repaso/pdf/MAT2514_apuntes_completos.pdf\#nameddest=teorema.18.14}{Observación 18.14})}
\begin{pasos}
\item[] $(0,1)$ es totalmente acotado y no es completo.
\item[] $\R$ es completo y no es totalmente acotado: ninguna familia finita de bolas de radio $1$ lo cubre, porque una unión finita de bolas es acotada.
\item[] $E=\{e^{(n)}:n\in\N\}\subseteq\ell^\infty$, el conjunto de las sucesiones con un $1$ en la posición $n$ y ceros en el resto, es completo y acotado, pero no es totalmente acotado.
\item[] Ninguno de los tres es compacto.
\end{pasos}
\bloque{\hypertarget{x-4}{}Ejercicio 4}
Sea $X$ completo y $A\subseteq X$ totalmente acotado. Pruebe que $\cl A$ es compacto.
\par{\small \hyperlink{s-4}{Ver solución} \quad·\quad \hyperlink{r-4}{volver a la respuesta de la pregunta 4}}
\needspace{6\baselineskip}
\subsection*{\hypertarget{e-5}{}Pregunta 5 \normalfont\small(\textsc{Compacidad y Hausdorff} · Clases \href{https://decostruttivismo.github.io/MAT2514-repaso/pdf/MAT2514_apuntes_completos.pdf\#nameddest=section.18}{18}, \href{https://decostruttivismo.github.io/MAT2514-repaso/pdf/MAT2514_apuntes_completos.pdf\#nameddest=section.36}{36})}
\bloque{Ejemplo}
\textit{La retícula del intervalo (\href{https://decostruttivismo.github.io/MAT2514-repaso/pdf/MAT2514_apuntes_completos.pdf\#nameddest=teorema.36.7}{Ejemplo 36.7})} {\small(en los apuntes: \href{https://decostruttivismo.github.io/MAT2514-repaso/pdf/MAT2514_apuntes_completos.pdf\#nameddest=teorema.36.7}{Ejemplo 36.7})}
\begin{pasos}
\item[] Sean $K=[0,1]$ y $F_n=\{\frac kn:\ 0\le k\le n\}$.
\item[(1)] Desviación de $F_n$: como $F_n\subseteq[0,1]$, vale $0$.
\item[(2)] Desviación de $[0,1]$: si $x\in[\frac kn,\frac{k+1}n]$, el punto de $F_n$ más cercano es uno de los dos extremos, y $d(x,F_n)=\min\{x-\frac kn,\ \frac{k+1}n-x\}$.
\item[(3)] Ese mínimo es a lo más $\frac1{2n}$, con igualdad en el punto medio del subintervalo.
\item[] Por (1), (2) y (3): $d_H([0,1],F_n)=\frac1{2n}$, que tiende a $0$.
\end{pasos}
\bloque{\hypertarget{x-5}{}Ejercicio 5}
Encuentre una $\frac14$-red finita de la circunferencia unitaria $S^1=\{x\in\R^2:\norm x=1\}$.
\par{\small \hyperlink{s-5}{Ver solución} \quad·\quad \hyperlink{r-5}{volver a la respuesta de la pregunta 5}}
\needspace{6\baselineskip}
\subsection*{\hypertarget{e-6}{}Pregunta 6 \normalfont\small(\textsc{Completitud} · Clases \href{https://decostruttivismo.github.io/MAT2514-repaso/pdf/MAT2514_apuntes_completos.pdf\#nameddest=section.15}{15}, \href{https://decostruttivismo.github.io/MAT2514-repaso/pdf/MAT2514_apuntes_completos.pdf\#nameddest=section.37}{37})}
\bloque{Ejemplo}
\textit{Ejemplos de la Clase 15}
\begin{pasos}
\item[] $[0,1]$ es completo, porque es cerrado en $\R$ y $\R$ es completo.
\item[] $(0,1)$ no es completo: $\frac1n$ es de Cauchy y su límite en $\R$, $0$, no está en $(0,1)$.
\item[] $\Q$ no es completo: una sucesión de racionales que converge a $\sqrt2$ en $\R$ es de Cauchy y no converge en $\Q$.
\item[] La esfera $S^{n-1}=\{x\in\R^n:\norm x_2=1\}$ es completa, porque es cerrada en $\R^n$.
\end{pasos}
\bloque{\hypertarget{x-6}{}Ejercicio 6}
Con la distancia de $\R$: ¿es completo $A=\{\frac1n:n\ge1\}\cup\{0\}$? ¿Y $B=\{\frac1n:n\ge1\}$?
\par{\small \hyperlink{s-6}{Ver solución} \quad·\quad \hyperlink{r-6}{volver a la respuesta de la pregunta 6}}
\needspace{6\baselineskip}
\subsection*{\hypertarget{e-7}{}Pregunta 7 \normalfont\small(\textsc{Distancia de Hausdorff} · Clase \href{https://decostruttivismo.github.io/MAT2514-repaso/pdf/MAT2514_apuntes_completos.pdf\#nameddest=section.36}{36})}
\bloque{Ejemplo}
\textit{Dos intervalos (\href{https://decostruttivismo.github.io/MAT2514-repaso/pdf/MAT2514_apuntes_completos.pdf\#nameddest=teorema.36.5}{Proposición 36.5})} {\small(en los apuntes: \href{https://decostruttivismo.github.io/MAT2514-repaso/pdf/MAT2514_apuntes_completos.pdf\#nameddest=teorema.36.5}{Proposición 36.5})}
\begin{pasos}
\item[] Sean $[a,b]$ y $[c,d]$ dos intervalos cerrados de $\R$.
\item[(1)] Para $x\in\R$, $d(x,[c,d])=\max\{c-x,\ 0,\ x-d\}$.
\item[(2)] Sobre $x\in[a,b]$, la expresión de (1) es máxima en $x=a$ o en $x=b$. Su máximo es $\max\{c-a,\ 0,\ b-d\}$.
\item[(3)] Intercambiando los papeles de los dos intervalos, la otra desviación es $\max\{a-c,\ 0,\ d-b\}$.
\item[(4)] Por (2) y (3): $d_H([a,b],[c,d])=\max\{\abs{a-c},\ \abs{b-d}\}$.
\item[] Por ejemplo, por (4), $d_H([0,1],[0,3])=\max\{0,2\}=2$.
\end{pasos}
\bloque{\hypertarget{x-7}{}Ejercicio 7}
Calcule $d_H(\{0,3\},[0,4])$.
\par{\small \hyperlink{s-7}{Ver solución} \quad·\quad \hyperlink{r-7}{volver a la respuesta de la pregunta 7}}
\needspace{6\baselineskip}
\subsection*{\hypertarget{e-8}{}Pregunta 8 \normalfont\small(\textsc{Distancia de Hausdorff} · Clase \href{https://decostruttivismo.github.io/MAT2514-repaso/pdf/MAT2514_apuntes_completos.pdf\#nameddest=section.36}{36})}
\bloque{Ejemplo}
\textit{Conjuntos que se tocan (\href{https://decostruttivismo.github.io/MAT2514-repaso/pdf/MAT2514_apuntes_completos.pdf\#nameddest=teorema.36.6}{Observación 36.6})} {\small(en los apuntes: \href{https://decostruttivismo.github.io/MAT2514-repaso/pdf/MAT2514_apuntes_completos.pdf\#nameddest=teorema.36.6}{Observación 36.6})}
\begin{pasos}
\item[] Sean $A=[0,1]$ y $B=[1,10]$ en $\R$.
\item[(1)] $A$ y $B$ tienen el punto $1$ en común. Luego $d(A,B)=0$.
\item[(2)] El punto $10\in B$ está a distancia $9$ de $A$.
\item[(3)] Por la \href{https://decostruttivismo.github.io/MAT2514-repaso/pdf/MAT2514_apuntes_completos.pdf\#nameddest=teorema.36.5}{Proposición 36.5}, $d_H(A,B)=\max\{\abs{0-1},\abs{1-10}\}=9$.
\item[] Por (1) y (3): $d(A,B)$ mide el par de puntos más cercano; $d_H(A,B)$ mide el punto peor ubicado.
\end{pasos}
\bloque{\hypertarget{x-8}{}Ejercicio 8}
Calcule $d_H([0,1],\{0\})$.
\par{\small \hyperlink{s-8}{Ver solución} \quad·\quad \hyperlink{r-8}{volver a la respuesta de la pregunta 8}}
\needspace{6\baselineskip}
\subsection*{\hypertarget{e-9}{}Pregunta 9 \normalfont\small(\textsc{Distancia de Hausdorff} · Clase \href{https://decostruttivismo.github.io/MAT2514-repaso/pdf/MAT2514_apuntes_completos.pdf\#nameddest=section.36}{36})}
\bloque{Ejemplo}
\textit{Racionales e intervalo (\href{https://decostruttivismo.github.io/MAT2514-repaso/pdf/MAT2514_apuntes_completos.pdf\#nameddest=teorema.36.4}{Observación 36.4})} {\small(en los apuntes: \href{https://decostruttivismo.github.io/MAT2514-repaso/pdf/MAT2514_apuntes_completos.pdf\#nameddest=teorema.36.4}{Observación 36.4})}
\begin{pasos}
\item[] Sean $A=\Q\cap[0,1]$ y $B=[0,1]$.
\item[(1)] $A\subseteq B$, así que la desviación de $A$ vale $0$.
\item[(2)] Todo $x\in[0,1]$ es límite de racionales de $[0,1]$, así que $d(x,A)=0$. La desviación de $B$ vale $0$.
\item[] Por (1) y (2), $d_H(A,B)=0$, aunque $A$ es numerable y $B$ no lo es. Los dos tienen la misma clausura.
\end{pasos}
\bloque{\hypertarget{x-9}{}Ejercicio 9}
Pruebe que si $A,B$ son acotados y $d_H(A,B)=0$, entonces $\cl A=\cl B$.
\par{\small \hyperlink{s-9}{Ver solución} \quad·\quad \hyperlink{r-9}{volver a la respuesta de la pregunta 9}}
\needspace{6\baselineskip}
\subsection*{\hypertarget{e-10}{}Pregunta 10 \normalfont\small(\textsc{Distancia de Hausdorff} · Clases \href{https://decostruttivismo.github.io/MAT2514-repaso/pdf/MAT2514_apuntes_completos.pdf\#nameddest=section.36}{36}, \href{https://decostruttivismo.github.io/MAT2514-repaso/pdf/MAT2514_apuntes_completos.pdf\#nameddest=section.20}{20})}
\bloque{Ejemplo}
\textit{Un caso que se puede dibujar}
\begin{pasos}
\item[] Sean $Z=[0,1]$, $A=\{0\}$ y $B=\{1\}$.
\item[(1)] $\varphi_A(z)=z$ y $\varphi_B(z)=1-z$.
\item[(2)] Por (1), $\abs{\varphi_A(z)-\varphi_B(z)}=\abs{2z-1}$, cuyo supremo en $[0,1]$ es $1$ (en $z=0$ y en $z=1$).
\item[(3)] Por otro lado, $d_H(\{0\},\{1\})=d(0,1)=1$.
\item[] Por (2) y (3), los dos lados de la igualdad valen $1$.
\end{pasos}
\bloque{\hypertarget{x-10}{}Ejercicio 10}
¿Qué propiedad común de las funciones $\varphi_A$ da la equicontinuidad en el Problema \hyperlink{prob-4}{4(a)}?
\par{\small \hyperlink{s-10}{Ver solución} \quad·\quad \hyperlink{r-10}{volver a la respuesta de la pregunta 10}}
\needspace{6\baselineskip}
\subsection*{\hypertarget{e-11}{}Pregunta 11 \normalfont\small(\textsc{Distancia de Hausdorff} · Clase \href{https://decostruttivismo.github.io/MAT2514-repaso/pdf/MAT2514_apuntes_completos.pdf\#nameddest=section.37}{37})}
\bloque{Ejemplo}
\textit{La diferencia entre suma y máximo}
\begin{pasos}
\item[] Sean $A_1=\{0\}$, $B_1=\{1\}$, $A_2=\{5\}$ y $B_2=\{6\}$ en $\R$.
\item[(1)] $d_H(A_1,B_1)=1$ y $d_H(A_2,B_2)=1$.
\item[(2)] En $A_1\cup A_2=\{0,5\}$ y $B_1\cup B_2=\{1,6\}$, cada punto está a distancia $1$ del otro conjunto. Luego $d_H(\{0,5\},\{1,6\})=1$.
\item[] Por (1) y (2): el valor es el máximo, $1$; la suma, $2$, es una cota mucho peor.
\end{pasos}
\bloque{\hypertarget{x-11}{}Ejercicio 11}
Pruebe que $d_H(A\cup B,\,A'\cup B)\le d_H(A,A')$.
\par{\small \hyperlink{s-11}{Ver solución} \quad·\quad \hyperlink{r-11}{volver a la respuesta de la pregunta 11}}
\needspace{6\baselineskip}
\subsection*{\hypertarget{e-12}{}Pregunta 12 \normalfont\small(\textsc{Distancia de Hausdorff} · Clases \href{https://decostruttivismo.github.io/MAT2514-repaso/pdf/MAT2514_apuntes_completos.pdf\#nameddest=section.37}{37}, \href{https://decostruttivismo.github.io/MAT2514-repaso/pdf/MAT2514_apuntes_completos.pdf\#nameddest=section.19}{19})}
\bloque{Ejemplo}
\textit{Un caso explícito}
\begin{pasos}
\item[] Sean $A_n=[0,\frac1n]$. Entonces $A=\bigcap_nA_n=\{0\}$.
\item[(1)] $A^\varepsilon=(-\varepsilon,\varepsilon)$.
\item[(2)] Por (1), $A_n\subseteq A^\varepsilon$ exactamente cuando $\frac1n\lt \varepsilon$.
\item[] Por (2), sirve cualquier $N\gt \frac1\varepsilon$.
\end{pasos}
\bloque{\hypertarget{x-12}{}Ejercicio 12}
¿Por qué el argumento falla para $A_n=[n,\infty)\subseteq\R$?
\par{\small \hyperlink{s-12}{Ver solución} \quad·\quad \hyperlink{r-12}{volver a la respuesta de la pregunta 12}}
\needspace{6\baselineskip}
\subsection*{\hypertarget{e-13}{}Pregunta 13 \normalfont\small(\textsc{Distancia de Hausdorff} · Clase \href{https://decostruttivismo.github.io/MAT2514-repaso/pdf/MAT2514_apuntes_completos.pdf\#nameddest=section.36}{36})}
\bloque{Ejemplo}
\textit{Una sucesión creciente que sí converge}
\begin{pasos}
\item[] Sean $B_n=\{0\}\cup\{1,\frac12,\dots,\frac1n\}$ y $B=\{0\}\cup\{\frac1k:k\ge1\}$. Todos están dentro de $[0,1]$, y $B$ es compacto.
\item[(1)] $B_n\subseteq B$, así que la desviación de $B_n$ vale $0$.
\item[(2)] Los puntos de $B$ que no están en $B_n$ son los $\frac1k$ con $k\ge n+1$. Cada uno está a distancia $\frac1k\le\frac1{n+1}$ de $0\in B_n$.
\item[(3)] Por (2), la desviación de $B$ es $\le\frac1{n+1}$.
\item[] Por (1) y (3), $d_H(B_n,B)\le\frac1{n+1}\to0$.
\end{pasos}
\bloque{\hypertarget{x-13}{}Ejercicio 13}
Calcule $d_H([0,n],[0,m])$ para $n\lt m$.
\par{\small \hyperlink{s-13}{Ver solución} \quad·\quad \hyperlink{r-13}{volver a la respuesta de la pregunta 13}}
\needspace{6\baselineskip}
\subsection*{\hypertarget{e-14}{}Pregunta 14 \normalfont\small(\textsc{Operador de Hutchinson} · Clase \href{https://decostruttivismo.github.io/MAT2514-repaso/pdf/MAT2514_apuntes_completos.pdf\#nameddest=section.37}{37})}
\bloque{Ejemplo}
\textit{El conjunto de Cantor como atractor (\href{https://decostruttivismo.github.io/MAT2514-repaso/pdf/MAT2514_apuntes_completos.pdf\#nameddest=teorema.37.7}{Ejemplo 37.7})} {\small(en los apuntes: \href{https://decostruttivismo.github.io/MAT2514-repaso/pdf/MAT2514_apuntes_completos.pdf\#nameddest=teorema.37.7}{Ejemplo 37.7})}
\begin{pasos}
\item[] Sean $f_1(x)=\frac x3$, $f_2(x)=\frac x3+\frac23$ en $[0,1]$, $C_0=[0,1]$ y $C_n=W^n(C_0)$.
\item[(1)] $C_1=[0,\frac13]\cup[\frac23,1]\subseteq C_0$.
\item[(2)] $W$ es monótono: si $A\subseteq B$, entonces $W(A)\subseteq W(B)$.
\item[(3)] Por (1), (2) e inducción, $(C_n)$ es decreciente.
\item[(4)] Por (3) y la \href{https://decostruttivismo.github.io/MAT2514-repaso/pdf/MAT2514_apuntes_completos.pdf\#nameddest=teorema.37.3}{Proposición 37.3}, $C_n\to C:=\bigcap_nC_n$, el conjunto de Cantor.
\item[(5)] Por el \href{https://decostruttivismo.github.io/MAT2514-repaso/pdf/MAT2514_apuntes_completos.pdf\#nameddest=teorema.37.6}{Teorema 37.6}, $C_n\to K$, el atractor.
\item[] Por (4), (5) y la unicidad del límite, $K=C$.
\end{pasos}
\bloque{\hypertarget{x-14}{}Ejercicio 14}
Sean $f_1(x)=\frac x4$, $f_2(x)=\frac x4+\frac34$, $A_0=[0,1]$ y $A_n=W^n(A_0)$. Pruebe que $(A_n)$ es decreciente.
\par{\small \hyperlink{s-14}{Ver solución} \quad·\quad \hyperlink{r-14}{volver a la respuesta de la pregunta 14}}
\needspace{6\baselineskip}
\subsection*{\hypertarget{e-15}{}Pregunta 15 \normalfont\small(\textsc{Operador de Hutchinson} · Clase \href{https://decostruttivismo.github.io/MAT2514-repaso/pdf/MAT2514_apuntes_completos.pdf\#nameddest=section.37}{37})}
\bloque{Ejemplo}
\textit{El mismo cálculo para Cantor (\href{https://decostruttivismo.github.io/MAT2514-repaso/pdf/MAT2514_apuntes_completos.pdf\#nameddest=teorema.37.7}{Ejemplo 37.7})} {\small(en los apuntes: \href{https://decostruttivismo.github.io/MAT2514-repaso/pdf/MAT2514_apuntes_completos.pdf\#nameddest=teorema.37.7}{Ejemplo 37.7})}
\begin{pasos}
\item[] Sean $C_0=[0,1]$ y $C_1=[0,\frac13]\cup[\frac23,1]$. El hueco $(\frac13,\frac23)$ mide $\frac13$.
\item[(1)] Para $x\in[\frac13,\frac23]$, $d(x,C_1)=\min\{x-\frac13,\ \frac23-x\}$, que es máximo en $x=\frac12$ y vale $\frac16$.
\item[(2)] Por (1), $d_H(C_0,C_1)=\frac16$.
\item[] Por (2), con $k=\frac13$, la cota a priori es $d_H(C_n,C)\le\frac{(1/3)^n}{2/3}\cdot\frac16=\frac{3^{-n}}4$.
\end{pasos}
\bloque{\hypertarget{x-15}{}Ejercicio 15}
En el Problema \hyperlink{prob-11}{11}, ¿qué $n$ garantiza $d_H(A_n,K)\lt 10^{-3}$?
\par{\small \hyperlink{s-15}{Ver solución} \quad·\quad \hyperlink{r-15}{volver a la respuesta de la pregunta 15}}
\needspace{6\baselineskip}
\subsection*{\hypertarget{e-16}{}Pregunta 16 \normalfont\small(\textsc{Punto fijo de Banach} · Clase \href{https://decostruttivismo.github.io/MAT2514-repaso/pdf/MAT2514_apuntes_completos.pdf\#nameddest=section.31}{31})}
\bloque{Ejemplo}
\textit{Una aplicación afín (\href{https://decostruttivismo.github.io/MAT2514-repaso/pdf/MAT2514_apuntes_completos.pdf\#nameddest=teorema.31.2}{Ejemplos 31.2} y 31.3)} {\small(en los apuntes: \href{https://decostruttivismo.github.io/MAT2514-repaso/pdf/MAT2514_apuntes_completos.pdf\#nameddest=teorema.31.5}{Ejemplo 31.5})}
\begin{pasos}
\item[] Sea $T(x)=\frac x3+2$ en $\R$, con $k=\frac13$, $x_0=0$ y $x_1=T(0)=2$.
\item[(1)] Cota a priori: $d(x_n,x^*)\le\frac{(1/3)^n}{2/3}\cdot2=3^{1-n}$.
\item[(2)] Para error $\lt 10^{-3}$ basta, por (1), que $3^{n-1}\gt 1000$.
\item[(3)] $3^6=729$ no cumple (2); $3^7=2187$ sí. Luego basta $n-1=7$, es decir, $n=8$.
\item[(4)] Cota a posteriori, ya calculado $x_8$: $d(x_8,x^*)\le\frac{1/3}{2/3}\abs{x_8-x_7}=\frac12\abs{x_8-x_7}$.
\end{pasos}
\bloque{\hypertarget{x-16}{}Ejercicio 16}
(\href{https://decostruttivismo.github.io/MAT2514-repaso/pdf/MAT2514_apuntes_completos.pdf\#nameddest=teorema.31.7}{Ejemplo 31.7}.) Sea $T(x)=\frac x2+\frac1x$ en $[1,2]$, con $k=\frac12$ y $x_0=1$. ¿Cuántas iteraciones garantizan error $\lt 10^{-6}$?
\par{\small \hyperlink{s-16}{Ver solución} \quad·\quad \hyperlink{r-16}{volver a la respuesta de la pregunta 16}}
\needspace{6\baselineskip}
\subsection*{\hypertarget{e-17}{}Pregunta 17 \normalfont\small(\textsc{Compacidad} · Clases \href{https://decostruttivismo.github.io/MAT2514-repaso/pdf/MAT2514_apuntes_completos.pdf\#nameddest=section.19}{19}, \href{https://decostruttivismo.github.io/MAT2514-repaso/pdf/MAT2514_apuntes_completos.pdf\#nameddest=section.20}{20})}
\bloque{Ejemplo}
\textit{Un cerrado y acotado no compacto en $\ell^\infty$ (\href{https://decostruttivismo.github.io/MAT2514-repaso/pdf/MAT2514_apuntes_completos.pdf\#nameddest=teorema.20.26}{Ejemplo 20.26})} {\small(en los apuntes: \href{https://decostruttivismo.github.io/MAT2514-repaso/pdf/MAT2514_apuntes_completos.pdf\#nameddest=teorema.20.26}{Ejemplo 20.26})}
\begin{pasos}
\item[] Sea $E=\{e^{(n)}:n\in\N\}\subseteq\ell^\infty$.
\item[(1)] $E$ es acotado: dos puntos distintos de $E$ están a distancia $1$, así que $\diam E=1$.
\item[(2)] $E$ es cerrado: una sucesión de puntos de $E$ que converge es de Cauchy; por (1), a partir de un índice es constante, y su límite está en $E$.
\item[(3)] $E$ no es compacto: las bolas $B(e^{(n)},\frac12)$ cubren $E$ y cada una contiene un solo punto de $E$, así que no hay subcubrimiento finito.
\end{pasos}
\bloque{\hypertarget{x-17}{}Ejercicio 17}
En el Problema \hyperlink{prob-12}{12(b)}, ¿por qué se toma un radio $R$ con un factor $\frac1{1-r}$?
\par{\small \hyperlink{s-17}{Ver solución} \quad·\quad \hyperlink{r-17}{volver a la respuesta de la pregunta 17}}
\needspace{6\baselineskip}
\subsection*{\hypertarget{e-18}{}Pregunta 18 \normalfont\small(\textsc{Punto fijo de Banach} · Clase \href{https://decostruttivismo.github.io/MAT2514-repaso/pdf/MAT2514_apuntes_completos.pdf\#nameddest=section.31}{31})}
\bloque{Ejemplo}
\textit{Las tres hipótesis de Banach (\href{https://decostruttivismo.github.io/MAT2514-repaso/pdf/MAT2514_apuntes_completos.pdf\#nameddest=teorema.31.6}{Observación 31.6})} {\small(en los apuntes: \href{https://decostruttivismo.github.io/MAT2514-repaso/pdf/MAT2514_apuntes_completos.pdf\#nameddest=teorema.31.6}{Observación 31.6})}
\begin{pasos}
\item[] Sin completitud: $T(x)=\frac{x+1}2$ en $(0,1)$ es una contracción de razón $\frac12$, y su único candidato a punto fijo, $1$, no está en $(0,1)$.
\item[] Con razón $k=1$: $T(x)=x+1$ en $\R$ cumple $\abs{T(x)-T(y)}=\abs{x-y}$ y no tiene punto fijo.
\item[] Sin razón uniforme: $T(x)=\sqrt{x^2+1}$ en $[0,\infty)$, como en la justificación.
\end{pasos}
\bloque{\hypertarget{x-18}{}Ejercicio 18}
Si además $X$ es compacto, ¿$T$ tiene un punto fijo?
\par{\small \hyperlink{s-18}{Ver solución} \quad·\quad \hyperlink{r-18}{volver a la respuesta de la pregunta 18}}
\needspace{6\baselineskip}
\subsection*{\hypertarget{e-19}{}Pregunta 19 \normalfont\small(\textsc{Funcionales lineales} · Clases \href{https://decostruttivismo.github.io/MAT2514-repaso/pdf/MAT2514_apuntes_completos.pdf\#nameddest=section.32}{32}, \href{https://decostruttivismo.github.io/MAT2514-repaso/pdf/MAT2514_apuntes_completos.pdf\#nameddest=section.34}{34})}
\bloque{Ejemplo}
\textit{$f\mapsto f(1)-f(0)$ (\href{https://decostruttivismo.github.io/MAT2514-repaso/pdf/MAT2514_apuntes_completos.pdf\#nameddest=teorema.34.5}{Ejemplo 34.5})} {\small(en los apuntes: \href{https://decostruttivismo.github.io/MAT2514-repaso/pdf/MAT2514_apuntes_completos.pdf\#nameddest=teorema.34.5}{Ejemplo 34.5})}
\begin{pasos}
\item[] Sea $\varphi(f)=f(1)-f(0)$ en $C([0,1],\R)$ con $\norm\cdot_\infty$.
\item[(1)] Cota: $\abs{\varphi(f)}\le\abs{f(1)}+\abs{f(0)}\le2\norm f_\infty$.
\item[(2)] La función $g(x)=2x-1$ cumple $-1\le g\le1$ en $[0,1]$, así que $\norm g_\infty=1$.
\item[(3)] $\varphi(g)=1-(-1)=2$.
\item[] Por (1), (2) y (3): $\norm\varphi=2$.
\item[] Con la constante $\mathbf 1$ se obtendría $\varphi(\mathbf 1)=0$: elegir bien la función de prueba es todo el cálculo.
\end{pasos}
\bloque{\hypertarget{x-19}{}Ejercicio 19}
(\href{https://decostruttivismo.github.io/MAT2514-repaso/pdf/MAT2514_apuntes_completos.pdf\#nameddest=teorema.32.10}{Observación 32.10}.) Calcule la norma de $\Lambda f=f(0)-2f(\frac12)+f(1)$ en $C([0,1],\R)$.
\par{\small \hyperlink{s-19}{Ver solución} \quad·\quad \hyperlink{r-19}{volver a la respuesta de la pregunta 19}}
\needspace{6\baselineskip}
\subsection*{\hypertarget{e-20}{}Pregunta 20 \normalfont\small(\textsc{Hahn–Banach} · Clase \href{https://decostruttivismo.github.io/MAT2514-repaso/pdf/MAT2514_apuntes_completos.pdf\#nameddest=section.35}{35})}
\bloque{Ejemplo}
\textit{Una distancia y su funcional (\href{https://decostruttivismo.github.io/MAT2514-repaso/pdf/MAT2514_apuntes_completos.pdf\#nameddest=teorema.35.8}{Ejemplo 35.8})} {\small(en los apuntes: \href{https://decostruttivismo.github.io/MAT2514-repaso/pdf/MAT2514_apuntes_completos.pdf\#nameddest=teorema.35.8}{Ejemplo 35.8})}
\begin{pasos}
\item[] En $(\R^2,\norm\cdot_\infty)$ sean $M=\{(t,t):t\in\R\}$ y $x_0=(1,0)$.
\item[(1)] $\dist(x_0,M)=\inf_t\max\{\abs{1-t},\abs t\}=\frac12$, alcanzado en $t=\frac12$.
\item[(2)] El funcional $F(y)=\frac12(y_1-y_2)$ anula $M$, porque $F(t,t)=0$.
\item[(3)] $\abs{F(y)}\le\frac12(\abs{y_1}+\abs{y_2})\le\norm y_\infty$, con igualdad en $y=(1,-1)$. Luego $\norm F=1$.
\item[] Por (1), (2) y (3): $F(x_0)=\frac12=\dist(x_0,M)$, como predice el \href{https://decostruttivismo.github.io/MAT2514-repaso/pdf/MAT2514_apuntes_completos.pdf\#nameddest=teorema.35.7}{Corolario 35.7}.
\end{pasos}
\bloque{\hypertarget{x-20}{}Ejercicio 20}
Calcule $\dist(f,M)$ para $f(x)=x$ en $C([0,1],\R)$, y para $f(x)=x^2$ en $C([-1,1],\R)$.
\par{\small \hyperlink{s-20}{Ver solución} \quad·\quad \hyperlink{r-20}{volver a la respuesta de la pregunta 20}}
\needspace{6\baselineskip}
\subsection*{\hypertarget{e-21}{}Pregunta 21 \normalfont\small(\textsc{Hahn–Banach} · Clase \href{https://decostruttivismo.github.io/MAT2514-repaso/pdf/MAT2514_apuntes_completos.pdf\#nameddest=section.35}{35})}
\bloque{Ejemplo}
\textit{Un subespacio denso}
\begin{pasos}
\item[] En $C([0,1],\R)$, sea $\mathcal P$ el subespacio de los polinomios.
\item[(1)] $\mathcal P$ es denso: $\cl{\mathcal P}=C([0,1],\R)$ (teorema de aproximación de Weierstrass, Clase 30).
\item[(2)] Sea $F\in X^*$ que anula $\mathcal P$. Por continuidad, $F$ anula $\cl{\mathcal P}$.
\item[] Por (1) y (2), $F=0$: ningún funcional continuo no nulo anula $\mathcal P$.
\end{pasos}
\bloque{\hypertarget{x-21}{}Ejercicio 21}
Pruebe que si $M$ es un subespacio de un espacio normado, entonces $\cl M$ también lo es.
\par{\small \hyperlink{s-21}{Ver solución} \quad·\quad \hyperlink{r-21}{volver a la respuesta de la pregunta 21}}
\needspace{6\baselineskip}
\subsection*{\hypertarget{e-22}{}Pregunta 22 \normalfont\small(\textsc{Funcionales lineales} · Clase \href{https://decostruttivismo.github.io/MAT2514-repaso/pdf/MAT2514_apuntes_completos.pdf\#nameddest=section.34}{34})}
\bloque{Ejemplo}
\textit{La evaluación fuera del intervalo (\href{https://decostruttivismo.github.io/MAT2514-repaso/pdf/MAT2514_apuntes_completos.pdf\#nameddest=teorema.34.6}{Ejemplo 34.6})} {\small(en los apuntes: \href{https://decostruttivismo.github.io/MAT2514-repaso/pdf/MAT2514_apuntes_completos.pdf\#nameddest=teorema.34.6}{Ejemplo 34.6})}
\begin{pasos}
\item[] En el espacio $\mathcal P$ de los polinomios con la norma $\norm p_\infty=\max_{[0,1]}\abs p$, sea $\psi(p)=p(2)$.
\item[(1)] $p_n(x)=x^n$ cumple $\norm{p_n}_\infty=1$ y $\psi(p_n)=2^n$. Luego $\psi$ no es acotado.
\item[(2)] $q_n(x)=(\frac x2)^n$ cumple $\norm{q_n}_\infty=2^{-n}\to0$.
\item[(3)] Pero $\psi(q_n)=1$ para todo $n$, que no tiende a $\psi(0)=0$.
\item[] Por (2) y (3), $\psi$ no es continuo en $0$.
\end{pasos}
\bloque{\hypertarget{x-22}{}Ejercicio 22}
¿Puede pasar lo mismo en $\R^n$ con $\norm\cdot_\infty$?
\par{\small \hyperlink{s-22}{Ver solución} \quad·\quad \hyperlink{r-22}{volver a la respuesta de la pregunta 22}}
\needspace{6\baselineskip}
\subsection*{\hypertarget{e-23}{}Pregunta 23 \normalfont\small(\textsc{Funcionales lineales} · Clase \href{https://decostruttivismo.github.io/MAT2514-repaso/pdf/MAT2514_apuntes_completos.pdf\#nameddest=section.34}{34})}
\bloque{Ejemplo}
\textit{La evaluación fuera del intervalo, en cualquier punto (\href{https://decostruttivismo.github.io/MAT2514-repaso/pdf/MAT2514_apuntes_completos.pdf\#nameddest=teorema.34.6}{Ejemplo 34.6})} {\small(en los apuntes: \href{https://decostruttivismo.github.io/MAT2514-repaso/pdf/MAT2514_apuntes_completos.pdf\#nameddest=teorema.34.6}{Ejemplo 34.6})}
\begin{pasos}
\item[] Sea $\psi(p)=p(2)$ en los polinomios con $\norm\cdot_\infty$ sobre $[0,1]$, y $q_n(x)=(\frac x2)^n$.
\item[(1)] En $0$: $q_n\to0$, pero $\psi(q_n)=1\not\to0$.
\item[(2)] Sea $p_0$ un polinomio cualquiera. Entonces $p_0+q_n\to p_0$.
\item[(3)] $\psi(p_0+q_n)=\psi(p_0)+1$, que no tiende a $\psi(p_0)$.
\item[] Por (2) y (3), $\psi$ no es continuo en $p_0$.
\end{pasos}
\bloque{\hypertarget{x-23}{}Ejercicio 23}
Pruebe que un funcional lineal acotado es Lipschitz.
\par{\small \hyperlink{s-23}{Ver solución} \quad·\quad \hyperlink{r-23}{volver a la respuesta de la pregunta 23}}
\needspace{6\baselineskip}
\subsection*{\hypertarget{e-24}{}Pregunta 24 \normalfont\small(\textsc{Funcionales lineales} · Clase \href{https://decostruttivismo.github.io/MAT2514-repaso/pdf/MAT2514_apuntes_completos.pdf\#nameddest=section.34}{34})}
\bloque{Ejemplo}
\textit{Un normado no completo (\href{https://decostruttivismo.github.io/MAT2514-repaso/pdf/MAT2514_apuntes_completos.pdf\#nameddest=teorema.32.3}{Ejemplo 32.3})} {\small(en los apuntes: \href{https://decostruttivismo.github.io/MAT2514-repaso/pdf/MAT2514_apuntes_completos.pdf\#nameddest=teorema.32.3}{Ejemplo 32.3})}
\begin{pasos}
\item[] En $c_{00}$ sea $x^{(m)}=(1,\frac12,\dots,\frac1m,0,0,\dots)$.
\item[(1)] Si $m\gt n$, entonces $\norm{x^{(m)}-x^{(n)}}_\infty=\frac1{n+1}$. Luego $(x^{(m)})$ es de Cauchy.
\item[(2)] Su único candidato a límite es $(\frac1n)_n$, que no es finalmente nula, así que no está en $c_{00}$.
\item[] Por (1) y (2), $c_{00}$ no es completo. Aun así, $c_{00}^*$ es de Banach.
\end{pasos}
\bloque{\hypertarget{x-24}{}Ejercicio 24}
¿En qué paso exacto de la demostración se usa la completitud de $\R$?
\par{\small \hyperlink{s-24}{Ver solución} \quad·\quad \hyperlink{r-24}{volver a la respuesta de la pregunta 24}}
\needspace{6\baselineskip}
\subsection*{\hypertarget{e-25}{}Pregunta 25 \normalfont\small(\textsc{Gromov–Hausdorff} · Clase \href{https://decostruttivismo.github.io/MAT2514-repaso/pdf/MAT2514_apuntes_completos.pdf\#nameddest=section.40}{40})}
\bloque{Ejemplo}
\textit{Un punto con dos compañeros (\href{https://decostruttivismo.github.io/MAT2514-repaso/pdf/MAT2514_apuntes_completos.pdf\#nameddest=teorema.40.2}{Ejemplo 40.2})} {\small(en los apuntes: \href{https://decostruttivismo.github.io/MAT2514-repaso/pdf/MAT2514_apuntes_completos.pdf\#nameddest=teorema.40.2}{Ejemplo 40.2})}
\begin{pasos}
\item[] Sean $X=\{0,1,2\}$ e $Y=\{0,2\}$, con la distancia de $\R$, y $R=\{(0,0),(1,0),(1,2),(2,2)\}$.
\item[(1)] $R$ es una correspondencia: todo punto de $X$ aparece como primera coordenada y todo punto de $Y$ como segunda.
\item[(2)] El punto $1\in X$ tiene dos compañeros en $Y$: $0$ y $2$. El par $(1,0),(1,2)$ aporta $2$.
\item[(3)] Los demás pares aportan a lo más $1$.
\item[] Por (2) y (3), $\dis R=2$, y el valor lo da justamente el par que comparte una coordenada.
\end{pasos}
\bloque{\hypertarget{x-25}{}Ejercicio 25}
Con $X=\{0,1,2\}$ e $Y=\{0,2\}$ (distancia de $\R$), calcule $\dis R$ para $R=\{(0,0),(1,0),(2,2)\}$.
\par{\small \hyperlink{s-25}{Ver solución} \quad·\quad \hyperlink{r-25}{volver a la respuesta de la pregunta 25}}
\needspace{6\baselineskip}
\subsection*{\hypertarget{e-26}{}Pregunta 26 \normalfont\small(\textsc{Gromov–Hausdorff} · Clases \href{https://decostruttivismo.github.io/MAT2514-repaso/pdf/MAT2514_apuntes_completos.pdf\#nameddest=section.40}{40}, \href{https://decostruttivismo.github.io/MAT2514-repaso/pdf/MAT2514_apuntes_completos.pdf\#nameddest=section.39}{39})}
\bloque{Ejemplo}
\textit{El triángulo y dos vértices (\href{https://decostruttivismo.github.io/MAT2514-repaso/pdf/MAT2514_apuntes_completos.pdf\#nameddest=teorema.40.7}{Ejemplo 40.7})} {\small(en los apuntes: \href{https://decostruttivismo.github.io/MAT2514-repaso/pdf/MAT2514_apuntes_completos.pdf\#nameddest=teorema.40.7}{Ejemplo 40.7})}
\begin{pasos}
\item[] Sean $T=\{u,v,w\}$ con las tres distancias iguales a $1$ e $Y=\{u,v\}$ con $d(u,v)=1$.
\item[] Sea $R$ una correspondencia cualquiera entre $T$ e $Y$.
\item[(1)] Cada uno de los tres puntos de $T$ tiene al menos un compañero en $Y$, que tiene solo dos puntos.
\item[(2)] Por (1) y el principio del palomar, hay dos puntos distintos $p\ne p'$ de $T$ con un mismo compañero $y\in Y$.
\item[(3)] Los elementos $(p,y),(p',y)\in R$ aportan $\abs{d(p,p')-d(y,y)}=\abs{1-0}=1$. Luego $\dis R\ge1$.
\item[(4)] (3) vale para toda correspondencia $R$. Por el \href{https://decostruttivismo.github.io/MAT2514-repaso/pdf/MAT2514_apuntes_completos.pdf\#nameddest=teorema.40.4}{Teorema 40.4}, $d_{GH}(T,Y)\ge\frac12$.
\item[] Con la cota superior $d_{GH}(T,Y)\le\frac12$ (\href{https://decostruttivismo.github.io/MAT2514-repaso/pdf/MAT2514_apuntes_completos.pdf\#nameddest=teorema.39.10}{Ejemplo 39.10}), resulta $d_{GH}(T,Y)=\frac12$.
\end{pasos}
\bloque{\hypertarget{x-26}{}Ejercicio 26}
Sea $X=Y=\{0,1\}$ con la distancia de $\R$. Muestre que hay una correspondencia de distorsión $1$, y que aun así $d_{GH}(X,Y)=0$.
\par{\small \hyperlink{s-26}{Ver solución} \quad·\quad \hyperlink{r-26}{volver a la respuesta de la pregunta 26}}
\needspace{6\baselineskip}
\subsection*{\hypertarget{e-27}{}Pregunta 27 \normalfont\small(\textsc{Gromov–Hausdorff} · Clases \href{https://decostruttivismo.github.io/MAT2514-repaso/pdf/MAT2514_apuntes_completos.pdf\#nameddest=section.39}{39}, \href{https://decostruttivismo.github.io/MAT2514-repaso/pdf/MAT2514_apuntes_completos.pdf\#nameddest=section.40}{40})}
\bloque{Ejemplo}
\textit{Mismo diámetro, espacios distintos (\href{https://decostruttivismo.github.io/MAT2514-repaso/pdf/MAT2514_apuntes_completos.pdf\#nameddest=teorema.39.11}{Observación 39.11})} {\small(en los apuntes: \href{https://decostruttivismo.github.io/MAT2514-repaso/pdf/MAT2514_apuntes_completos.pdf\#nameddest=teorema.39.11}{Observación 39.11})}
\begin{pasos}
\item[] Sean $X=[0,1]$ e $Y=\{0,1\}$, con la distancia de $\R$.
\item[(1)] $\diam X=1$ y $\diam Y=1$. La cota por diámetros da $d_{GH}(X,Y)\ge0$.
\item[(2)] $X$ es infinito e $Y$ tiene dos puntos, así que no son isométricos.
\item[(3)] Por (2) y la rigidez (\href{https://decostruttivismo.github.io/MAT2514-repaso/pdf/MAT2514_apuntes_completos.pdf\#nameddest=teorema.39.12}{Teorema 39.12}), $d_{GH}(X,Y)\gt 0$.
\item[] Por (1) y (3), los diámetros no detectan una distancia que sí es positiva.
\end{pasos}
\bloque{\hypertarget{x-27}{}Ejercicio 27}
(\href{https://decostruttivismo.github.io/MAT2514-repaso/pdf/MAT2514_apuntes_completos.pdf\#nameddest=teorema.39.7}{Ejemplo 39.7}.) Calcule $d_{GH}([0,1],[0,3])$.
\par{\small \hyperlink{s-27}{Ver solución} \quad·\quad \hyperlink{r-27}{volver a la respuesta de la pregunta 27}}
\needspace{6\baselineskip}
\subsection*{\hypertarget{e-28}{}Pregunta 28 \normalfont\small(\textsc{Gromov–Hausdorff} · Clase \href{https://decostruttivismo.github.io/MAT2514-repaso/pdf/MAT2514_apuntes_completos.pdf\#nameddest=section.39}{39})}
\bloque{Ejemplo}
\textit{La desigualdad puede ser estricta (\href{https://decostruttivismo.github.io/MAT2514-repaso/pdf/MAT2514_apuntes_completos.pdf\#nameddest=teorema.39.5}{Ejemplo 39.5})} {\small(en los apuntes: \href{https://decostruttivismo.github.io/MAT2514-repaso/pdf/MAT2514_apuntes_completos.pdf\#nameddest=teorema.39.5}{Ejemplo 39.5})}
\begin{pasos}
\item[] Sean $A=\{0,1,3\}$ y $B=\{0,2,3\}$ en $\R$.
\item[(1)] $d_H(A,B)=1$: el punto $1\in A$ está a distancia $1$ de $B$, el punto $2\in B$ está a distancia $1$ de $A$, y los demás puntos están en ambos conjuntos.
\item[(2)] La aplicación $t\mapsto3-t$ manda $A$ sobre $B$ ($0\mapsto3$, $1\mapsto2$, $3\mapsto0$) y preserva distancias. Luego $A$ y $B$ son isométricos y $d_{GH}(A,B)=0$.
\item[] Por (1) y (2): $d_{GH}(A,B)=0\lt 1=d_H(A,B)$.
\end{pasos}
\bloque{\hypertarget{x-28}{}Ejercicio 28}
Pruebe que si $S\subseteq X$ es una $\varepsilon$-red de $X$ (con $X$ compacto), entonces $d_{GH}(S,X)\le\varepsilon$.
\par{\small \hyperlink{s-28}{Ver solución} \quad·\quad \hyperlink{r-28}{volver a la respuesta de la pregunta 28}}
\needspace{6\baselineskip}
\subsection*{\hypertarget{e-29}{}Pregunta 29 \normalfont\small(\textsc{Gromov–Hausdorff} · Clase \href{https://decostruttivismo.github.io/MAT2514-repaso/pdf/MAT2514_apuntes_completos.pdf\#nameddest=section.41}{41})}
\bloque{Ejemplo}
\textit{Dos convergencias (\href{https://decostruttivismo.github.io/MAT2514-repaso/pdf/MAT2514_apuntes_completos.pdf\#nameddest=teorema.41.7}{Ejemplo 41.7})} {\small(en los apuntes: \href{https://decostruttivismo.github.io/MAT2514-repaso/pdf/MAT2514_apuntes_completos.pdf\#nameddest=teorema.41.7}{Ejemplo 41.7})}
\begin{pasos}
\item[] (a) La grilla $X_n=\{0,\frac1n,\dots,1\}$ cumple $d_{GH}(X_n,[0,1])\le\frac1{2n}\to0$: espacios finitos convergen a un segmento.
\item[] (b) Los rectángulos $Q_n=[0,1]\times[0,\frac1n]\subseteq\R^2$, con la distancia euclidiana, convergen a $[0,1]$.
\item[] En (b), la dimensión baja de $2$ a $1$ en el límite.
\end{pasos}
\bloque{\hypertarget{x-29}{}Ejercicio 29}
Pruebe que $X_n=\{0,\frac1n,\dots,1\}$ es disconexo para $n\ge1$.
\par{\small \hyperlink{s-29}{Ver solución} \quad·\quad \hyperlink{r-29}{volver a la respuesta de la pregunta 29}}
\needspace{6\baselineskip}
\subsection*{\hypertarget{e-30}{}Pregunta 30 \normalfont\small(\textsc{Gromov–Hausdorff} · Clase \href{https://decostruttivismo.github.io/MAT2514-repaso/pdf/MAT2514_apuntes_completos.pdf\#nameddest=section.41}{41})}
\bloque{Ejemplo}
\textit{Redondear a la grilla (\href{https://decostruttivismo.github.io/MAT2514-repaso/pdf/MAT2514_apuntes_completos.pdf\#nameddest=teorema.41.2}{Ejemplo 41.2})} {\small(en los apuntes: \href{https://decostruttivismo.github.io/MAT2514-repaso/pdf/MAT2514_apuntes_completos.pdf\#nameddest=teorema.41.2}{Ejemplo 41.2})}
\begin{pasos}
\item[] Sean $X_n=\{0,\frac1n,\dots,1\}$ y $f\colon[0,1]\to X_n$, $f(x)=\lfloor nx\rfloor/n$.
\item[(1)] Para todo $x$, $0\le x-f(x)\lt \frac1n$. De ahí sale $\dis f\le\frac1n$ (ver el ejercicio de la pregunta 31).
\item[(2)] $f$ es sobreyectiva: $f(\frac kn)=\frac kn$. Luego $f([0,1])=X_n$, que es una $\frac1n$-red de $X_n$.
\item[] Por (1) y (2), $f$ es una $\frac1n$-isometría.
\end{pasos}
\bloque{\hypertarget{x-30}{}Ejercicio 30}
(Problema \hyperlink{prob-17}{17(b)}.) Con $f_n(i,a)=\frac in$ en la escalera $G_n$, ¿es $f_n(G_n)$ una $\frac1n$-red de $[0,1]$?
\par{\small \hyperlink{s-30}{Ver solución} \quad·\quad \hyperlink{r-30}{volver a la respuesta de la pregunta 30}}
\needspace{6\baselineskip}
\subsection*{\hypertarget{e-31}{}Pregunta 31 \normalfont\small(\textsc{Gromov–Hausdorff} · Clase \href{https://decostruttivismo.github.io/MAT2514-repaso/pdf/MAT2514_apuntes_completos.pdf\#nameddest=section.41}{41})}
\bloque{Ejemplo}
\textit{Redondear a la grilla (\href{https://decostruttivismo.github.io/MAT2514-repaso/pdf/MAT2514_apuntes_completos.pdf\#nameddest=teorema.41.2}{Ejemplo 41.2})} {\small(en los apuntes: \href{https://decostruttivismo.github.io/MAT2514-repaso/pdf/MAT2514_apuntes_completos.pdf\#nameddest=teorema.41.2}{Ejemplo 41.2})}
\begin{pasos}
\item[(1)] $f(x)=\lfloor nx\rfloor/n$ salta en cada uno de los puntos $\frac1n,\frac2n,\dots,1$.
\item[(2)] Aun así $\dis f\le\frac1n$ y $f$ es sobreyectiva sobre $X_n$.
\item[] Por (2), $f$ es una $\frac1n$-isometría, aunque por (1) no es continua.
\end{pasos}
\bloque{\hypertarget{x-31}{}Ejercicio 31}
Para el redondeo $f(x)=\lfloor nx\rfloor/n$, pruebe que $\dis f\le\frac1n$.
\par{\small \hyperlink{s-31}{Ver solución} \quad·\quad \hyperlink{r-31}{volver a la respuesta de la pregunta 31}}
\needspace{6\baselineskip}
\subsection*{\hypertarget{e-32}{}Pregunta 32 \normalfont\small(\textsc{Gromov–Hausdorff} · Clases \href{https://decostruttivismo.github.io/MAT2514-repaso/pdf/MAT2514_apuntes_completos.pdf\#nameddest=section.41}{41}, \href{https://decostruttivismo.github.io/MAT2514-repaso/pdf/MAT2514_apuntes_completos.pdf\#nameddest=section.19}{19})}
\bloque{Ejemplo}
\textit{El mismo fenómeno en $\R$}
\begin{pasos}
\item[(1)] $\{\frac1n:n\ge1\}$ no es compacto: la sucesión $\frac1n$ converge a $0$, que no está en el conjunto, y ninguna subsucesión converge dentro del conjunto.
\item[(2)] $\{\frac1n:n\ge1\}\cup\{0\}$ sí es compacto: es cerrado y acotado en $\R$.
\item[] Por (1) y (2): lo que hace falta es agregar el límite.
\end{pasos}
\bloque{\hypertarget{x-32}{}Ejercicio 32}
Pruebe que en un espacio métrico, si $x_n\to x$, entonces $C=\{x_n:n\in\N\}\cup\{x\}$ es compacto.
\par{\small \hyperlink{s-32}{Ver solución} \quad·\quad \hyperlink{r-32}{volver a la respuesta de la pregunta 32}}
\needspace{6\baselineskip}
\subsection*{\hypertarget{e-33}{}Pregunta 33 \normalfont\small(\textsc{Gromov–Hausdorff} · Clases \href{https://decostruttivismo.github.io/MAT2514-repaso/pdf/MAT2514_apuntes_completos.pdf\#nameddest=section.41}{41}, \href{https://decostruttivismo.github.io/MAT2514-repaso/pdf/MAT2514_apuntes_completos.pdf\#nameddest=section.40}{40})}
\bloque{Ejemplo}
\textit{Los rectángulos delgados (\href{https://decostruttivismo.github.io/MAT2514-repaso/pdf/MAT2514_apuntes_completos.pdf\#nameddest=teorema.41.10}{Ejemplo 41.10})} {\small(en los apuntes: \href{https://decostruttivismo.github.io/MAT2514-repaso/pdf/MAT2514_apuntes_completos.pdf\#nameddest=teorema.41.10}{Ejemplo 41.10})}
\begin{pasos}
\item[] Sea $Q_n=[0,1]\times[0,\frac1n]\subseteq\R^2$, con la distancia euclidiana.
\item[(1)] Los puntos más lejanos de $Q_n$ son esquinas opuestas. Luego $\diam Q_n=\sqrt{1+n^{-2}}$.
\item[(2)] $\diam[0,1]=1$.
\item[(3)] Por (1), (2) y la proposición: $d_{GH}(Q_n,[0,1])\ge\frac12\bigl(\sqrt{1+n^{-2}}-1\bigr)\gt 0$.
\item[] Además $\diam Q_n\to1=\diam[0,1]$, como debe pasar si $Q_n\xrightarrow{GH}[0,1]$.
\end{pasos}
\bloque{\hypertarget{x-33}{}Ejercicio 33}
Para la escalera $G_n$ del Problema \hyperlink{prob-17}{17}, con $d_n((i,a),(j,b))=\frac{\abs{i-j}+\abs{a-b}}n$, calcule $\diam G_n$ y verifique que tiende a $1$.
\par{\small \hyperlink{s-33}{Ver solución} \quad·\quad \hyperlink{r-33}{volver a la respuesta de la pregunta 33}}
\newpage\parte{Parte 4. Soluciones de los ejercicios}
\needspace{6\baselineskip}
\subsection*{\hypertarget{s-1}{}Solución del ejercicio 1}
\emph{(\href{https://decostruttivismo.github.io/MAT2514-repaso/pdf/MAT2514_apuntes_completos.pdf\#nameddest=ejercicio.A12B.6}{Ejercicio A12B.6}.) Sea $K\subseteq X$ compacto no vacío y $x\in X$. Demuestre que existe $y_0\in K$ con $d(x,y_0)=d(x,K)$.}
\begin{pasos}
\item[] Llamamos $m=d(x,K)=\inf_{y\in K}d(x,y)$.
\item[(1)] Para cada $n\ge1$, el número $m+\frac1n$ no es cota inferior de $\{d(x,y):y\in K\}$.
\item[(2)] Por (1), existe $y_n\in K$ con $m\le d(x,y_n)\lt m+\frac1n$. Así $d(x,y_n)\to m$.
\item[(3)] $K$ es compacto, luego secuencialmente compacto: existe una subsucesión $(y_{n_k})_k$ que converge a algún $y_0\in K$.
\item[(4)] Por la triangular inversa, $\abs{d(x,y_{n_k})-d(x,y_0)}\le d(y_{n_k},y_0)$.
\item[(5)] Por (3) y (4), $d(x,y_{n_k})\to d(x,y_0)$.
\item[(6)] Por (2), la subsucesión $(d(x,y_{n_k}))_k$ también converge a $m$.
\item[] Por (5), (6) y la unicidad del límite en $\R$: $d(x,y_0)=m=d(x,K)$.
\end{pasos}
\par{\small \hyperlink{x-1}{volver al ejercicio} \quad·\quad \hyperlink{r-1}{respuesta de la pregunta 1}}
\needspace{6\baselineskip}
\subsection*{\hypertarget{s-2}{}Solución del ejercicio 2}
\emph{Demuestre: si $K$ es compacto, $F$ es cerrado, ambos son no vacíos y $K\cap F=\varnothing$, entonces $d(K,F)\gt 0$.}
\begin{pasos}
\item[] Supongamos, por contradicción, que $d(K,F)=0$.
\item[(1)] Para cada $n\ge1$ existen $y_n\in K$ y $z_n\in F$ con $d(y_n,z_n)\lt \frac1n$.
\item[(2)] $K$ es secuencialmente compacto: alguna subsucesión $(y_{n_k})_k$ converge a un punto $y\in K$.
\item[(3)] Por la triangular y (1): $d(y,z_{n_k})\le d(y,y_{n_k})+d(y_{n_k},z_{n_k})\lt d(y,y_{n_k})+\frac1{n_k}$.
\item[(4)] Por (2), el lado derecho de (3) tiende a $0$. Luego $z_{n_k}\to y$.
\item[(5)] Los $z_{n_k}$ están en $F$, que es cerrado. Por (4), el límite $y$ está en $F$.
\item[(6)] Por (2) y (5), $y\in K\cap F$.
\item[] (6) contradice $K\cap F=\varnothing$. Luego $d(K,F)\gt 0$.
\end{pasos}
\par{\small \hyperlink{x-2}{volver al ejercicio} \quad·\quad \hyperlink{r-2}{respuesta de la pregunta 2}}
\needspace{6\baselineskip}
\subsection*{\hypertarget{s-3}{}Solución del ejercicio 3}
\emph{Sea $K=[0,1]\cup[2,3]$ y sea $L\subseteq\R$ compacto con $d_H(L,K)\lt \frac13$. Pruebe que $L$ es disconexo.}
\begin{pasos}
\item[] Sean $K_1=[0,1]$ y $K_2=[2,3]$, de modo que $d(K_1,K_2)=1$.
\item[] Definimos $L_1=\{y\in L:\ d(y,K_1)\lt \frac13\}$ y $L_2=\{y\in L:\ d(y,K_2)\lt \frac13\}$.
\item[(1)] $L_1\cup L_2=L$.
\begin{subpasos}
\item Sea $y\in L$. Como $d_H(L,K)\lt \frac13$, se tiene $d(y,K)\lt \frac13$.
\item Además $d(y,K)=\min\{d(y,K_1),\,d(y,K_2)\}$. Luego $d(y,K_1)\lt \frac13$ o $d(y,K_2)\lt \frac13$.
\end{subpasos}
\item[(2)] $L_1\cap L_2=\varnothing$.
\begin{subpasos}
\item Si $y\in L_1\cap L_2$, hay $a\in K_1$ con $d(y,a)\lt \frac13$ y $b\in K_2$ con $d(y,b)\lt \frac13$.
\item Entonces $d(a,b)\le d(a,y)+d(y,b)\lt \frac23$.
\item Pero $d(a,b)\ge d(K_1,K_2)=1$. Contradicción.
\end{subpasos}
\item[(3)] $L_1$ y $L_2$ son abiertos en $L$: la función $y\mapsto d(y,K_i)$ es continua, y $L_i$ es la preimagen del abierto $(-\infty,\frac13)$ por su restricción a $L$.
\item[(4)] $L_1\ne\varnothing$.
\begin{subpasos}
\item Sea $a=0\in K$. Como $d_H(L,K)\lt \frac13$, se tiene $d(a,L)\lt \frac13$: existe $y\in L$ con $d(y,a)\lt \frac13$.
\item Como $a\in K_1$, ese $y$ cumple $d(y,K_1)\lt \frac13$, es decir, $y\in L_1$.
\end{subpasos}
\item[(5)] $L_2\ne\varnothing$: igual que en (4), partiendo de $a=2\in K_2$.
\item[] Por (1)–(5), $(L_1,L_2)$ es una disconexión de $L$.
\end{pasos}
\par{\small \hyperlink{x-3}{volver al ejercicio} \quad·\quad \hyperlink{r-3}{respuesta de la pregunta 3}}
\needspace{6\baselineskip}
\subsection*{\hypertarget{s-4}{}Solución del ejercicio 4}
\emph{Sea $X$ completo y $A\subseteq X$ totalmente acotado. Pruebe que $\cl A$ es compacto.}
\begin{pasos}
\item[(1)] $\cl A$ es completo: es cerrado dentro del completo $X$ (\href{https://decostruttivismo.github.io/MAT2514-repaso/pdf/MAT2514_apuntes_completos.pdf\#nameddest=teorema.15.12}{Proposición 15.12}).
\item[(2)] $\cl A$ es totalmente acotado.
\begin{subpasos}
\item Sea $\varepsilon\gt 0$. Como $A$ es totalmente acotado, tiene una $\frac\varepsilon2$-red finita $F$.
\item Sea $z\in\cl A$. Existe $a\in A$ con $d(z,a)\lt \frac\varepsilon2$.
\item Como $F$ es una $\frac\varepsilon2$-red de $A$, existe $y\in F$ con $d(a,y)\lt \frac\varepsilon2$.
\item Por la triangular, $d(z,y)\lt \varepsilon$. Luego $F$ es una $\varepsilon$-red finita de $\cl A$.
\end{subpasos}
\item[] Por (1), (2) y el \href{https://decostruttivismo.github.io/MAT2514-repaso/pdf/MAT2514_apuntes_completos.pdf\#nameddest=teorema.19.6}{Corolario 19.6}, $\cl A$ es compacto.
\end{pasos}
\par{\small \hyperlink{x-4}{volver al ejercicio} \quad·\quad \hyperlink{r-4}{respuesta de la pregunta 4}}
\needspace{6\baselineskip}
\subsection*{\hypertarget{s-5}{}Solución del ejercicio 5}
\emph{Encuentre una $\frac14$-red finita de la circunferencia unitaria $S^1=\{x\in\R^2:\norm x=1\}$.}
\begin{pasos}
\item[] Sea $C_n$ el conjunto de los $n$ puntos de $S^1$ equiespaciados en ángulo, como en el \href{https://decostruttivismo.github.io/MAT2514-repaso/pdf/MAT2514_apuntes_completos.pdf\#nameddest=teorema.41.8}{Ejemplo 41.8}.
\item[(1)] Dos puntos consecutivos de $C_n$ forman un ángulo $\frac{2\pi}n$. Luego todo punto $p\in S^1$ forma un ángulo $\alpha\le\frac\pi n$ con algún punto $c\in C_n$.
\item[(2)] La distancia de $p$ a $c$ es la cuerda $\norm{p-c}=2\sin\frac\alpha2$.
\item[(3)] Como $\sin u\le u$ para $u\ge0$, por (1) y (2): $\norm{p-c}\le\alpha\le\frac\pi n$.
\item[(4)] Queremos $\frac\pi n\lt \frac14$, es decir $n\gt 4\pi\approx12{,}57$.
\item[] Por (3) y (4), con $n=13$ el conjunto $C_{13}$ es una $\frac14$-red finita de $S^1$.
\end{pasos}
\par{\small \hyperlink{x-5}{volver al ejercicio} \quad·\quad \hyperlink{r-5}{respuesta de la pregunta 5}}
\needspace{6\baselineskip}
\subsection*{\hypertarget{s-6}{}Solución del ejercicio 6}
\emph{Con la distancia de $\R$: ¿es completo $A=\{\frac1n:n\ge1\}\cup\{0\}$? ¿Y $B=\{\frac1n:n\ge1\}$?}
\begin{pasos}
\item[(1)] $A$ es cerrado en $\R$: su único punto de acumulación es $0$, y $0\in A$.
\item[(2)] Por (1) y la afirmación de esta pregunta, $A$ es completo.
\item[(3)] En $B$, la sucesión $\frac1n$ es de Cauchy, porque converge en $\R$.
\item[(4)] Su límite en $\R$ es $0$, que no está en $B$. Luego no converge en $B$.
\item[] Por (3) y (4), $B$ no es completo. (Así se resuelven las partes (c) y (d) del \href{https://decostruttivismo.github.io/MAT2514-repaso/pdf/MAT2514_apuntes_completos.pdf\#nameddest=ejercicio.A12B.4}{Ejercicio A12B.4}.)
\end{pasos}
\par{\small \hyperlink{x-6}{volver al ejercicio} \quad·\quad \hyperlink{r-6}{respuesta de la pregunta 6}}
\needspace{6\baselineskip}
\subsection*{\hypertarget{s-7}{}Solución del ejercicio 7}
\emph{Calcule $d_H(\{0,3\},[0,4])$.}
\begin{pasos}
\item[(1)] Desviación de $\{0,3\}$: los puntos $0$ y $3$ están en $[0,4]$. Vale $0$.
\item[(2)] Para $x\in[0,\frac32]$, el punto de $\{0,3\}$ más cercano es $0$, y $d(x,\{0,3\})=x\le\frac32$.
\item[(3)] Para $x\in[\frac32,3]$, el más cercano es $3$, y $d(x,\{0,3\})=3-x\le\frac32$.
\item[(4)] Para $x\in[3,4]$, el más cercano es $3$, y $d(x,\{0,3\})=x-3\le1$.
\item[(5)] Por (2), (3) y (4), la desviación de $[0,4]$ es $\frac32$, alcanzada en $x=\frac32$.
\item[] Por (1) y (5): $d_H(\{0,3\},[0,4])=\max\{0,\frac32\}=\frac32$.
\end{pasos}
\par{\small \hyperlink{x-7}{volver al ejercicio} \quad·\quad \hyperlink{r-7}{respuesta de la pregunta 7}}
\needspace{6\baselineskip}
\subsection*{\hypertarget{s-8}{}Solución del ejercicio 8}
\emph{Calcule $d_H([0,1],\{0\})$.}
\begin{pasos}
\item[(1)] Desviación de $\{0\}$: como $0\in[0,1]$, vale $0$.
\item[(2)] Desviación de $[0,1]$: para $x\in[0,1]$, $d(x,\{0\})=x$. Su supremo es $1$, en $x=1$.
\item[] Por (1) y (2): $d_H([0,1],\{0\})=1$, aunque $\{0\}\subseteq[0,1]$.
\end{pasos}
\par{\small \hyperlink{x-8}{volver al ejercicio} \quad·\quad \hyperlink{r-8}{respuesta de la pregunta 8}}
\needspace{6\baselineskip}
\subsection*{\hypertarget{s-9}{}Solución del ejercicio 9}
\emph{Pruebe que si $A,B$ son acotados y $d_H(A,B)=0$, entonces $\cl A=\cl B$.}
\begin{pasos}
\item[(1)] Como $d_H(A,B)=0$, la desviación de $A$ es $0$: cada $a\in A$ cumple $d(a,B)=0$.
\item[(2)] $d(a,B)=0$ equivale a $a\in\cl B$ (\href{https://decostruttivismo.github.io/MAT2514-repaso/pdf/MAT2514_apuntes_completos.pdf\#nameddest=teorema.20.14}{Proposición 20.14}). Por (1), $A\subseteq\cl B$.
\item[(3)] $\cl B$ es cerrado y contiene a $A$. Por (2), $\cl A\subseteq\cl B$.
\item[(4)] Intercambiando $A$ y $B$ en (1)–(3), $\cl B\subseteq\cl A$.
\item[] Por (3) y (4), $\cl A=\cl B$.
\end{pasos}
\par{\small \hyperlink{x-9}{volver al ejercicio} \quad·\quad \hyperlink{r-9}{respuesta de la pregunta 9}}
\needspace{6\baselineskip}
\subsection*{\hypertarget{s-10}{}Solución del ejercicio 10}
\emph{¿Qué propiedad común de las funciones $\varphi_A$ da la equicontinuidad en el Problema \hyperlink{prob-4}{4(a)}?}
\begin{pasos}
\item[(1)] Cada $\varphi_A$ es $1$-Lipschitz (\href{https://decostruttivismo.github.io/MAT2514-repaso/pdf/MAT2514_apuntes_completos.pdf\#nameddest=teorema.20.11}{Proposición 20.11}): $\abs{\varphi_A(x)-\varphi_A(y)}\le d(x,y)$ para todos $x,y$.
\item[(2)] La constante $1$ en (1) es la misma para todos los conjuntos $A$.
\item[(3)] Dado $\varepsilon\gt 0$, tomamos $\delta=\varepsilon$. Por (1), si $d(x,y)\lt \delta$ entonces $\abs{\varphi_A(x)-\varphi_A(y)}\lt \varepsilon$.
\item[] Por (2), el mismo $\delta$ de (3) sirve para todas las $\varphi_A$ a la vez: eso es la equicontinuidad (\href{https://decostruttivismo.github.io/MAT2514-repaso/pdf/MAT2514_apuntes_completos.pdf\#nameddest=teorema.29.1}{Definición 29.1}).
\end{pasos}
\par{\small \hyperlink{x-10}{volver al ejercicio} \quad·\quad \hyperlink{r-10}{respuesta de la pregunta 10}}
\needspace{6\baselineskip}
\subsection*{\hypertarget{s-11}{}Solución del ejercicio 11}
\emph{Pruebe que $d_H(A\cup B,\,A'\cup B)\le d_H(A,A')$.}
\begin{pasos}
\item[(1)] Aplicamos el \href{https://decostruttivismo.github.io/MAT2514-repaso/pdf/MAT2514_apuntes_completos.pdf\#nameddest=teorema.37.5}{Lema 37.5}(a) con $A_1=A$, $B_1=A'$, $A_2=B$ y $B_2=B$.
\item[(2)] $d_H(B,B)=0$.
\item[] Por (1) y (2): $d_H(A\cup B,A'\cup B)\le\max\{d_H(A,A'),0\}=d_H(A,A')$.
\end{pasos}
\par{\small \hyperlink{x-11}{volver al ejercicio} \quad·\quad \hyperlink{r-11}{respuesta de la pregunta 11}}
\needspace{6\baselineskip}
\subsection*{\hypertarget{s-12}{}Solución del ejercicio 12}
\emph{¿Por qué el argumento falla para $A_n=[n,\infty)\subseteq\R$?}
\begin{pasos}
\item[(1)] Los $A_n$ son cerrados pero no compactos, porque no son acotados.
\item[(2)] $\bigcap_nA_n=\varnothing$: ningún número real es mayor o igual que todo $n$.
\item[] Por (1) y (2): el \href{https://decostruttivismo.github.io/MAT2514-repaso/pdf/MAT2514_apuntes_completos.pdf\#nameddest=teorema.19.13}{Corolario 19.13} pide compactos, y con cerrados no acotados la intersección puede ser vacía. No hay un conjunto límite $A$.
\end{pasos}
\par{\small \hyperlink{x-12}{volver al ejercicio} \quad·\quad \hyperlink{r-12}{respuesta de la pregunta 12}}
\needspace{6\baselineskip}
\subsection*{\hypertarget{s-13}{}Solución del ejercicio 13}
\emph{Calcule $d_H([0,n],[0,m])$ para $n\lt m$.}
\begin{pasos}
\item[(1)] Por la \href{https://decostruttivismo.github.io/MAT2514-repaso/pdf/MAT2514_apuntes_completos.pdf\#nameddest=teorema.36.5}{Proposición 36.5}, $d_H([0,n],[0,m])=\max\{\abs{0-0},\abs{n-m}\}=m-n$.
\item[] Por (1), $d_H([0,n],[0,n+1])=1$ para todo $n$: la sucesión ni siquiera es de Cauchy en $\mathcal K(\R)$.
\end{pasos}
\par{\small \hyperlink{x-13}{volver al ejercicio} \quad·\quad \hyperlink{r-13}{respuesta de la pregunta 13}}
\needspace{6\baselineskip}
\subsection*{\hypertarget{s-14}{}Solución del ejercicio 14}
\emph{Sean $f_1(x)=\frac x4$, $f_2(x)=\frac x4+\frac34$, $A_0=[0,1]$ y $A_n=W^n(A_0)$. Pruebe que $(A_n)$ es decreciente.}
\begin{pasos}
\item[(1)] $A_1=f_1([0,1])\cup f_2([0,1])=[0,\frac14]\cup[\frac34,1]\subseteq[0,1]=A_0$.
\item[(2)] $W$ es monótono: si $A\subseteq B$, entonces $f_i(A)\subseteq f_i(B)$ para cada $i$, y por lo tanto $W(A)\subseteq W(B)$.
\item[(3)] Supongamos $A_n\subseteq A_{n-1}$. Por (2), $A_{n+1}=W(A_n)\subseteq W(A_{n-1})=A_n$.
\item[] Por (1), (3) e inducción, $A_{n+1}\subseteq A_n$ para todo $n$.
\end{pasos}
\par{\small \hyperlink{x-14}{volver al ejercicio} \quad·\quad \hyperlink{r-14}{respuesta de la pregunta 14}}
\needspace{6\baselineskip}
\subsection*{\hypertarget{s-15}{}Solución del ejercicio 15}
\emph{En el Problema \hyperlink{prob-11}{11}, ¿qué $n$ garantiza $d_H(A_n,K)\lt 10^{-3}$?}
\begin{pasos}
\item[(1)] La cota es $d_H(A_n,K)\le\frac{4^{-n}}3$. Basta pedir $\frac{4^{-n}}3\lt 10^{-3}$.
\item[(2)] La condición de (1) equivale a $4^n\gt \frac{1000}3\approx333{,}3$.
\item[(3)] $4^4=256$ no cumple (2); $4^5=1024$ sí.
\item[] Por (3), basta $n=5$.
\end{pasos}
\par{\small \hyperlink{x-15}{volver al ejercicio} \quad·\quad \hyperlink{r-15}{respuesta de la pregunta 15}}
\needspace{6\baselineskip}
\subsection*{\hypertarget{s-16}{}Solución del ejercicio 16}
\emph{(\href{https://decostruttivismo.github.io/MAT2514-repaso/pdf/MAT2514_apuntes_completos.pdf\#nameddest=teorema.31.7}{Ejemplo 31.7}.) Sea $T(x)=\frac x2+\frac1x$ en $[1,2]$, con $k=\frac12$ y $x_0=1$. ¿Cuántas iteraciones garantizan error $\lt 10^{-6}$?}
\begin{pasos}
\item[(1)] $x_1=T(1)=\frac12+1=\frac32$. Luego $d(x_1,x_0)=\frac12$.
\item[(2)] Por (1), la cota a priori es $d(x_n,\sqrt2)\le\frac{(1/2)^n}{1/2}\cdot\frac12=2^{-n}$.
\item[(3)] Por (2), basta $2^{-n}\lt 10^{-6}$, es decir, $2^n\gt 10^6$.
\item[(4)] $2^{19}=524288$ no cumple (3); $2^{20}=1048576$ sí.
\item[] Por (4), basta $n=20$.
\end{pasos}
\par{\small \hyperlink{x-16}{volver al ejercicio} \quad·\quad \hyperlink{r-16}{respuesta de la pregunta 16}}
\needspace{6\baselineskip}
\subsection*{\hypertarget{s-17}{}Solución del ejercicio 17}
\emph{En el Problema \hyperlink{prob-12}{12(b)}, ¿por qué se toma un radio $R$ con un factor $\frac1{1-r}$?}
\begin{pasos}
\item[] Queremos $f_i\bigl(\cl B(x_0,R)\bigr)\subseteq\cl B(x_0,R)$ para cada $i$, donde cada $f_i$ es una contracción de razón $r\lt 1$.
\item[] Sea $x$ con $\norm{x-x_0}\le R$.
\item[(1)] Por la triangular, $\norm{f_i(x)-x_0}\le\norm{f_i(x)-f_i(x_0)}+\norm{f_i(x_0)-x_0}$.
\item[(2)] Como $f_i$ es contracción de razón $r$, $\norm{f_i(x)-f_i(x_0)}\le r\norm{x-x_0}\le rR$.
\item[(3)] Por (1) y (2), $\norm{f_i(x)-x_0}\le rR+\norm{f_i(x_0)-x_0}$.
\item[(4)] El lado derecho de (3) es $\le R$ exactamente cuando $\norm{f_i(x_0)-x_0}\le(1-r)R$.
\item[] Por (4), basta $R\ge\frac{\max_i\norm{f_i(x_0)-x_0}}{1-r}$.
\end{pasos}
\par{\small \hyperlink{x-17}{volver al ejercicio} \quad·\quad \hyperlink{r-17}{respuesta de la pregunta 17}}
\needspace{6\baselineskip}
\subsection*{\hypertarget{s-18}{}Solución del ejercicio 18}
\emph{Si además $X$ es compacto, ¿$T$ tiene un punto fijo?}
\begin{pasos}
\item[] Sí. Sea $g(x)=d(x,Tx)$.
\item[(1)] $g$ es continua: $T$ es continua (es $1$-Lipschitz) y la distancia es continua.
\item[(2)] Por (1) y el teorema de Weierstrass, $g$ alcanza su mínimo en algún $x_0\in X$.
\item[(3)] Supongamos $g(x_0)\gt 0$, es decir, $Tx_0\ne x_0$.
\item[(4)] Por la hipótesis aplicada a $x_0\ne Tx_0$: $g(Tx_0)=d(Tx_0,T(Tx_0))\lt d(x_0,Tx_0)=g(x_0)$.
\item[(5)] (4) contradice que $g(x_0)$ sea el mínimo de $g$. Luego (3) es falso y $g(x_0)=0$.
\item[] Por (5), $Tx_0=x_0$: $x_0$ es un punto fijo.
\end{pasos}
\par{\small \hyperlink{x-18}{volver al ejercicio} \quad·\quad \hyperlink{r-18}{respuesta de la pregunta 18}}
\needspace{6\baselineskip}
\subsection*{\hypertarget{s-19}{}Solución del ejercicio 19}
\emph{(\href{https://decostruttivismo.github.io/MAT2514-repaso/pdf/MAT2514_apuntes_completos.pdf\#nameddest=teorema.32.10}{Observación 32.10}.) Calcule la norma de $\Lambda f=f(0)-2f(\frac12)+f(1)$ en $C([0,1],\R)$.}
\begin{pasos}
\item[(1)] Por la cota de la justificación, $\norm\Lambda\le1+2+1=4$.
\item[(2)] Los signos que convienen son $+1$ en $0$, $-1$ en $\frac12$ y $+1$ en $1$.
\item[(3)] La poligonal con esos valores es $g(x)=\abs{4x-2}-1$: vale $1$ en $0$, $-1$ en $\frac12$ y $1$ en $1$.
\item[(4)] $g$ es continua y $\norm g_\infty=1$.
\item[(5)] $\Lambda g=1-2\cdot(-1)+1=4$.
\item[] Por (1), (4) y (5): $\norm\Lambda=4$.
\end{pasos}
\par{\small \hyperlink{x-19}{volver al ejercicio} \quad·\quad \hyperlink{r-19}{respuesta de la pregunta 19}}
\needspace{6\baselineskip}
\subsection*{\hypertarget{s-20}{}Solución del ejercicio 20}
\emph{Calcule $\dist(f,M)$ para $f(x)=x$ en $C([0,1],\R)$, y para $f(x)=x^2$ en $C([-1,1],\R)$.}
\begin{pasos}
\item[(1)] Para $f(x)=x$ en $[0,1]$: $\max f=1$ y $\min f=0$.
\item[] Por (1) y la fórmula de la justificación, $\dist(f,M)=\frac12$, con la constante óptima $\frac12$.
\item[(2)] Para $f(x)=x^2$ en $[-1,1]$: $\max f=1$ (en $x=\pm1$) y $\min f=0$ (en $x=0$).
\item[] Por (2), $\dist(f,M)=\frac12$, de nuevo con la constante $\frac12$.
\end{pasos}
\par{\small \hyperlink{x-20}{volver al ejercicio} \quad·\quad \hyperlink{r-20}{respuesta de la pregunta 20}}
\needspace{6\baselineskip}
\subsection*{\hypertarget{s-21}{}Solución del ejercicio 21}
\emph{Pruebe que si $M$ es un subespacio de un espacio normado, entonces $\cl M$ también lo es.}
\begin{pasos}
\item[] Sean $x,y\in\cl M$ y $\lambda\in\R$.
\item[(1)] Existen sucesiones $x_n\to x$ e $y_n\to y$ con $x_n,y_n\in M$.
\item[(2)] $x_n+\lambda y_n\in M$, porque $M$ es subespacio.
\item[(3)] $x_n+\lambda y_n\to x+\lambda y$, porque la suma y el producto por escalar son continuos.
\item[] Por (2) y (3), $x+\lambda y$ es límite de puntos de $M$, es decir, $x+\lambda y\in\cl M$.
\end{pasos}
\par{\small \hyperlink{x-21}{volver al ejercicio} \quad·\quad \hyperlink{r-21}{respuesta de la pregunta 21}}
\needspace{6\baselineskip}
\subsection*{\hypertarget{s-22}{}Solución del ejercicio 22}
\emph{¿Puede pasar lo mismo en $\R^n$ con $\norm\cdot_\infty$?}
\begin{pasos}
\item[] No.
\item[(1)] Todo funcional lineal en $\R^n$ es de la forma $f(x)=\sum_ic_ix_i$, con $c_i=f(e_i)$.
\item[(2)] Por la triangular, $\abs{f(x)}\le\sum_i\abs{c_i}\abs{x_i}\le\bigl(\sum_i\abs{c_i}\bigr)\norm x_\infty$.
\item[] Por (2), $f$ es acotado.
\end{pasos}
\par{\small \hyperlink{x-22}{volver al ejercicio} \quad·\quad \hyperlink{r-22}{respuesta de la pregunta 22}}
\needspace{6\baselineskip}
\subsection*{\hypertarget{s-23}{}Solución del ejercicio 23}
\emph{Pruebe que un funcional lineal acotado es Lipschitz.}
\begin{pasos}
\item[(1)] Por el \href{https://decostruttivismo.github.io/MAT2514-repaso/pdf/MAT2514_apuntes_completos.pdf\#nameddest=teorema.34.2}{Lema 34.2}, $\abs{f(z)}\le\norm f\,\norm z$ para todo $z$.
\item[(2)] Por linealidad, $f(x)-f(y)=f(x-y)$.
\item[] Por (1) con $z=x-y$ y (2): $\abs{f(x)-f(y)}\le\norm f\,\norm{x-y}$. Es $\norm f$-Lipschitz.
\end{pasos}
\par{\small \hyperlink{x-23}{volver al ejercicio} \quad·\quad \hyperlink{r-23}{respuesta de la pregunta 23}}
\needspace{6\baselineskip}
\subsection*{\hypertarget{s-24}{}Solución del ejercicio 24}
\emph{¿En qué paso exacto de la demostración se usa la completitud de $\R$?}
\begin{pasos}
\item[] En el paso (3) de la justificación: al definir $f(x)=\lim_nf_n(x)$.
\item[] Ese límite existe porque $(f_n(x))_n$ es de Cauchy en $\R$ y $\R$ es completo.
\end{pasos}
\par{\small \hyperlink{x-24}{volver al ejercicio} \quad·\quad \hyperlink{r-24}{respuesta de la pregunta 24}}
\needspace{6\baselineskip}
\subsection*{\hypertarget{s-25}{}Solución del ejercicio 25}
\emph{Con $X=\{0,1,2\}$ e $Y=\{0,2\}$ (distancia de $\R$), calcule $\dis R$ para $R=\{(0,0),(1,0),(2,2)\}$.}
\begin{pasos}
\item[(1)] $R$ es una correspondencia: las primeras coordenadas son $0,1,2$ y las segundas son $0,2$.
\item[(2)] El par $(0,0),(1,0)$ aporta $\abs{d(0,1)-d(0,0)}=\abs{1-0}=1$.
\item[(3)] El par $(0,0),(2,2)$ aporta $\abs{d(0,2)-d(0,2)}=\abs{2-2}=0$.
\item[(4)] El par $(1,0),(2,2)$ aporta $\abs{d(1,2)-d(0,2)}=\abs{1-2}=1$.
\item[] Por (2), (3) y (4), $\dis R=1$. Con el \href{https://decostruttivismo.github.io/MAT2514-repaso/pdf/MAT2514_apuntes_completos.pdf\#nameddest=teorema.40.4}{Teorema 40.4}, $d_{GH}(X,Y)\le\frac12$.
\end{pasos}
\par{\small \hyperlink{x-25}{volver al ejercicio} \quad·\quad \hyperlink{r-25}{respuesta de la pregunta 25}}
\needspace{6\baselineskip}
\subsection*{\hypertarget{s-26}{}Solución del ejercicio 26}
\emph{Sea $X=Y=\{0,1\}$ con la distancia de $\R$. Muestre que hay una correspondencia de distorsión $1$, y que aun así $d_{GH}(X,Y)=0$.}
\begin{pasos}
\item[(1)] $R=X\times Y$ es una correspondencia.
\item[(2)] Los elementos $(0,0)$ y $(0,1)$ de $R$ aportan $\abs{d(0,0)-d(0,1)}=\abs{0-1}=1$. Luego $\dis R\ge1$; y ningún par aporta más de $1$, así que $\dis R=1$.
\item[(3)] La correspondencia $R'=\{(0,0),(1,1)\}$ tiene distorsión $0$. Luego $d_{GH}(X,Y)=0$.
\item[] Por (2) y (3): una correspondencia mala no dice nada sobre la cota inferior.
\end{pasos}
\par{\small \hyperlink{x-26}{volver al ejercicio} \quad·\quad \hyperlink{r-26}{respuesta de la pregunta 26}}
\needspace{6\baselineskip}
\subsection*{\hypertarget{s-27}{}Solución del ejercicio 27}
\emph{(\href{https://decostruttivismo.github.io/MAT2514-repaso/pdf/MAT2514_apuntes_completos.pdf\#nameddest=teorema.39.7}{Ejemplo 39.7}.) Calcule $d_{GH}([0,1],[0,3])$.}
\begin{pasos}
\item[(1)] Cota inferior por diámetros: $d_{GH}\ge\frac12\abs{3-1}=1$.
\item[(2)] Cota superior: trasladamos los intervalos para centrarlos en $0$, obteniendo $[-\frac12,\frac12]$ y $[-\frac32,\frac32]$ dentro de $\R$. Las traslaciones son isometrías.
\item[(3)] Por la \href{https://decostruttivismo.github.io/MAT2514-repaso/pdf/MAT2514_apuntes_completos.pdf\#nameddest=teorema.36.5}{Proposición 36.5}, $d_H([-\frac12,\frac12],[-\frac32,\frac32])=\max\{1,1\}=1$. Por (2) y la \href{https://decostruttivismo.github.io/MAT2514-repaso/pdf/MAT2514_apuntes_completos.pdf\#nameddest=teorema.39.4}{Proposición 39.4}, $d_{GH}\le1$.
\item[] Por (1) y (3), $d_{GH}([0,1],[0,3])=1$.
\end{pasos}
\par{\small \hyperlink{x-27}{volver al ejercicio} \quad·\quad \hyperlink{r-27}{respuesta de la pregunta 27}}
\needspace{6\baselineskip}
\subsection*{\hypertarget{s-28}{}Solución del ejercicio 28}
\emph{Pruebe que si $S\subseteq X$ es una $\varepsilon$-red de $X$ (con $X$ compacto), entonces $d_{GH}(S,X)\le\varepsilon$.}
\begin{pasos}
\item[(1)] Ambiente común: $Z=X$, con $S\subseteq X$ con la métrica restringida.
\item[(2)] Por (1) y la \href{https://decostruttivismo.github.io/MAT2514-repaso/pdf/MAT2514_apuntes_completos.pdf\#nameddest=teorema.39.4}{Proposición 39.4}, $d_{GH}(S,X)\le d_H(S,X)$.
\item[(3)] Desviación de $S$: como $S\subseteq X$, vale $0$.
\item[(4)] Desviación de $X$: como $S$ es una $\varepsilon$-red, cada $x\in X$ cumple $d(x,S)\le\varepsilon$.
\item[] Por (2), (3) y (4): $d_{GH}(S,X)\le d_H(S,X)\le\varepsilon$.
\end{pasos}
\par{\small \hyperlink{x-28}{volver al ejercicio} \quad·\quad \hyperlink{r-28}{respuesta de la pregunta 28}}
\needspace{6\baselineskip}
\subsection*{\hypertarget{s-29}{}Solución del ejercicio 29}
\emph{Pruebe que $X_n=\{0,\frac1n,\dots,1\}$ es disconexo para $n\ge1$.}
\begin{pasos}
\item[] Sean $U=X_n\cap(-\infty,\frac1{2n})$ y $V=X_n\cap(\frac1{2n},\infty)$.
\item[(1)] $U=\{0\}$ y $V=\{\frac1n,\dots,1\}$. Ambos son no vacíos.
\item[(2)] $U$ y $V$ son abiertos en $X_n$, por ser intersecciones de $X_n$ con abiertos de $\R$.
\item[(3)] $U\cap V=\varnothing$, porque los intervalos $(-\infty,\frac1{2n})$ y $(\frac1{2n},\infty)$ son disjuntos.
\item[(4)] $U\cup V=X_n$, porque ningún punto de $X_n$ es igual a $\frac1{2n}$.
\item[] Por (1)–(4), $(U,V)$ es una disconexión de $X_n$ (\href{https://decostruttivismo.github.io/MAT2514-repaso/pdf/MAT2514_apuntes_completos.pdf\#nameddest=teorema.24.1}{Definición 24.1}).
\end{pasos}
\par{\small \hyperlink{x-29}{volver al ejercicio} \quad·\quad \hyperlink{r-29}{respuesta de la pregunta 29}}
\needspace{6\baselineskip}
\subsection*{\hypertarget{s-30}{}Solución del ejercicio 30}
\emph{(Problema \hyperlink{prob-17}{17(b)}.) Con $f_n(i,a)=\frac in$ en la escalera $G_n$, ¿es $f_n(G_n)$ una $\frac1n$-red de $[0,1]$?}
\begin{pasos}
\item[(1)] La imagen es $f_n(G_n)=\{0,\frac1n,\dots,1\}$.
\item[(2)] Dado $x\in[0,1]$, sea $k=\lfloor nx\rfloor$, de modo que $\frac kn\le x\le\frac{k+1}n$.
\item[(3)] Las distancias $x-\frac kn$ y $\frac{k+1}n-x$ son no negativas y suman $\frac1n$. Luego la menor es $\le\frac1{2n}$.
\item[] Por (1), (2) y (3), $f_n(G_n)$ es una $\frac1{2n}$-red de $[0,1]$, y en particular una $\frac1n$-red.
\end{pasos}
\par{\small \hyperlink{x-30}{volver al ejercicio} \quad·\quad \hyperlink{r-30}{respuesta de la pregunta 30}}
\needspace{6\baselineskip}
\subsection*{\hypertarget{s-31}{}Solución del ejercicio 31}
\emph{Para el redondeo $f(x)=\lfloor nx\rfloor/n$, pruebe que $\dis f\le\frac1n$.}
\begin{pasos}
\item[] Sean $x,x'\in[0,1]$. Llamamos $u=x-f(x)$ y $u'=x'-f(x')$.
\item[(1)] Por la definición de la parte entera, $0\le u\lt \frac1n$ y $0\le u'\lt \frac1n$.
\item[(2)] $(f(x)-f(x'))-(x-x')=u'-u$.
\item[(3)] Por (1), $u'-u$ es la diferencia de dos números de $[0,\frac1n)$. Luego $\abs{u'-u}\lt \frac1n$.
\item[(4)] Por la triangular inversa, $\bigl\lvert\abs{f(x)-f(x')}-\abs{x-x'}\bigr\rvert\le\abs{(f(x)-f(x'))-(x-x')}$.
\item[] Por (2), (3) y (4): $\bigl\lvert\abs{f(x)-f(x')}-\abs{x-x'}\bigr\rvert\lt \frac1n$ para todos $x,x'$. Luego $\dis f\le\frac1n$.
\end{pasos}
\par{\small \hyperlink{x-31}{volver al ejercicio} \quad·\quad \hyperlink{r-31}{respuesta de la pregunta 31}}
\needspace{6\baselineskip}
\subsection*{\hypertarget{s-32}{}Solución del ejercicio 32}
\emph{Pruebe que en un espacio métrico, si $x_n\to x$, entonces $C=\{x_n:n\in\N\}\cup\{x\}$ es compacto.}
\begin{pasos}
\item[] Sea $\{U_i\}_{i\in I}$ un cubrimiento abierto de $C$.
\item[(1)] Como $x\in C$, existe $i_0\in I$ con $x\in U_{i_0}$.
\item[(2)] $U_{i_0}$ es abierto y $x_n\to x$. Luego existe $N$ con $x_n\in U_{i_0}$ para todo $n\ge N$.
\item[(3)] Cada uno de los finitos puntos $x_1,\dots,x_{N-1}$ está en algún $U_{i_k}$, $k=1,\dots,N-1$.
\item[] Por (1), (2) y (3), $U_{i_0},U_{i_1},\dots,U_{i_{N-1}}$ es un subcubrimiento finito de $C$. Luego $C$ es compacto.
\end{pasos}
\par{\small \hyperlink{x-32}{volver al ejercicio} \quad·\quad \hyperlink{r-32}{respuesta de la pregunta 32}}
\needspace{6\baselineskip}
\subsection*{\hypertarget{s-33}{}Solución del ejercicio 33}
\emph{Para la escalera $G_n$ del Problema \hyperlink{prob-17}{17}, con $d_n((i,a),(j,b))=\frac{\abs{i-j}+\abs{a-b}}n$, calcule $\diam G_n$ y verifique que tiende a $1$.}
\begin{pasos}
\item[(1)] $\abs{i-j}\le n$, porque $0\le i,j\le n$; y $\abs{a-b}\le1$, porque $a,b\in\{0,1\}$.
\item[(2)] Por (1), $d_n((i,a),(j,b))\le\frac{n+1}n$.
\item[(3)] Los vértices $(0,0)$ y $(n,1)$ están a distancia exactamente $\frac{n+1}n$.
\item[(4)] Por (2) y (3), $\diam G_n=\frac{n+1}n=1+\frac1n$.
\item[] Por (4), $\diam G_n\to1=\diam[0,1]$, como predice la proposición.
\end{pasos}
\par{\small \hyperlink{x-33}{volver al ejercicio} \quad·\quad \hyperlink{r-33}{respuesta de la pregunta 33}}
\newpage\parte{Apéndice. Problemas de la guía del examen}
La guía tiene 17 problemas. Los Problemas~3, 7 y 8 se resolvieron en el taller de la Clase 43 y no estarán en el examen; los Problemas 15 y 16 son complementarios.
\needspace{5\baselineskip}
\subsection*{\hypertarget{prob-1}{}Problema 1. Redes finitas, distancia a un compacto y $d_H$}
Sea $(X,d)$ un espacio métrico y recordemos que para $\emptyset\ne A\subseteq X$ se escribe $d(x,A):=\inf_{a\in A}d(x,a)$.
\begin{enumerate}[label=(\alph*),leftmargin=2.2em,itemsep=2pt]
\item[(a)] Sea $A\subseteq X$ compacto y no vacío. Pruebe que para cada $x\in X$ el ínfimo $d(x,A)$ se alcanza: existe $a_0\in A$ con $d(x,a_0)=d(x,A)$.
\item[(b)] Sea $K\in\mathcal K(X)$ y sea $S\subseteq K$ una $\varepsilon$-red de $K$, es decir un conjunto con $d(x,S)\le\varepsilon$ para todo $x\in K$. Pruebe que $d_H(S,K)\le\varepsilon$.
\item[(c)] Deduzca que todo $K\in\mathcal K(X)$ es el límite, en $(\mathcal K(X),d_H)$, de una sucesión de subconjuntos finitos de $K$.
\end{enumerate}
{\small\itshape Sugerencia. En (a) use que $a\mapsto d(x,a)$ es continua y Weierstrass; en (c) use que un compacto es totalmente acotado.}
\par{\small\textbf{Preguntas del quiz relacionadas:} \hyperlink{r-1}{1}, \hyperlink{r-4}{4}, \hyperlink{r-5}{5}, \hyperlink{r-7}{7}}
\needspace{5\baselineskip}
\subsection*{\hypertarget{prob-2}{}Problema 2. Los continuos forman un cerrado de $\mathcal K(X)$}
Sea $(X,d)$ un espacio métrico. Para $A,B\subseteq X$ no vacíos escribimos $d(A,B):=\inf\{d(a,b):a\in A, b\in B\}$.
\begin{enumerate}[label=(\alph*),leftmargin=2.2em,itemsep=2pt]
\item[(a)] Sea $K\in\mathcal K(X)$ disconexo. Pruebe que existen $K_1,K_2\in\mathcal K(X)$ disjuntos con $K=K_1\cup K_2$, y que $\eta:=d(K_1,K_2)\gt 0$.
\item[(b)] Con la notación de (a), pruebe que todo $L\in\mathcal K(X)$ con $d_H(L,K)\lt \eta/3$ es disconexo.
\item[(c)] Deduzca que $\mathcal{CN}(X):=\{K\in\mathcal K(X):K\ \text{conexo}\}$ es cerrado en $(\mathcal K(X),d_H)$ y que, si $X$ es completo, $(\mathcal{CN}(X),d_H)$ es un espacio métrico completo.
\end{enumerate}
{\small\itshape Sugerencia. En (b) separe $L$ según cuál de los dos $K_i$ está cerca, y use que $d(\cdot,K_i)$ es continua.}
\par{\small\textbf{Preguntas del quiz relacionadas:} \hyperlink{r-2}{2}, \hyperlink{r-3}{3}, \hyperlink{r-6}{6}}
\needspace{5\baselineskip}
\subsection*{\hypertarget{prob-4}{}Problema 4. $d_H$ como norma del supremo, y Arzelà–Ascoli}
Sea $(Z,d)$ un espacio métrico compacto no vacío. Para $A\in\mathcal K(Z)$ definamos $\varphi_A\colon Z\to\mathbb{R}$ por $\varphi_A(z):=d(z,A)$.
\begin{enumerate}[label=(\alph*),leftmargin=2.2em,itemsep=2pt]
\item[(a)] Pruebe que $\mathcal F:=\{\varphi_A:A\in\mathcal K(Z)\}$ es un subconjunto equicontinuo y uniformemente acotado de $C(Z,\mathbb{R})$, y concluya por Arzelà–Ascoli que $\mathcal F$ es relativamente compacto en $(C(Z,\mathbb{R}),d_\infty)$.
\item[(b)] Pruebe que $d_\infty(\varphi_A,\varphi_B)=d_H(A,B)$ para todos $A,B\in\mathcal K(Z)$; es decir, $A\mapsto\varphi_A$ es un encaje isométrico de $(\mathcal K(Z),d_H)$ en $(C(Z,\mathbb{R}),d_\infty)$.
\item[(c)] Deduzca de (a) y (b) que $(\mathcal K(Z),d_H)$ es totalmente acotado.
\end{enumerate}
{\small\itshape Sugerencia. Para la desigualdad $\ge$ de (b) evalúe $\varphi_A-\varphi_B$ en los puntos de $A$ y en los de $B$.}
\par{\small\textbf{Preguntas del quiz relacionadas:} \hyperlink{r-4}{4}, \hyperlink{r-10}{10}}
\needspace{5\baselineskip}
\subsection*{\hypertarget{prob-5}{}Problema 5. El dual de $C(X,\mathbb{R})$ y la distancia a las constantes}
Sea $X$ un espacio métrico compacto con al menos dos puntos y $C(X,\mathbb{R})$ con la norma $\|\cdot\|_\infty$.
\begin{enumerate}[label=(\alph*),leftmargin=2.2em,itemsep=2pt]
\item[(a)] Sean $x_1,\dots,x_n\in X$ distintos y $c_1,\dots,c_n\in\mathbb{R}$. Pruebe que $\Lambda(f):=\sum_{i=1}^{n}c_i f(x_i)$ define un funcional lineal acotado sobre $C(X,\mathbb{R})$ con $\|\Lambda\|\le\sum_{i=1}^{n}|c_i|$.
\item[(b)] Pruebe que de hecho $\|\Lambda\|=\sum_{i=1}^{n}|c_i|$.
\item[(c)] Sea $M:=\{c\mathbf 1:c\in\mathbb{R}\}$ el subespacio de las funciones constantes. Pruebe que $M$ es cerrado, que $\operatorname{dist}(f,M)=\frac12\bigl(\max_X f-\min_X f\bigr)$ para toda $f\in C(X,\mathbb{R})$, y que para cada $f$ no constante existe $F\in C(X,\mathbb{R})^{*}$ explícito con $F|_M=0$, $\|F\|=1$ y $F(f)=\operatorname{dist}(f,M)$.
\end{enumerate}
{\small\itshape Sugerencia. En (b) construya $f$ con $\|f\|_\infty\le1$ y $f(x_i)=\operatorname{sign}(c_i)$ usando las funciones $\max\{0,1-d(\cdot,x_i)/\rho\}$ con $\rho$ pequeño.}
\par{\small\textbf{Preguntas del quiz relacionadas:} \hyperlink{r-19}{19}, \hyperlink{r-20}{20}, \hyperlink{r-21}{21}}
\needspace{5\baselineskip}
\subsection*{\hypertarget{prob-6}{}Problema 6. Un funcional no acotado y el dual de un espacio incompleto}
Sea $\ell^{\infty}$ el espacio de las sucesiones reales acotadas con $\|x\|_\infty=\sup_n|x_n|$, y sea $X_{\mathrm{fin}}\subseteq\ell^{\infty}$ el subespacio de las sucesiones de soporte finito (las que tienen $x_n=0$ salvo para finitos índices), con la norma heredada.
\begin{enumerate}[label=(\alph*),leftmargin=2.2em,itemsep=2pt]
\item[(a)] Pruebe que $X_{\mathrm{fin}}$ no es completo, exhibiendo una sucesión de Cauchy en $X_{\mathrm{fin}}$ que no converge en $X_{\mathrm{fin}}$.
\item[(b)] Pruebe que $f(x):=\sum_{n\ge1}n x_n$ está bien definido sobre $X_{\mathrm{fin}}$, es lineal y no es acotado; concluya que $f$ no es continuo en ningún punto de $X_{\mathrm{fin}}$.
\item[(c)] Pruebe, pese a (a), que $X_{\mathrm{fin}}^{*}$ es un espacio de Banach, indicando con precisión en qué paso de la demostración se usa una hipótesis de completitud y de qué espacio es esa completitud.
\end{enumerate}
{\small\itshape Sugerencia. Para (a) sirven las sumas parciales de la sucesión $(1/n)_n$; para (b), los vectores de la base canónica $e^{(m)}$.}
\par{\small\textbf{Preguntas del quiz relacionadas:} \hyperlink{r-22}{22}, \hyperlink{r-23}{23}, \hyperlink{r-24}{24}}
\needspace{5\baselineskip}
\subsection*{\hypertarget{prob-9}{}Problema 9. Cálculos con la distancia de Hausdorff}
\begin{enumerate}[label=(\alph*),leftmargin=2.2em,itemsep=2pt]
\item[(a)] En $\mathbb{R}$ con la métrica usual, calcule $d_H\bigl([0,1],[a,a+1]\bigr)$ para $a\ge0$ y $d_H\bigl([0,1],\{0,1\}\bigr)$, justificando cada supremo.
\item[(b)] En $\mathbb{R}^{2}$ con la norma euclidiana, sean $D:=\{x:\|x\|\le1\}$, $C:=\{x:\|x\|=1\}$ y $\{0\}$ el origen. Calcule $d_H(D,C)$ y $d_H(C,\{0\})$.
\item[(c)] Pruebe que para $A,B\subseteq X$ acotados y no vacíos se tiene $d_H(A,B)=0$ si y solo si $\overline{A}=\overline{B}$, y exhiba dos subconjuntos acotados distintos de $\mathbb{R}$ con $d_H=0$.
\end{enumerate}
{\small\itshape Sugerencia. En (b) recuerde que para $x\in D$ el punto de $C$ más cercano a $x\ne0$ es $x/\|x\|$.}
\par{\small\textbf{Preguntas del quiz relacionadas:} \hyperlink{r-7}{7}, \hyperlink{r-8}{8}, \hyperlink{r-9}{9}}
\needspace{5\baselineskip}
\subsection*{\hypertarget{prob-10}{}Problema 10. Uniones, sucesiones decrecientes y sucesiones crecientes en $\mathcal K(X)$}
Sea $(X,d)$ un espacio métrico.
\begin{enumerate}[label=(\alph*),leftmargin=2.2em,itemsep=2pt]
\item[(a)] Pruebe que para $A_1,A_2,B_1,B_2\in\mathcal K(X)$, $d_H\bigl(A_1\cup A_2, B_1\cup B_2\bigr) \le\max\bigl\{d_H(A_1,B_1), d_H(A_2,B_2)\bigr\}$, y deduzca que, para $B\in\mathcal K(X)$ fijo, la aplicación $A\mapsto A\cup B$ es $1$-Lipschitz de $(\mathcal K(X),d_H)$ en sí mismo.
\item[(b)] Pruebe que si $A_1\supseteq A_2\supseteq A_3\supseteq\cdots$ es decreciente en $\mathcal K(X)$, entonces $A:=\bigcap_{n}A_n\in\mathcal K(X)$ y $d_H(A_n,A)\to0$.
\item[(c)] Dé un ejemplo de sucesión creciente $B_1\subseteq B_2\subseteq\cdots$ en $\mathcal K(\mathbb{R})$ que no converja en $d_H$, y otro que sí converja, identificando el límite.
\end{enumerate}
{\small\itshape Sugerencia. En (b) aplique la propiedad de los compactos encajados a los restos $A_n\setminus A^{\varepsilon}$.}
\par{\small\textbf{Preguntas del quiz relacionadas:} \hyperlink{r-11}{11}, \hyperlink{r-12}{12}, \hyperlink{r-13}{13}}
\needspace{5\baselineskip}
\subsection*{\hypertarget{prob-11}{}Problema 11. Un IFS concreto, su atractor y cuántas iteraciones hacen falta}
Sea $X=[0,1]$ con la métrica usual y sean $f_1(x):=\frac x4,\qquad f_2(x):=\frac x4+\frac34$. Sea $W(A)=f_1(A)\cup f_2(A)$ el operador de Hutchinson asociado.
\begin{enumerate}[label=(\alph*),leftmargin=2.2em,itemsep=2pt]
\item[(a)] Pruebe que $f_1,f_2$ es un sistema de funciones iteradas sobre el espacio métrico completo $[0,1]$, con razón $r=\frac14$, que $W\bigl(\mathcal K([0,1])\bigr)\subseteq\mathcal K([0,1])$ y que existe un único $K\in\mathcal K([0,1])$ con $W(K)=K$.
\item[(b)] Sea $A_0:=[0,1]$ y $A_n:=W^{n}(A_0)$. Pruebe que $A_n$ es unión de $2^{n}$ intervalos cerrados de longitud $4^{-n}$, que la sucesión $(A_n)$ es decreciente, y que $K=\bigcap_{n}A_n$.
\item[(c)] Calcule $d_H\bigl(A_0,W(A_0)\bigr)$ y use la estimación de error a priori para dar un $n$ explícito que garantice $d_H(A_n,K)\lt 10^{-3}$.
\end{enumerate}
{\small\itshape Sugerencia. La estimación a priori de la Clase 31, leída en $(\mathcal K(X),d_H)$, es $d_H(A_n,K)\le\frac{r^{n}}{1-r} d_H(A_0,W(A_0))$.}
\par{\small\textbf{Preguntas del quiz relacionadas:} \hyperlink{r-6}{6}, \hyperlink{r-12}{12}, \hyperlink{r-14}{14}, \hyperlink{r-15}{15}, \hyperlink{r-16}{16}, \hyperlink{r-18}{18}}
\needspace{5\baselineskip}
\subsection*{\hypertarget{prob-12}{}Problema 12. Identificar un atractor y su dimensión de similaridad}
Sea $f_1,\dots,f_m$ un IFS sobre $\mathbb{R}^{d}$ con razones $r_i$ y $r:=\max_ir_i\lt 1$, sea $W$ su operador de Hutchinson y sea $K$ su atractor.
\begin{enumerate}[label=(\alph*),leftmargin=2.2em,itemsep=2pt]
\item[(a)] Pruebe que si $A_0\in\mathcal K(\mathbb{R}^{d})$ cumple $W(A_0)\subseteq A_0$, entonces la sucesión $A_n:=W^{n}(A_0)$ es decreciente y $K=\bigcap_{n}A_n$.
\item[(b)] Pruebe que un tal $A_0$ siempre existe: fijado $x_0\in\mathbb{R}^d$, si $R\ge\dfrac{\max_i\|f_i(x_0)-x_0\|}{1-r}$ entonces la bola cerrada $A_0:=\overline{B}(x_0,R)$ cumple $W(A_0)\subseteq A_0$ y $A_0\in\mathcal K(\mathbb{R}^{d})$.
\item[(c)] Aplíquelo al IFS de $\mathbb{R}^{2}$ dado por $f_i(x)=\frac{x+p_i}2$, $i=1,2,3$, con $p_1,p_2,p_3$ los vértices de un triángulo: identifique el atractor como intersección de una sucesión decreciente explícita y escriba su dimensión de similaridad, usando sin demostrarla la fórmula $\dim_{S}=\log m/\log(1/r)$ para un IFS de $m$ semejanzas de razón $r$.
\end{enumerate}
{\small\itshape Sugerencia. En (b) estime $\|f_i(x)-x_0\|$ separando $f_i(x)-f_i(x_0)$ de $f_i(x_0)-x_0$.}
\par{\small\textbf{Preguntas del quiz relacionadas:} \hyperlink{r-12}{12}, \hyperlink{r-14}{14}, \hyperlink{r-17}{17}}
\needspace{5\baselineskip}
\subsection*{\hypertarget{prob-13}{}Problema 13. $d_{GH}$ entre espacios finitos pequeños}
Usaremos libremente la fórmula $d_{GH}(X,Y)=\frac12\inf_R\operatorname{dis} R$, donde el ínfimo recorre las correspondencias $R\subseteq X\times Y$.
\begin{enumerate}[label=(\alph*),leftmargin=2.2em,itemsep=2pt]
\item[(a)] Sean $X=\{x_1,x_2\}$ con $d_X(x_1,x_2)=a\gt 0$ e $Y=\{y_1,y_2\}$ con $d_Y(y_1,y_2)=b\gt 0$. Pruebe que $d_{GH}(X,Y)=\frac12|a-b|$.
\item[(b)] Sea $T=\{A,B,C\}$ con las tres distancias mutuas iguales a $1$ y sea $P=\{p_1,p_2\}$ con $d(p_1,p_2)=1$. Pruebe que $d_{GH}(T,P)=\frac12$. Observe que la cota $\left|\operatorname{diam} T-\operatorname{diam} P\right|\le2 d_{GH}(T,P)$ no sirve aquí.
\item[(c)] Generalice: si $X$ consta de $n$ puntos mutuamente a distancia $1$ e $Y$ es un espacio métrico con menos de $n$ puntos, pruebe que $d_{GH}(X,Y)\ge\frac12$. Deduzca $d_{GH}(T,P)=d_{GH}(T,\{q\})$ para un espacio $\{q\}$ de un punto.
\end{enumerate}
{\small\itshape Sugerencia. En (b) y (c) mire, para cada $y\in Y$, el conjunto $S_y:=\{x\in X:(x,y)\in R\}$ y cuente.}
\par{\small\textbf{Preguntas del quiz relacionadas:} \hyperlink{r-25}{25}, \hyperlink{r-26}{26}, \hyperlink{r-27}{27}}
\needspace{5\baselineskip}
\subsection*{\hypertarget{prob-14}{}Problema 14. Redes finitas, la grilla y qué sobrevive al límite GH}
\begin{enumerate}[label=(\alph*),leftmargin=2.2em,itemsep=2pt]
\item[(a)] Sea $X$ un espacio métrico compacto no vacío y $S\subseteq X$ una $\varepsilon$-red de $X$. Pruebe que $d_{GH}(S,X)\le\varepsilon$.
\item[(b)] Deduzca que todo espacio métrico compacto no vacío es límite de Gromov–Hausdorff de espacios métricos finitos, y que para $X_n:=\frac1n\mathbb{Z}\cap[0,1]$ se tiene $d_{GH}\bigl(X_n,[0,1]\bigr)\le\frac1{2n}$, de modo que $X_n\to[0,1]$.
\item[(c)] Deduzca que la convergencia GH no preserva ni la conexidad ni la dimensión topológica; en cambio, pruebe que $\left|\operatorname{diam} X-\operatorname{diam} Y\right|\le2 d_{GH}(X,Y)$ y concluya que $\operatorname{diam} X_n\to\operatorname{diam}[0,1]$.
\end{enumerate}
{\small\itshape Sugerencia. En (a) use el ambiente común $Z=X$ con la inclusión de $S$ y la identidad de $X$.}
\par{\small\textbf{Preguntas del quiz relacionadas:} \hyperlink{r-28}{28}, \hyperlink{r-29}{29}, \hyperlink{r-33}{33}}
\needspace{5\baselineskip}
\subsection*{\hypertarget{prob-17}{}Problema 17. $\varepsilon$-isometrías: la escalera converge al segmento}
\begin{enumerate}[label=(\alph*),leftmargin=2.2em,itemsep=2pt]
\item[(a)] Pruebe que si existe una $\varepsilon$-isometría $f\colon X\to Y$ entre espacios métricos compactos no vacíos, entonces $d_{GH}(X,Y)\le\frac32\varepsilon$.
\item[(b)] Sea $G_n$ el grafo escalera con vértices $V_n=\{(i,a):0\le i\le n, a\in\{0,1\}\}$, aristas entre $(i,a)$ y $(i+1,a)$ y entre $(i,0)$ y $(i,1)$, todas de longitud $\frac1n$, dotado de la métrica del camino $d_n$. Pruebe que $d_n\bigl((i,a),(j,b)\bigr)=\bigl(|i-j|+|a-b|\bigr)/n$, que $f_n(i,a):=\frac in$ es una $\frac1n$-isometría de $G_n$ en $[0,1]$, y que $G_n\to[0,1]$ en GH.
\item[(c)] Calcule $\operatorname{diam} G_n$ y verifique la estimación $\left|\operatorname{diam} G_n-\operatorname{diam}[0,1]\right|\le2 d_{GH}(G_n,[0,1])$. Pruebe además que el conjunto $\mathcal C:=\{[G_n]:n\in\mathbb{N}\}\cup\{[[0,1]]\}$ es compacto en $(\mathbb M,d_{GH})$, pese a que todos los $G_n$ son finitos y $[0,1]$ no lo es.
\end{enumerate}
{\small\itshape Sugerencia. En (b), un camino entre $(i,a)$ y $(j,b)$ usa al menos $|i-j|$ largueros y al menos $|a-b|$ peldaños.}
\par{\small\textbf{Preguntas del quiz relacionadas:} \hyperlink{r-29}{29}, \hyperlink{r-30}{30}, \hyperlink{r-31}{31}, \hyperlink{r-32}{32}, \hyperlink{r-33}{33}}
\needspace{5\baselineskip}
\subsection*{\hypertarget{prob-15}{}Problema 15. Dos familias precompactas y dos que no lo son {\normalfont\small(complementario)}}
Recordemos que $\mathcal C\subseteq\mathbb M$ es uniformemente totalmente acotada si hay $D\ge0$ con $\operatorname{diam} X\le D$ para todo $X\in\mathcal C$ y, para cada $\varepsilon\gt 0$, un $N(\varepsilon)\in\mathbb{N}$ con $N(\varepsilon,X)\le N(\varepsilon)$ para todo $X\in\mathcal C$.
\begin{enumerate}[label=(\alph*),leftmargin=2.2em,itemsep=2pt]
\item[(a)] Sea $Z$ un espacio métrico compacto no vacío y $\mathcal C_Z:=\{[A]\in\mathbb M: A\subseteq Z\ \text{compacto no vacío}\}$. Pruebe que $\mathcal C_Z$ es uniformemente totalmente acotada y concluya que es precompacta en $(\mathbb M,d_{GH})$.
\item[(b)] Sea $X$ compacto no vacío y $\mathcal C:=\{[\lambda X]:\lambda\in(0,1]\}$, donde $\lambda X=(X,\lambda d)$. Pruebe que $\mathcal C$ es precompacta y que $\lambda_n\to0^{+}$ implica $\lambda_nX\to\{q\}$, el espacio de un punto.
\item[(c)] Pruebe que $\{[0,n]:n\in\mathbb{N}\}$ y $\{\Delta_n:n\in\mathbb{N}\}$, con $\Delta_n$ el espacio de $n$ puntos mutuamente a distancia $1$, no son precompactas, señalando en cada caso qué condición falla.
\end{enumerate}
{\small\itshape Sugerencia. En (a) parta de una $(\varepsilon/2)$-red de $Z$ y quédese con los centros cuyas bolas cortan a $A$, reemplazándolos por un punto de $A$.}
\needspace{5\baselineskip}
\subsection*{\hypertarget{prob-16}{}Problema 16. Arzelà–Ascoli y Gromov son el mismo teorema {\normalfont\small(complementario)}}
\begin{enumerate}[label=(\alph*),leftmargin=2.2em,itemsep=2pt]
\item[(a)] Sea $X$ un espacio métrico compacto no vacío, $x_0\in X$ y $\mathcal F:=\bigl\{f\in C(X,\mathbb{R}): f\ \text{es}\ 1\text{-Lipschitz}, f(x_0)=0\bigr\}$. Pruebe que $\mathcal F$ es equicontinua, uniformemente acotada por $\operatorname{diam} X$ y cerrada en $(C(X,\mathbb{R}),d_\infty)$, y concluya por Arzelà–Ascoli que $\mathcal F$ es compacta.
\item[(b)] Sean $D\gt 0$, $N\in\mathbb{N}$ y $\mathcal C_{N,D}:=\{[S]\in\mathbb M:\ \#S\le N,\ \operatorname{diam} S\le D\}$. Usando el conjunto de matrices $\mathcal A_N=\{A\in[0,D]^{N\times N}: a_{ii}=0, a_{ij}=a_{ji},\ a_{ij}\le a_{ik}+a_{kj}\}$ y la aplicación $\Phi$ que a cada $A\in\mathcal A_N$ le asocia la clase de isometría del espacio $\{1,\dots,N\}$ con la semimétrica $(a_{ij})$, identificando los puntos a distancia $0$, pruebe que $\mathcal C_{N,D}$ es compacta en $(\mathbb M,d_{GH})$.
\item[(c)] Explicite el paralelo probando las dos estimaciones \guillemotleft{}datos cercanos $\Rightarrow$ objetos cercanos\guillemotright{}: si $f,g\in\mathcal F$, si $\{y_1,\dots,y_m\}$ es una $\delta$-red de $X$ y si $|f(y_i)-g(y_i)|\le\varepsilon$ para todo $i$, entonces $d_\infty(f,g)\le\varepsilon+2\delta$; y si $A,A'\in\mathcal A_N$ cumplen $\rho(A,A'):=\max_{i,j}|a_{ij}-a'_{ij}|\le2\varepsilon$, entonces $d_{GH}\bigl(\Phi(A),\Phi(A')\bigr)\le\varepsilon$.
\end{enumerate}
{\small\itshape Sugerencia. En (b) recuerde que $\mathcal A_N$ es compacto por Heine–Borel y que $\Phi(\mathcal A_N)=\mathcal C_{N,D}$.}
\end{document}
